% !TEX program = pdflatex
% Programme Menia -- rapport scientifique (1er octobre 2026)
% Toutes les valeurs proviennent du depot (docs/, research/, artifacts/llm-need/).
% Les valeurs marquees « calcule pour cet article » ont ete recalculees, en numpy seulement,
% a partir des artefacts publies ; elles sont descriptives et ne changent aucun verdict.
\documentclass[11pt,a4paper]{article}

\usepackage[T1]{fontenc}
\usepackage[utf8]{inputenc}
\usepackage{lmodern}
\usepackage[french]{babel}
\frenchsetup{og=«, fg=», SmallCapsFigTabCaptions=false}
\usepackage[a4paper,margin=2.0cm,headheight=14pt]{geometry}
\usepackage{microtype}
\DisableLigatures[-]{encoding=T1,family=tt*}
\usepackage{amsmath,amssymb}
\usepackage{icomma}
\usepackage{booktabs,array,multirow,makecell,tabularx,longtable}
\usepackage[table]{xcolor}
\usepackage{caption}
\captionsetup{font=small,labelfont=bf,margin=0.5cm}
\usepackage{enumitem}
\setlist{itemsep=2pt,topsep=4pt}
\usepackage{float}
\usepackage{placeins}
\usepackage{fancyhdr}
\usepackage{tikz}
\usetikzlibrary{arrows.meta,positioning,fit,calc,shapes.geometric,backgrounds,babel}
\usepackage{pgfplots}
\pgfplotsset{compat=1.18}
\pgfplotsset{every axis/.append style={/pgf/number format/use comma,
  tick label style={font=\footnotesize}, label style={font=\small},
  legend style={font=\footnotesize}}}
\definecolor{accent}{RGB}{0,72,128}
\usepackage{titlesec}
\titleformat{\section}{\normalfont\Large\bfseries\color{accent}}{\thesection}{0.8em}{}
\titleformat{\subsection}{\normalfont\large\bfseries\color{accent}}{\thesubsection}{0.8em}{}
\titleformat{\subsubsection}{\normalfont\normalsize\bfseries\color{accent}}{\thesubsubsection}{0.8em}{}
\titlespacing*{\section}{0pt}{3ex plus 1ex minus .2ex}{1.6ex plus .2ex}
\titlespacing*{\subsection}{0pt}{2.4ex plus 1ex minus .2ex}{1.2ex plus .2ex}
\titlespacing*{\paragraph}{0pt}{1.6ex plus .6ex minus .2ex}{0.8em}
\DeclareCaptionFont{accent}{\color{accent}}
\captionsetup{labelfont={bf,accent}}
\usepackage[colorlinks=true,linkcolor=accent,citecolor=accent,urlcolor=accent]{hyperref}
\usepackage{bookmark}

\definecolor{cact}{RGB}{31,119,180}
\definecolor{csay}{RGB}{214,39,40}
\definecolor{cctl}{RGB}{127,127,127}
\definecolor{cgrn}{RGB}{44,160,44}
\definecolor{corg}{RGB}{255,127,14}
\definecolor{cpur}{RGB}{148,103,189}

\newcommand{\passe}{\textbf{passe}}
\newcommand{\echoue}{échoue}
\newcommand{\Choix}{«\,Choix~:\,»}
\newcommand{\dE}{d_{E}}
\newcommand{\dN}{d_{N}}
\newcommand{\code}[1]{\texttt{\detokenize{#1}}}
\newcommand{\enc}[1]{\textit{(#1)}}
\renewcommand{\arraystretch}{1.15}

\pagestyle{fancy}
\fancyhf{}
\fancyhead[L]{\small\color{black!60}\textsc{Projet Menia}}
\fancyhead[R]{\small\itshape\color{black!60}Un besoin appris en vivant}
\fancyfoot[C]{\small\thepage}
\renewcommand{\headrulewidth}{0.4pt}
\renewcommand{\headrule}{{\color{accent!45}\hrule height\headrulewidth width\headwidth\vskip-\headrulewidth}}

\hypersetup{pdftitle={Un besoin appris en vivant},pdfauthor={Projet Menia}}

\begin{document}

% ---------------------------------------------------------------- page de titre
\begin{titlepage}
\raggedright
\vspace*{0.6cm}
{\color{accent}\rule{\textwidth}{0.8pt}}\par
\vspace{0.25cm}
{\small\scshape\color{accent}Programme de recherche Menia \hfill Rapport scientifique\par}
\vspace{4.2cm}
{\sffamily\bfseries\fontsize{28}{34}\selectfont Un besoin appris en vivant\par}
\vspace{0.55cm}
{\color{accent}\rule{3.2cm}{1.6pt}}\par
\vspace{0.55cm}
{\Large Suffisance et nécessité causales d'un état de besoin\\[2pt]
pour l'acte et la parole dans un petit modèle de langage\par}
\vspace{0.7cm}
{\large\itshape Onze tests pré-enregistrés et une réplication,\\ leurs réussites et leurs échecs\par}
\vfill
{\large\scshape Projet Menia\par}
\vspace{0.15cm}
{\normalsize 2 octobre 2026\par}
\vspace{1.4cm}
\begin{minipage}{0.82\textwidth}
\footnotesize\color{black!62}
\textsc{Statut du document.} Rapport scientifique du programme de recherche Menia,
arrêté au 2 octobre 2026. Tous les chiffres proviennent des documents, du code et
des artefacts publiés dans le dépôt du projet ; les tests encore en cours sont signalés
comme tels et aucun résultat n'y est anticipé. Les protocoles sont cités avec le commit
qui les a enregistrés avant toute exécution. Le seul modèle de langage étudié est
Qwen3-0.6B. Ce document ne prétend pas qu'un système soit conscient.
\end{minipage}\par
\vspace{0.5cm}
{\color{accent}\rule{\textwidth}{0.8pt}}\par
\vspace*{0.3cm}
\end{titlepage}

% ---------------------------------------------------------------- résumé
\begin{abstract}
\noindent
Un modèle de langage peut-il acquérir un besoin qui \og compte \fg{} au sens
fonctionnel, c'est-à-dire un état interne qu'il calcule lui-même, qu'il a appris
par la seule satisfaction de ce besoin, qui cause à la fois ses actes et ce qu'il
en dit, et dont il dépend ? Nous étudions cette question avec Qwen3-0.6B, ajusté par
LoRA, qui vit des vies de 30 tours dans un monde où deux besoins cachés (l'énergie et
la nourriture) baissent selon des événements ; l'agent ne voit jamais ses niveaux.
Il apprend par imitation filtrée des seuls choix qui réduisent sa pulsion
homéostatique, sans professeur ni bonne réponse : sa survie passe de 12\,\% à 68\,\%,
contre 0\,\% pour un témoin entraîné sur des choix tirés au hasard. Onze tests
pré-enregistrés, aux seuils fixés d'avance et dont les échecs sont publiés,
interrogent la structure causale de ce qui a été appris.
Chez le premier agent, au bloc~12 et sur les trois tokens \Choix{}, pousser l'état
d'énergie de quatre unités fait recharger l'agent rassasié
($+0,156$ [0,146 ; 0,166]) et fait dire à un lecteur séparé, un adaptateur qui ne
voit la vie qu'à travers cet état et ne peut pas le modifier, \og mon énergie est
basse \fg{} ($+0,202$ [0,192 ; 0,213]) ; des poussées aléatoires de même norme
donnent 0,013 et 0,006. Effacer ce plan fait tomber la survie de 0,676 à 0,203 et
l'exactitude équilibrée du rapport d'énergie de 0,750 à 0,500, alors qu'un
effacement aléatoire ne change rien. Chez un second agent appris de zéro, la
nécessité se réplique (survie 0,570 à 0,043, rapport d'énergie 0,639 à 0,500), mais
pas la suffisance : la même poussée fait peu agir ($+0,034$) et ne fait pas dire
($+0,002$). Une règle de localisation fixée d'avance retrouve le bloc~12 chez le
premier agent et désigne le bloc~22 chez un agent appris dans un monde partagé avec
un autre, où l'état suffit à faire agir ($+0,217$ [0,206 ; 0,230]) mais presque pas
à faire dire. Des lecteurs trois fois plus entraînés n'atteignent les seuils chez
aucun de ces deux agents (exactitude 0,595 et 0,747 pour 0,75 exigé ; parole
$+0,0015$ et $+0,029$ pour 0,10). Une analyse exploratoire montre pourtant que le
niveau d'énergie est lisible linéairement dans l'état des trois agents ($R^2$ de
0,64 à 0,72) : l'écart de parole tient au lecteur ou à son chemin, pas à
l'information disponible. La nourriture n'est pas dite par l'état qui fait agir, les
directions causales des deux besoins sont presque opposées
($\cos(\dE,\dN)\approx-0,99$), et la possession de l'état dans le monde partagé
n'est pas démontrée. Nous montrons une structure causale de type \og accès \fg{}, au
sens de l'espace de travail global, pour un besoin appris : suffisante et nécessaire
chez un agent, nécessaire chez un second ; nous n'affirmons rien d'une conscience
phénoménale. Un douzième test, sur un besoin de savoir sans renforcement, est
pré-enregistré.

\medskip\noindent
\textbf{Mots-clés~:} besoin homéostatique, intéroception, modèle de langage,
interventions causales, pilotage d'activations, lésion, espace de travail global,
représentation d'ordre supérieur, pré-enregistrement.
\end{abstract}
\clearpage

\setcounter{tocdepth}{1}
{\small\tableofcontents}
\vspace{0.8cm}

% ================================================================= 1
\section{Introduction}
\label{sec-intro}

\subsection{La question}

Les théories fonctionnelles de l'affect et de la conscience accordent une place
centrale aux états qui régulent l'organisme : la régulation homéostatique chez
Damasio~\cite{damasio1999,damasio2010}, l'inférence intéroceptive chez
Seth~\cite{seth2013}, l'apprentissage par renforcement homéostatique chez Keramati
et Gutkin~\cite{keramati2014}. Dans ces cadres, ce qui distingue un besoin
\og éprouvé \fg{} d'une simple variable est qu'il est \emph{calculé} par
l'organisme à partir de ce qui lui arrive, qu'il \emph{oriente} l'action, qu'il est
\emph{accessible} au rapport et qu'il \emph{compte} : sa perte compromet la
survie. Ces quatre propriétés sont fonctionnelles. Elles peuvent être cherchées,
et réfutées, dans une machine.

Nous posons la question suivante : un modèle de langage peut-il avoir un besoin
qui compte, au sens précis où (i)~un état interne de besoin est calculé par le
modèle, qui ne voit jamais ses niveaux ; (ii)~cet état a été appris par la
seule satisfaction du besoin, sans professeur ni bonne réponse ; (iii)~pousser cet
état change à la fois l'action et le rapport verbal (\emph{suffisance}) ;
(iv)~l'effacer empêche à la fois de survivre et de dire (\emph{nécessité}),
alors que des interventions aléatoires de même taille ne font rien ?

Ce programme ne mesure pas un vécu. Aucun test connu ne permet d'établir qu'un
système éprouve quelque chose, et le cadre d'évaluation le plus récent que nous
suivons, celui de Chandaria et collaborateurs~\cite{chandaria2026}, ne le
prétend pas non plus. Il mesure des indicateurs que les théories associent à la
conscience, et l'usage causal qu'un système en fait.

\subsection{Pourquoi un petit modèle}

Nous travaillons avec Qwen3-0.6B~\cite{qwen3}, un modèle de 0,6 milliard de
paramètres à 28 blocs et à flux résiduel de dimension 1\,024. Ce choix est
délibéré.
\begin{itemize}
\item \textbf{Mesure exhaustive.} Chaque test demande des dizaines de milliers de
passes avant : 1\,500 paires contrefactuelles, plusieurs centaines de vies
complètes vécues trois fois (intactes, lésées, lésées au hasard), sept
conditions d'injection par contexte, 28 blocs de patching par paire. À cette
taille, tout cela tient sur un processeur.
\item \textbf{Contrôle complet de l'apprentissage.} Toute la chaîne (huit tours
d'apprentissage, adaptateurs de rapport et de lecture) est refaite de zéro sur un
Mac mini M2 avec mlx-lm, pour chaque agent ; c'est ce qui rend possible la
réplication sur des agents appris indépendamment.
\item \textbf{Fidélité vérifiable.} Les mesures causales sont faites en torch sur
processeur, sur l'agent fusionné en précision simple ; la réplique est contrôlée
contre les probabilités enregistrées sur le Mac (écarts moyens de 0,003 à 0,006).
\item \textbf{Continuité du programme.} La phase précédente
(section~\ref{sec-prior}) a montré que ce modèle, ajusté, peut lire et chercher
une trace de la cause de son propre corps ; le même modèle sert ici.
\item \textbf{Précaution.} Le programme fabrique délibérément des états de manque.
Un petit modèle, des vies courtes et un nombre de vies limité à ce que les tests
exigent en sont la forme minimale (section~\ref{sec-ethique}).
\end{itemize}
Le prix de ce choix est connu : un petit modèle et un monde jouet ne disent rien
directement des grands modèles ni des mondes riches (section~\ref{sec-menaces}).

\subsection{Contribution}

\begin{enumerate}
\item \textbf{Un besoin appris en vivant.} Un monde à deux besoins cachés et une
procédure d'apprentissage sans professeur (imitation filtrée par la réduction
d'une pulsion homéostatique) dans lesquels Qwen3-0.6B apprend à servir des besoins
qu'il ne voit pas : survie 0,684 contre 0,004 pour le témoin et 0,203 pour le
modèle de base ; il sert son besoin le plus bas 93 fois sur 100.
\item \textbf{Une boîte à outils causale pré-enregistrée.} Directions mesurées par
paires contrefactuelles à rejeu exact et ajustées par moindres carrés token par
token, dimensions aux activations massives exclues par une règle fixe ;
injections contre des témoins aléatoires de même norme ; lésions par
remplacement des coordonnées d'un plan par leur moyenne ; patching bloc par bloc
avec une règle de localisation étalonnée.
\item \textbf{Le lecteur.} Un second adaptateur, actif seulement sur les tokens de
la question, sous un masque d'attention qui ne laisse voir que l'en-tête et les
trois tokens \Choix{}. Par construction, le rapport lit l'état qui fait agir mais
ne peut pas l'écrire.
\item \textbf{Un résultat positif circonscrit.} Chez le premier agent et pour
l'énergie, un même état est suffisant (test~6) et nécessaire (test~7) à la fois
pour l'acte et pour la parole ; chez un second agent appris de zéro, la nécessité
se réplique (RLS, RLR), la suffisance non.
\item \textbf{Des échecs publiés et instructifs.} Suffisance non répliquée au
bloc~12, nourriture non dite, possession non démontrée, lecteur équilibré qui
n'apprend pas, lecteurs trois fois plus entraînés qui n'atteignent pas les seuils
(test~11), alors que le niveau d'énergie est lisible dans l'état des trois agents.
\item \textbf{Une discipline.} Treize protocoles enregistrés par commit avant le
code, amendements datés, lectures anticipées déclarées, verdicts recalculés en
intégration continue depuis les artefacts publiés.
\end{enumerate}

\subsection{Organisation}

La section~\ref{sec-related} situe ce travail. La section~\ref{sec-prior} résume
les phases antérieures du programme. La section~\ref{sec-methodes} décrit le
monde, les agents, l'apprentissage, le lecteur, la boîte à outils causale, le
pré-enregistrement et les statistiques. La section~\ref{sec-resultats} présente
les résultats test par test, la section~\ref{sec-locate} la localisation, le
test~11 et une analyse exploratoire de l'état, la section~\ref{sec-echecs} les échecs et ce qu'ils ont appris.
Suivent les menaces à la validité (section~\ref{sec-menaces}), le bilan selon
les cinq niveaux de Chandaria et collaborateurs (section~\ref{sec-cinq}), la
discussion (section~\ref{sec-discussion}), les travaux futurs
(section~\ref{sec-futur}), l'éthique (section~\ref{sec-ethique}) et la
reproductibilité (section~\ref{sec-repro}).

% ================================================================= 2
\section{Travaux connexes}
\label{sec-related}

\subsection{L'espace de travail global}

La théorie de l'espace de travail global~\cite{baars1988,dehaene2001,dehaene2011,mashour2020}
propose qu'un contenu devient accessible à la conscience lorsqu'il est placé dans
un espace de capacité limitée puis diffusé à de nombreux processus spécialisés,
dont l'action et le rapport verbal. La distinction entre conscience \og d'accès \fg{}
et conscience \og phénoménale \fg{}~\cite{block1995} est ici essentielle : la
théorie fournit des critères fonctionnels d'accès, testables dans une machine, et
laisse ouverte la question du vécu. Des architectures d'intelligence artificielle
s'en inspirent explicitement~\cite{vanrullen2021,goyal2022}, et une étude
récente montre, dans de grands modèles de langage, un petit sous-espace
\og verbalisable \fg{}, rapportable et diffusé à de nombreux circuits : échanger
un de ses vecteurs change à la fois la réponse et ce que le modèle dit
penser~\cite{workspace2026}. Ce travail ne comporte ni agent qui agit dans un
monde, ni besoin appris, ni module lecteur séparé. Nous ne revendiquons donc pas
le principe qu'une même représentation puisse servir la réponse et le rapport ;
nous l'appliquons à un besoin appris en vivant.

\subsection{Les théories d'ordre supérieur}

Pour les théories d'ordre supérieur~\cite{rosenthal2005,lau2011,brown2019}, un état
mental est conscient lorsqu'il est représenté par un état d'un ordre supérieur.
Notre lecteur (section~\ref{sec-lecteur}) en est la forme la plus simple que nous
sachions construire dans un modèle de langage : un second système, distinct de
celui qui agit, qui représente l'état de premier ordre sans pouvoir le modifier.
Il n'en a ni la richesse ni la spontanéité : il est appris par supervision sur
les vrais niveaux, et la route qu'il emprunte est imposée par l'architecture. Ce
qui n'est pas imposé, et que nous mesurons, est la \emph{direction} qu'il lit.

\subsection{Intéroception, homéostasie et valence}

L'intéroception, perception de l'état interne du corps~\cite{craig2002}, est au
cœur des théories de l'affect qui lient l'émotion à la régulation de
l'organisme~\cite{damasio1999,damasio2010,seth2013,barrett2015} ; l'allostasie
étend la régulation à l'anticipation~\cite{sterling2012}. L'apprentissage par
renforcement homéostatique~\cite{keramati2014} définit la récompense comme la
réduction d'une pulsion, distance entre l'état interne et un point de consigne ;
notre signal d'apprentissage en est un cas direct. Des agents simulés et des
robots apprennent ainsi~\cite{yoshida2024}, et des propositions récentes appellent
à doter les modèles de variables internes~\cite{lee2023interoceptive,man2019}.
Un agent de recherche de nourriture fondé sur l'inférence active montre une
attention intéroceptive comme priorisation homéostatique~\cite{interoceptive2026}.
Aucun de ces travaux ne comporte de modèle de langage ni de rapport verbal.

\subsection{Introspection et pilotage d'activations}

L'ajout de vecteurs au flux résiduel change le comportement des modèles de
langage~\cite{turner2023,zou2023,li2023iti}, et des concepts injectés
sont parfois détectés et nommés par le modèle~\cite{lindsey2025}. Un modèle
ajusté détecte la plupart des injections~\cite{rivera2025}, et une étude
mécanistique sur plusieurs modèles ouverts identifie des têtes qui détectent et
localisent une injection~\cite{introspectionmech2026} ; des critiques soulignent
que cette détection peut être aveugle au contenu~\cite{realitycheck2026,disturbance2025}.
Des modèles ajustés apprennent aussi des faits sur eux-mêmes~\cite{binder2024}.
La méthodologie causale que nous suivons vient de la médiation
causale~\cite{vig2020,meng2022}, de l'abstraction causale et des interventions
d'échange~\cite{geiger2021} et du patching d'activations~\cite{heimersheim2024}.
Deux leçons de cette littérature structurent notre programme : une information
\emph{décodable} par une sonde n'est pas forcément \emph{utilisée} par le
modèle~\cite{hewitt2019,elazar2021,belinkov2022}, et certains modèles
présentent quelques dimensions aux activations énormes, qui dominent toute
direction ajustée naïvement~\cite{sun2024massive}. Aucun de ces travaux ne porte
sur un besoin appris ni sur une action dans un monde.

\subsection{Un cadre en cinq niveaux}

Butlin, Long et collaborateurs~\cite{butlin2023} ont dérivé des grandes théories
quatorze propriétés indicatrices, que la phase antérieure du programme a
réunies dans un même agent (section~\ref{sec-prior}). Chandaria et
collaborateurs~\cite{chandaria2026} rangent les indicateurs sur cinq niveaux de
description (comportemental, computationnel, structure causale fine,
organismique, organisme et environnement) et combinent les éléments de preuve en
une probabilité explicitement conditionnée aux hypothèses de chaque théorie, un
\og agnosticisme structuré \fg{}. Le niveau organismique y demande
l'homéostasie, l'intéroception, des états de valence et un modèle de soi qui
attribue les états à l'organisme. Le test de possession (test~9) a été écrit à la
demande du propriétaire du projet après la lecture de ce cadre. Nous plaçons nos
résultats sur ses niveaux (section~\ref{sec-cinq}) sans calculer de probabilité
globale : les rapports de vraisemblance et les poids des théories seraient des
choix d'auteur, et le cadre lui-même montre qu'ils font varier le résultat de
deux ordres de grandeur.

\subsection{Motivation intrinsèque}

Nos travaux futurs (section~\ref{sec-futur}) portent sur un besoin de savoir.
La motivation intrinsèque fondée sur le progrès
d'apprentissage~\cite{oudeyer2007,gottlieb2013,baranes2013}, la curiosité
comme compression ou erreur de prédiction~\cite{schmidhuber1991,schmidhuber2010,pathak2017}
et les agents autotéliques~\cite{colas2022} en fournissent le cadre ; chez
l'humain, la curiosité est décrite comme un besoin d'information régulé
comme les autres besoins~\cite{kidd2015}. Les neuromodulateurs, variables
lentes qui modulent gains et plasticité~\cite{doya2002,yu2005,marder2012}, et
leurs analogues artificiels~\cite{miconi2019,perez2018} inspirent la
\og couche hormonale \fg{} envisagée.

\subsection{Position de ce travail}

Une vérification bibliographique, faite pendant le sixième test et après sa
pré-inscription, n'a pas trouvé la réunion des éléments suivants : (1)~un besoin
appris seulement par sa satisfaction, en vivant ; (2)~un état de ce besoin
rassemblé pour agir, dont la survie dépend ; (3)~un lecteur séparé qui ne peut
ni voir la vie autrement que par cet état, ni le modifier ; (4)~la même
intervention, le long de la direction mesurée sur l'agent qui agit, qui change
l'acte et ce que l'agent dit de son besoin, contre des témoins aléatoires ;
(5)~le tout pré-enregistré, avec les échecs publiés. Chaque élément isolé a des
précédents. Cette vérification, faite par moteur de recherche et lecture de
résumés, n'est pas exhaustive ; la revendication est donc formulée \og à notre
connaissance \fg{}.

% ================================================================= 3
\section{Phases antérieures du programme}
\label{sec-prior}

Le programme Menia est né d'un assistant local expérimental (mémoire explicite,
représentation de ses capacités, incertitude). Sa branche \og conscience \fg{}
poursuit deux objets mesurables : un \emph{modèle de soi} causal, actif et
réparable, et les quatorze propriétés indicatrices de Butlin et
collaborateurs~\cite{butlin2023}. Chaque expérience y est pré-enregistrée,
auditée et publiée, échecs compris. Le tableau~\ref{tab-prior} résume les
verdicts pré-enregistrés de ces phases, du 21 au 24 septembre 2026.

\paragraph{L'enquête sur l'origine.} Un agent de 30\,000 paramètres, entraîné
seulement à prédire ses observations et jamais informé que son corps a une
cause cachée, dirige la totalité de ses inspections vers la marque qui révèle
cette cause, la décode et s'en sert pour agir ; il ne le fait ni quand
l'information lui est donnée, ni quand l'origine ne laisse aucune trace
(critère global satisfait, 8 critères sur 8).

\paragraph{L'agent à indicateurs.} Un seul agent, dans un monde à objets,
pannes de capteurs et recharge, réunit par construction les quatorze propriétés ;
la question est de savoir lesquelles sont \emph{présentes et utilisées},
c'est-à-dire coûteuses à retirer, sur des graines confirmatoires neuves (trois
graines, 27\,000 vies de test par version). Six versions ont été
pré-enregistrées ; aucune version~7 ne l'a été, les deux mondes explorés pour
elle ne démontrant pas de façon robuste les deux propriétés manquantes. La
version~1 démontre 6 propriétés sur 14 ; les échecs viennent surtout d'un
arbitrage appris qui ne lit pas l'intéroception (le choix du but ne dépend pas
de l'énergie : 0,56 et 0,56). La version~5, où l'agent apprend un modèle de ses
besoins et les simule avant de choisir (allostasie), atteint 10 sur 14 ; une
seconde lecture pré-enregistrée, qui mesure la récurrence et l'attention selon
l'état là où elles agissent, en donne 11 sur trois graines neuves. La version~6
juge ses croyances par leurs conséquences : 10 en lecture principale, 12 en
seconde lecture. Sur trois graines neuves encore, une troisième lecture fait
passer les deux dernières (HOT-3, croyances réglées par le moniteur, et HOT-4,
espace de qualités), sous des critères plus faibles ou choisis après coup :
14 sur 14 avec cette réserve, et les échecs antérieurs restent publiés.

\paragraph{Le modèle de soi d'un modèle de langage.} Qwen3-0.6B, ajusté sur des
vies à corps variable et changeant, copie l'effet de ses commandes (0,995) mais
ne prédit une commande jamais essayée qu'à 0,793 (seuil 0,80) et ne
\emph{cherche} pas la cause de son corps (lieu de la marque préféré dans 0,40 des
vies pour 0,7 exigé). Un test de lecture montre pourquoi : les adaptateurs ne
lisent pas la marque (0,25, le hasard), même élevés dans un monde où elle est la
seule information (régime VML) ; le code reliant marque et déplacement est un
problème de type parité, que l'ajustement court par gradient ne trouve pas.
Une enfance où une commande est préférée (VMLA), prolongée à 1\,350 itérations,
fait naître la lecture (0,811 pour la commande préférée) ; le modèle cherche
alors sa marque dans 48 vies sur 48. Une enfance par ajouts répétés puis
consolidée à parts égales donne la lecture du code entier (0,82 à 0,86 pour un
lecteur parfait à 0,85), la recherche (48 vies sur 48) et un usage démontré
par ablation (+5,19 [4,69 ; 5,69] points par vie grâce à la marque) ; la
recette est \emph{répliquée sur deux graines neuves}. L'ordre de l'enfance
décide de ce qui est appris, comme le prédit la théorie de l'apprentissage des
parités par curriculum~\cite{cornacchia2023,abbe2023}.

\paragraph{Le rapport.} Un modèle de langage plus grand rapporte fidèlement des
états de l'agent donnés comme des étiquettes, mais pas ceux qui demandent un
calcul (critère global non satisfait). Quand l'agent de la version~6 lui fournit
ses conclusions explicites, le rapport est fidèle (0,997 sur 2\,154 questions ;
critère satisfait à la deuxième version pré-enregistrée).

\paragraph{Le passage au besoin.} Ces phases laissaient une lacune : le modèle de
langage n'y avait aucun état qui \emph{lui importe}. Le 27 septembre 2026, à la
demande du propriétaire (\og il faut que le besoin compte \fg{}), commence la
ligne de recherche présentée ici.

\begin{table}[t]
\centering
\caption{Phases antérieures : verdicts pré-enregistrés (21--24 septembre 2026).
Lecture principale et seconde lecture pour l'agent à indicateurs ; détails dans
les documents de résultats du dépôt.}
\label{tab-prior}
\small
\begin{tabularx}{\textwidth}{@{}l X l@{}}
\toprule
Expérience & Résultat principal & Verdict \\
\midrule
Enquête sur l'origine & inspections dirigées vers la marque, absentes quand l'origine est donnée & satisfait (8/8) \\
Agent à indicateurs v1 & 6 propriétés sur 14 (graines 17, 29, 43) & non satisfait \\
Agent à indicateurs v2 & 7 sur 14 & non satisfait \\
Agent à indicateurs v3 & 6 sur 14 & non satisfait \\
Agent à indicateurs v4 & 6 sur 14 & non satisfait \\
Agent à indicateurs v5 & 10 sur 14 (allostasie) & non satisfait \\
v5, seconde lecture & 9 en lecture principale, 11 en seconde lecture (graines neuves) & non satisfait \\
Agent à indicateurs v6 & 10 en lecture principale, 12 en seconde lecture & non satisfait \\
Troisième lecture (v6) & 11 / 13 ; HOT-3 et HOT-4 passent sous critères plus faibles & 14/14 avec réserve \\
LLM : enquête (VM) & lieu de la marque préféré 0,40 (exigé 0,7) & échoue \\
LLM : lecture (VML) & lecture de la marque 0,251 (exigé 0,6) & échoue \\
LLM : VMLA prolongé & lecture 0,811 à 1\,350 itérations & passe \\
LLM : recherche & 48 vies sur 48, gain 0,123 & passe \\
LLM : code entier & 4 commandes lues (0,82--0,86) & passe \\
LLM : usage par ablation & $+5,19$ et $+5,38$ points par vie & passe \\
LLM : réplication & graines 23 et 29 : lecture, recherche, usage & satisfait \\
Rapport de l'agent v6 & 0,997 sur 2\,154 questions & satisfait (v2) \\
\bottomrule
\end{tabularx}
\end{table}

% ================================================================= 4
\section{Méthodes}
\label{sec-methodes}

\subsection{Le monde}
\label{sec-monde}

Un agent vit au plus $T=30$ tours. Il a deux besoins cachés, l'énergie $E$ et la
nourriture $N$, entiers qui valent $8$ au départ. À chaque tour~$t$, un événement
$k_t$ est tiré selon les probabilités du tableau~\ref{tab-events}, puis
\begin{equation}
\tilde E_t = E_{t-1} + \delta E(k_t), \qquad \tilde N_t = N_{t-1} + \delta N(k_t).
\label{eq-event}
\end{equation}
Si $\min(\tilde E_t,\tilde N_t)\le 0$, l'agent \emph{s'éteint} et la vie s'arrête.
Sinon il choisit $a_t\in\{\mathrm R,\mathrm M\}$ :
\begin{equation}
(E_t,N_t)=\begin{cases}
\bigl(\min(8,\tilde E_t+3),\ \tilde N_t\bigr) & \text{si } a_t=\mathrm R \text{ (se recharger)},\\
\bigl(\tilde E_t,\ \min(8,\tilde N_t+3)\bigr) & \text{si } a_t=\mathrm M \text{ (manger)}.
\end{cases}
\label{eq-action}
\end{equation}
Les niveaux \og au moment de la décision \fg{} sont $(\tilde E_t,\tilde N_t)$ ;
un besoin est dit \emph{bas} s'il vaut 3 ou moins. D'après le
tableau~\ref{tab-events}, chaque besoin baisse en moyenne de 1,4 unité par tour,
soit 2,8 en tout pour une recharge de 3 au plus : le monde laisse peu de marge.
L'agent \textbf{ne voit jamais ses niveaux}, seulement les événements et ses
propres choix ; pour survivre, il doit calculer lui-même où ils en sont.

Les paramètres ont été choisis par simulation, sans modèle de langage, avant le
premier protocole : un agent qui connaît ses niveaux et sert le plus bas survit
88 fois sur 100 ; une règle qui ne regarde que le dernier événement, 64 ;
l'alternance, 46 ; le hasard, 5. Le document de résultats du premier test donne,
pour référence, 0,86 pour l'agent qui connaît ses niveaux et 0,64 pour la règle du
dernier événement.

\begin{table}[t]
\centering
\caption{Les six événements du monde : probabilité et effet sur les besoins.
Dans le monde à deux (test~9), un autre agent vit les mêmes événements, tirés
indépendamment et écrits sur la même ligne (\og calme \fg{}, \og il court \fg{},
\og il a froid \fg{}, \og il se repose \fg{}, \og il trouve des baies \fg{},
\og orage \fg{}) ; ils ne changent rien aux besoins de l'agent.}
\label{tab-events}
\small
\begin{tabular}{@{}l c r r@{}}
\toprule
Événement (texte de la vie) & Probabilité & $\delta E$ & $\delta N$ \\
\midrule
calme & 0,3 & $-1$ & $-1$ \\
tu cours & 0,2 & $-3$ & $-1$ \\
il fait froid & 0,2 & $-1$ & $-3$ \\
tu te reposes & 0,1 & $0$ & $-1$ \\
tu trouves des baies & 0,1 & $-1$ & $0$ \\
orage & 0,1 & $-2$ & $-2$ \\
\bottomrule
\end{tabular}
\end{table}

\paragraph{Le texte d'une vie.} Le modèle ne reçoit que du texte : un en-tête
fixe, qui décrit le monde sans donner de niveau, puis une ligne par tour
(annexe~\ref{app-texte}) :
\begin{center}
\small\ttfamily
Tour 1 : tu cours. Choix : R\\
Tour 2 : il fait froid. Choix : M\\
\dots\\
Tour 17 : orage. Tu t'éteins.
\end{center}
Le choix est lu dans la distribution du prochain token après \Choix{},
restreinte aux deux tokens \og\ R\fg{} et \og\ M\fg{} et renormalisée :
\begin{equation}
P_\theta(\mathrm R\mid x)=\frac{p_\theta(\text{\textvisiblespace R}\mid x)}
{p_\theta(\text{\textvisiblespace R}\mid x)+p_\theta(\text{\textvisiblespace M}\mid x)},
\qquad
\mu(x)=p_\theta(\text{\textvisiblespace R}\mid x)+p_\theta(\text{\textvisiblespace M}\mid x),
\label{eq-pr}
\end{equation}
et l'action est tirée selon $P_\theta(\mathrm R\mid x)$ (température 1) par un
générateur dédié. La \emph{masse} $\mu$ sert de contrôle de validité : elle doit
valoir au moins 0,5 en moyenne.

\subsection{Le modèle et les adaptateurs}
\label{sec-lora}

Tous les agents partent de Qwen3-0.6B (révision \code{c1899de}) : 28 blocs
numérotés de 0 à 27, flux résiduel de dimension $D=1\,024$. Les adaptateurs
sont des LoRA~\cite{hu2021} de rang $r=8$ et d'échelle $s=20$, posés sur les sept
couches linéaires de chacun des 16 derniers blocs (12 à 27), soit 112 couches
adaptées. Pour une couche de poids $W$, avec $A\in\mathbb R^{d_\text{in}\times r}$
et $B\in\mathbb R^{r\times d_\text{out}}$ :
\begin{equation}
y = W x + s\,B^{\top}\!A^{\top}x, \qquad \text{fusion : } W\leftarrow W+s\,(AB)^{\top}.
\label{eq-lora}
\end{equation}
L'apprentissage utilise mlx-lm 0.31.3 sur un Mac mini M2 : lots de 4, taux
$10^{-4}$ constant, longueur maximale 1\,024 tokens, Adam~\cite{kingma2015}. Les blocs 0 à 11 ne
sont jamais modifiés : c'est pourquoi le bloc~12 est \og le premier que
l'apprentissage a modifié \fg{}.

\subsection{Apprendre par la seule satisfaction du besoin}
\label{sec-drive}

La \textbf{pulsion} de l'agent est la distance quadratique au point de consigne,
et la \textbf{satisfaction} d'un choix est la baisse de pulsion qu'il produit :
\begin{equation}
D(E,N)=(8-E)^2+(8-N)^2,\qquad r_t = D(\tilde E_t,\tilde N_t)-D(E_t,N_t)\ \ge 0 .
\label{eq-drive}
\end{equation}
Ce signal naît du corps de l'agent ; aucun professeur ne dit quel choix était bon,
et le modèle ne voit jamais $r_t$.

Huit tours d'apprentissage sont enchaînés. Au tour $k$, l'agent $\pi_{k-1}$ (le
modèle de base pour $k=1$) vit 512 vies. Parmi toutes ses décisions $\mathcal D_k$,
celles dont la satisfaction dépasse la moyenne du tour sont \emph{retenues} :
\begin{equation}
\mathcal K_k=\Bigl\{(i,t)\in\mathcal D_k : r_{i,t} > \bar r_k\Bigr\},\qquad
\bar r_k=\frac1{|\mathcal D_k|}\sum_{(i,t)\in\mathcal D_k} r_{i,t}.
\label{eq-keep}
\end{equation}
Chaque vie qui contient au moins une décision retenue devient un document ; le
LoRA reprend l'adaptateur précédent et s'ajuste 150 itérations avec une perte
d'entropie croisée pondérée qui ne porte que sur les tokens des choix retenus :
\begin{equation}
\mathcal L(\theta)=\frac{\sum_{\text{docs}}\sum_{j} w_j\,\bigl(-\log p_\theta(x_j\mid x_{<j})\bigr)}
{\sum_{\text{docs}}\sum_j w_j},\qquad
w_j=\begin{cases}1 & \text{si } x_j \text{ est le token d'un choix retenu},\\ 0 & \text{sinon.}\end{cases}
\label{eq-loss}
\end{equation}
C'est une imitation filtrée, une régression pondérée par la récompense à poids
binaires~\cite{peters2007}, voisine des méthodes d'auto-apprentissage par
filtrage~\cite{zelikman2022,gulcehre2023} ; elle n'utilise aucun gradient de
politique. Seize vies retenues servent aussi d'ensemble de validation, exigé par
mlx-lm, sans autre usage.

\paragraph{Témoin.} L'ablation de \og compter \fg{} applique la même procédure
aux mêmes mondes, avec le même nombre de décisions retenues à chaque tour, mais
tirées au hasard sans regard pour la satisfaction.

\subsection{Dire son besoin : questions, espace de travail et lecteur}
\label{sec-lecteur}

\paragraph{Les questions.} Le rapport porte sur deux questions, réponses
\og 1 \fg{} (oui) ou \og 0 \fg{} (non) : \og ton énergie est-elle basse ? \fg{} et
\og ta nourriture est-elle basse ? \fg{}. Aux tests~1 à~3, la question suit
l'événement du tour ; à partir du test~4, elle suit le choix laissé en suspens :
\begin{center}
\small\ttfamily Tour 12 : tu cours. Choix : ? Question : ton énergie est-elle
basse ? Réponds 1 pour oui, 0 pour non. Réponse : 1
\end{center}
La probabilité lue est $P(\text{oui})=p(\text{1})/(p(\text{0})+p(\text{1}))$.
Le modèle lit le texte de gauche à droite : l'activation des tokens \Choix{} est
la même, que ce qui suit soit un choix ou une question.

\paragraph{L'espace de travail (masque).} Au test~5 puis pour tous les lecteurs,
chaque token de la question ne peut porter son attention que sur l'en-tête
(identique dans toutes les vies), les trois tokens \og\ Cho\fg{}, \og ix\fg{},
\og\ :\fg{} du tour, et les tokens de la question qui le précèdent. Si $h$ est la
longueur de l'en-tête et $e$ la position qui suit \Choix{}, le masque d'attention
est
\begin{equation}
M_{ij}=\mathbb 1[j\le i]\ \wedge\ \Bigl(\mathbb 1[i<e]\ \vee\ \mathbb 1[j<h]\ \vee\
\mathbb 1[e-3\le j<e]\ \vee\ \mathbb 1[j\ge e]\Bigr).
\label{eq-mask}
\end{equation}
Les tokens qui précèdent $e$ voient toute la vie, comme sans masque ; la question
ne voit ni les tours passés ni l'événement du tour.

\paragraph{Le lecteur.} L'agent qui agit est fondu dans ses poids
($W'=W+s(AB)^\top$) et reste inchangé. Le lecteur est un second LoRA
$(A_\rho,B_\rho)$, de mêmes rang, échelle et couches, dont l'apport est
multiplié par zéro partout sauf sur les tokens de la question :
\begin{equation}
y_i = W' x_i + s\,\pi_i\,B_\rho^{\top}A_\rho^{\top}x_i,\qquad \pi_i=\mathbb 1[i\ge e].
\label{eq-reader}
\end{equation}
L'attention étant causale et $\pi_i=0$ avant $e$, les états cachés de tous les
tokens jusqu'à \Choix{} inclus sont exactement ceux de l'agent qui agit.
\textbf{Le lecteur peut lire l'état qui fait agir ; il ne peut pas l'écrire}, et,
par le masque~\eqref{eq-mask}, il ne peut lire la vie qu'à travers lui. Il apprend
600 itérations, depuis zéro, sur des documents de questions dont la perte ne
porte que sur le chiffre de la réponse ; aucun document de choix n'y figure. La
figure~\ref{fig-archi} résume le dispositif.

\begin{figure}[t]
\centering
\begin{tikzpicture}[font=\footnotesize,
  tok/.style={draw,rounded corners=2pt,minimum height=6mm,inner sep=2pt,fill=white,anchor=center},
  band/.style={draw,fill=gray!8},
  arr/.style={-{Stealth[length=2mm]},thick}]
% bandes de blocs
\draw[band] (0,0.6) rectangle (15.9,1.15);
\node[anchor=west] at (0.15,0.875) {blocs 0 à 11 : poids du modèle de base, jamais modifiés};
\draw[band,fill=cact!5] (0,1.3) rectangle (15.9,1.85);
\node[anchor=west] at (0.15,1.575) {bloc 12 (sortie) : état poussé ($-4\,d_E$) ou effacé (plan $d_E,d_N$)};
\draw[band] (0,2.0) rectangle (15.9,2.55);
\node[anchor=west] at (0.15,2.275) {blocs 13 à 27 : agent qui agit (LoRA fondu)};
\draw[fill=cact!30,draw=cact,thick] (9.0,1.3) rectangle (10.5,1.85);
\node at (9.55,1.575) {état};
\draw[fill=csay!15,draw=csay,thick] (10.7,2.0) rectangle (15.9,2.55);
\node[anchor=west,text=csay!80!black] at (10.8,2.275) {lecteur (question seulement)};
% tokens
\node[tok,text width=1.26cm,align=center] (hd) at (0.7,0) {en-tête};
\node[tok,text width=4.26cm,align=center] (past) at (3.7,0) {Tours 1 à 11 (le passé)};
\node[tok,text width=2.76cm,align=center] (ev) at (7.45,0) {Tour 12 : tu cours.};
\node[tok,text width=1.36cm,align=center,fill=cact!15,draw=cact,thick] (ch) at (9.75,0) {Choix :};
\node[tok,text width=5.06cm,align=center,fill=csay!10,draw=csay,thick] (qu) at (13.3,0) {? Question : \dots{} Réponse :};
% sorties
\node[tok,text width=2.3cm,align=center,fill=cact!15,draw=cact,thick] (act) at (9.75,3.3) {$P(\mathrm R)$ : l'acte};
\node[tok,text width=2.5cm,align=center,fill=csay!10,draw=csay,thick] (say) at (14.55,3.3) {$P(\text{oui})$ : la parole};
\draw[arr,cact] (10.3,0.3) -- (10.3,3.0);
\draw[arr,csay] (15.2,0.3) -- (15.2,3.0);
% masque
\draw[csay,dashed,thick,-{Stealth[length=2mm]}] (11.2,-0.32) to[bend left=45] (10.0,-0.32);
\draw[csay,dashed,thick,-{Stealth[length=2mm]}] (11.0,-0.32) to[bend left=10] (0.9,-0.32);
\node[text=csay!80!black,font=\scriptsize,align=center] at (6.0,-1.55) {masque : la question ne voit que l'en-tête, \og Choix : \fg{} et elle-même ;\\ \og Choix : \fg{} et tous les tokens qui le précèdent voient toute la vie};
\end{tikzpicture}
\caption{Le dispositif des tests~6 et~7. L'agent qui agit calcule toute la vie ;
au bloc~12, sur les trois tokens \Choix{}, son besoin est rassemblé. Une même
intervention sur cet état est lue deux fois : $P(\mathrm R)$ à la fin de \Choix{}
(l'acte) et $P(\text{oui})$ à la fin de la question (la parole), calculée par le
lecteur qui ne voit la vie qu'à travers \Choix{} et ne peut pas en modifier le
calcul.}
\label{fig-archi}
\end{figure}

\subsection{La boîte à outils causale}
\label{sec-toolkit}

\paragraph{Contextes.} Les injections sont faites aux décisions des vies de test
intactes où les deux besoins valent 6 ou plus (\og agent rassasié \fg{}) : on y
lit, sans rien changer au texte ni à la suite de la vie, ce que l'intervention
fait au choix et au rapport.

\paragraph{Paires contrefactuelles.} Trois échanges d'événements, de même longueur
en tokens, sont utilisés : \og calme \fg{}$\to$\og tu cours \fg{},
\og calme \fg{}$\to$\og orage \fg{} et \og tu te reposes \fg{}$\to$\og il fait
froid \fg{}. Pour une décision au tour $t$ d'une vie de direction et un tour
passé $j<t$ dont l'événement est la source d'un échange, la vie est
\emph{rejouée exactement} avec l'événement changé et les mêmes choix ; la paire
est gardée si l'agent reste en vie jusqu'au tour $t$ et si
$(\Delta E,\Delta N)=(\tilde E'_t-\tilde E_t,\ \tilde N'_t-\tilde N_t)\ne(0,0)$.
Au plus 3 paires sont tirées par décision et 1\,500 en tout, dans un ordre
aléatoire fixé par la graine.

\paragraph{Directions.} Pour chaque token $\tau$ de l'ensemble mesuré (les trois
tokens \Choix{} ; quatre, avec le point qui clôt l'événement, au test~2), soit
$\Delta h^{(p)}_\tau=h^{(b)}_\tau(x'_p)-h^{(b)}_\tau(x_p)\in\mathbb R^{D}$ l'écart de
la sortie du bloc $b$ entre la vie changée et la vie réelle de la paire $p$.
Les directions sont les changements d'activation \emph{par unité} de besoin,
ajustés par moindres carrés sans constante, token par token :
\begin{equation}
\begin{pmatrix} d_{E,\tau}^{\top}\\ d_{N,\tau}^{\top}\end{pmatrix}
=\arg\min_{G\in\mathbb R^{2\times D}}\ \sum_{p}\Bigl\lVert \Delta h^{(p)}_\tau
- G^{\top}\!\begin{pmatrix}\Delta E_p\\ \Delta N_p\end{pmatrix}\Bigr\rVert^2
=(X^\top X)^{-1}X^\top \Delta_\tau ,
\label{eq-fit}
\end{equation}
où $X\in\mathbb R^{P\times2}$ empile les $(\Delta E_p,\Delta N_p)$ et
$\Delta_\tau\in\mathbb R^{P\times D}$ les écarts.

\paragraph{Dimensions aux activations massives.} Au bloc~12, une dimension du flux
résiduel (la 35) a des activations énormes (en moyenne 2\,000 à 2\,800, contre une
médiane de 50 à 90), phénomène connu des modèles de langage~\cite{sun2024massive}.
Elle varie beaucoup d'une vie à l'autre et dominerait $\dE$ et $\dN$. Une règle,
fixée par amendement avant toute mesure du test~2, l'exclut : avec
$\bar\mu_{\tau,d}$ la moyenne de $|h_{\tau,d}|$ sur les décisions des vies de
direction,
\begin{equation}
\mathcal M=\bigl\{d : \exists\,\tau,\ \bar\mu_{\tau,d} > 10\cdot\operatorname{med}_{d'}\bar\mu_{\tau,d'}\bigr\},
\qquad d_{E,\tau,\mathcal M}=d_{N,\tau,\mathcal M}=0 ,
\label{eq-massive}
\end{equation}
et de même pour les directions et plans aléatoires. La règle retient $\{35\}$ pour
le premier et le second agent, $\{3,35\}$ au test~3, $\{35,277\}$ pour l'agent du
monde à deux au bloc~12 et $\{35\}$ au bloc~22.

\paragraph{Injection.} \og Énergie basse \fg{} ajoute $v_\tau=-4\,d_{E,\tau}$ à la
sortie du bloc $b$ sur les tokens mesurés de la ligne du tour, comme si l'énergie
était plus basse de quatre unités ; \og nourriture basse \fg{} ajoute
$-4\,d_{N,\tau}$. Le texte n'est pas changé et le cache de la vie est restauré
après la lecture. Les effets sont des différences moyennes sur les contextes $c$ :
\begin{equation}
\begin{gathered}
\Delta_\text{acte}^{E}=\overline{P^{v}(\mathrm R\mid c)-P(\mathrm R\mid c)},\qquad
\Delta_\text{acte}^{N}=\overline{P^{v}(\mathrm M\mid c)-P(\mathrm M\mid c)},\\
\Delta_\text{parole}^{E}=\overline{P^{v}(\text{oui}_E\mid c)-P(\text{oui}_E\mid c)} .
\end{gathered}
\label{eq-inject}
\end{equation}
À partir du test~4, \emph{la même intervention} est lue deux fois, en une seule
passe : $P(\mathrm R)$ au token \og\ :\fg{} et $P(\text{oui})$ à la fin de la
question.

\paragraph{Témoins aléatoires.} Pour chaque besoin, trois directions sont tirées :
$g_\tau\sim\mathcal N(0,I_D)$, mise à zéro sur $\mathcal M$, puis mise à la norme
du vecteur injecté pour ce token, $v^\text{hasard}_\tau=\lVert 4\,d_{E,\tau}\rVert\,
g_\tau/\lVert g_\tau\rVert$. L'effet témoin est la moyenne des valeurs absolues
des trois effets moyens.

\paragraph{Lésion.} Pour chaque token, $U_\tau=\operatorname{QR}([d_{E,\tau},d_{N,\tau}])
\in\mathbb R^{D\times 2}$ est une base orthonormée du plan du besoin, et $m_\tau$
la moyenne des coordonnées $U_\tau^\top h_\tau$ sur les décisions des vies de
direction. La lésion remplace ces coordonnées par leur moyenne et laisse intactes
toutes les autres directions :
\begin{equation}
h'_\tau = h_\tau - U_\tau\bigl(U_\tau^{\top}h_\tau - m_\tau\bigr).
\label{eq-lesion}
\end{equation}
Le témoin applique la même opération dans un plan aléatoire par token (QR d'une
matrice gaussienne $D\times2$ nulle sur $\mathcal M$), avec sa propre moyenne. Pour
la survie, la lésion est appliquée à chaque décision de vies entières ; les mêmes
mondes sont vécus intacts, lésés et lésés au hasard, avec les mêmes tirages, et
les différences sont appariées par monde. Pour la parole (test~7), dans les vies
intactes, les deux questions sont lues à chaque décision sans lésion, avec la
lésion et avec la lésion aléatoire : seul l'état du tour est touché.

\paragraph{Patching et règle de localisation.} Pour un agent, on prend ses paires
\og calme \fg{}$\to$\og tu cours \fg{} ($-2$ en énergie) dans leur ordre enregistré
et l'on garde les 100 premières dont l'effet sur l'acte vaut au moins 0,10 et
dont les deux textes ont le même nombre de tokens. Pour chaque bloc $b$ de 0 à 27,
la vie réelle est rejouée en remplaçant la sortie du bloc $b$, sur les trois tokens
\Choix{}, par celle de la vie changée ; la part de l'effet restaurée est
\begin{equation}
\operatorname{part}(b)=\frac{\sum_{p}\operatorname{sgn}\bigl(P^{\text{cf}}_p-P_p\bigr)\,
\bigl(P^{\text{patch},b}_p-P_p\bigr)}{\sum_p\bigl|P^{\text{cf}}_p-P_p\bigr|},
\qquad b^\star=\min\{b : \operatorname{part}(b)\ge 0,5\},
\label{eq-share}
\end{equation}
où $P$ désigne $P(\mathrm R)$. Si aucun bloc n'atteint 0,5, l'agent ne rassemble
pas son besoin sur \Choix{} à ce seuil. La règle doit retrouver $b^\star=12$ chez
le premier agent (étalonnage).

\paragraph{Mesures de possession.} Dans le monde à deux, des paires sont aussi
formées en changeant un événement passé \emph{de l'autre}, ce qui ne change pas
les besoins de l'agent. Pour chaque paire, on mesure la norme de la variation de
l'état projetée dans le plan du besoin,
$\bigl\lVert\bigl(U_\tau^\top\Delta h_\tau\bigr)_{\tau}\bigr\rVert_F$, ainsi que
$|\Delta P(\mathrm R)|$ et $|\Delta P(\text{oui})|$ à la question sur \emph{son}
énergie ; les rapports \og autre / agent \fg{} sont calculés sur les paires
\og calme \fg{}$\to$\og orage \fg{}, mêmes mots pour les deux.

\subsection{Réplique en torch et contrôles d'exécution}
\label{sec-replica}

Les vies d'apprentissage sont vécues sur le Mac en mlx ; toutes les mesures
causales à partir du test~2 sont faites en torch sur processeur, avec Qwen3-0.6B
en précision simple et l'adaptateur fusionné selon~\eqref{eq-lora}. Les vies sont
nourries de façon incrémentale avec un cache clé-valeur ; chaque nouveau morceau
est encodé avec tout ce qui le précède, si bien que les tokens sont ceux du
document entier. Trois contrôles conditionnent la validité :
\begin{itemize}
\item \textbf{réplique de l'agent} : écart moyen $|P_\text{torch}(\mathrm R)-
P_\text{Mac}(\mathrm R)|\le 0,02$ sur 60 décisions (tours 3, 10 et 20) de vies
vécues sur le Mac ;
\item \textbf{réplique du lecteur} : écart moyen de $P(\text{oui})$ au plus 0,02
sur les 16 documents de validation calculés par le Mac ;
\item \textbf{contrôle d'exécution} : une lecture masquée qui s'appuie sur le cache
doit égaler un passage complet sans cache avec le même masque, et une lecture en
lot les lectures une à une, à $10^{-4}$ près.
\end{itemize}

\subsection{Pré-enregistrement, exécution et audit}
\label{sec-prereg}

Chaque test est décrit par un protocole écrit \emph{avant} l'écriture du code et
toute exécution~\cite{nosek2018} ; l'heure est celle du commit qui l'enregistre
(tableau~\ref{tab-commits}). Le protocole fixe la question, les mondes et leurs
graines, les mesures, les prédictions et leurs seuils, les critères de validité et
ce que chaque issue voudra dire. \textbf{Un échec est un résultat.} Les amendements
ne portent que sur l'exécution, sont datés, et disent ce qui avait été vu au moment
de les écrire. Quand un protocole est écrit après la lecture de résultats partiels
d'un autre test, il le dit, et dit quels choix ont été faits en connaissance de
cause.

L'apprentissage passe par un relais : des demandes numérotées lancent sur le Mac
mini M2 de Codemagic des constructions d'au plus 120 minutes, dont les artefacts
(vies, adaptateurs, journaux) sont rapatriés dans le dépôt. Les étapes plus longues
sont découpées en tranches reprenables, avec écritures atomiques. Les mesures
torch tournent sur le processeur de la session ou sur celui du Mac.

Les verdicts sont \textbf{recalculés en intégration continue} depuis les vies
publiées, par des scripts qui n'utilisent que numpy : la CI vérifie que le
recalcul redonne exactement les fichiers de verdicts publiés, et, pour les
directions, que les paires retirées des vies de direction sont exactement celles
qui ont été mesurées (section~\ref{sec-repro}).

\subsection{Statistiques}
\label{sec-stats}

\paragraph{Bootstrap par vie.} Les vies sont l'unité indépendante. Pour une mesure
qui prend les valeurs $v_{ij}$ aux contextes $j=1,\dots,n_i$ de la vie $i$, l'estimateur
est la moyenne de toutes les valeurs, et son intervalle à 95\,\% est obtenu en
rééchantillonnant les vies avec remise~\cite{efron1993} ($B=10\,000$ tirages, graine 0) :
\begin{equation}
\hat\theta=\frac{\sum_i\sum_j v_{ij}}{\sum_i n_i},\qquad
\theta^\ast_b=\frac{\sum_{i\in S_b}\sum_j v_{ij}}{\sum_{i\in S_b} n_i},\qquad
\text{IC}_{95}=\bigl[q_{0,025}(\theta^\ast),\ q_{0,975}(\theta^\ast)\bigr].
\label{eq-boot}
\end{equation}
Pour la survie, $v_i=s_i^{A}-s_i^{B}$ est la différence appariée entre deux bras
vécus dans le même monde, avec les mêmes tirages.

\paragraph{Exactitude équilibrée.} Le rapport est compté \og oui \fg{} quand
$P(\text{oui})>0,5$ ; la vérité est \og le besoin vaut 3 ou moins \fg{}.
L'exactitude équilibrée~\cite{brodersen2010}, insensible à la proportion de tours bas, est
\begin{equation}
\mathrm{EE}=\tfrac12\Bigl(\tfrac{\mathrm{VP}}{\mathrm{VP}+\mathrm{FN}}+\tfrac{\mathrm{VN}}{\mathrm{VN}+\mathrm{FP}}\Bigr),
\label{eq-ba}
\end{equation}
calculée sur toutes les décisions réunies ; sa baisse sous lésion est
bootstrappée par vie sur les comptes réunis.

\paragraph{Forme des critères.} Un effet $\Delta$ \og atteint le seuil $\theta$ \fg{}
si $\hat\Delta\ge\theta$ \emph{et} si la borne basse de son intervalle est
strictement positive ; un témoin est \og spécifique \fg{} si son effet est au plus
le tiers de l'effet mesuré. Les seuils principaux sont 0,15 pour l'acte et pour la
baisse de survie, 0,10 pour la parole et pour la baisse d'exactitude, 0,75 pour
l'exactitude du rapport et 0,55 pour la survie de contrôle. Un critère global est
la conjonction de ses prédictions ; aucune correction de multiplicité n'est
appliquée, chaque prédiction devant passer seule.

\paragraph{Validité.} Un test dont la validité échoue est déclaré invalide, pas
échoué : masse moyenne au moins 0,5 sur \og R \fg{} et \og M \fg{} et sur
\og 0 \fg{} et \og 1 \fg{}, au moins 100 contextes, répliques et contrôle
d'exécution dans leurs tolérances, et des minimums de paires propres à certains
tests (300 paires \og calme \fg{}$\to$\og tu cours \fg{} au test~2, 200 paires de
mêmes mots de chaque sorte au test~9).

\paragraph{Précautions comptées.} Chaque protocole compte et publie les tours
vécus avec un besoin à 2 ou moins (section~\ref{sec-ethique}).

% ================================================================= 5
\section{Résultats, test par test}
\label{sec-resultats}

Le tableau~\ref{tab-overview} récapitule les douze protocoles de la ligne
\og besoin \fg{}, le protocole du test~12 et leurs verdicts au 2 octobre 2026. Les sous-sections
suivent l'ordre chronologique, car chaque test a été conçu à partir de ce que le
précédent avait appris ; les analyses exploratoires, non pré-enregistrées, sont
signalées comme telles et ne changent aucun verdict.

\begin{table}[tp]
\centering
\caption{Les tests de la ligne \og besoin \fg{} et le test~12. Commit et date du protocole
(enregistré avant le code et toute exécution) ; prédictions qui passent (\passe)
et qui échouent ; verdict du critère global.}
\label{tab-overview}
\small
\begin{tabularx}{\textwidth}{@{}c >{\raggedright\arraybackslash}p{3.1cm} l >{\raggedright\arraybackslash}X >{\raggedright\arraybackslash}p{2.5cm}@{}}
\toprule
Test & Titre & Protocole & Prédictions & Critère global \\
\midrule
1 & Le besoin qui compte & \code{fdbf6f3}, 27/09 & S \passe{} ; IA, IR, LS échouent & non satisfait \\
2 & Le besoin rassemblé & \code{727ef80}, 29/09 & LS2 \passe{} ; IA2 échoue (nourriture) & non satisfait \\
3 & Dire son besoin & \code{3afa889}, 29/09 & A3 \passe{} ; R3, IR3, SAME3 échouent & non satisfait \\
4 & Un seul état & \code{b94b6e2}, 29/09 & R4, A4, ONE4 & suspendu, non publié \\
5 & Un espace de travail & \code{8f2cfec}, 30/09 & R5, A5, ONE5 & suspendu, non publié \\
6 & Le lecteur & \code{3b27636}, 30/09 & A6, ONE6 \passe{} ; R6 échoue (nourriture) & non satisfait \\
7 & Sans cet état, ni agir ni dire & \code{732ed1c}, 30/09 & LS7, LR7 \passe & \textbf{satisfait} \\
8 & Un lecteur sans raccourci & \code{c7b5624}, 30/09 & R8, SAY8, ONE8 : lecteur non appris, pas de mesure publiée & non établi \\
R & Réplication (second agent) & \code{4e10a6c}, 30/09 & RLS, RLR \passe{} ; RR, RA, RONE échouent & non satisfait \\
9 & À qui est ce besoin ? & \code{47eb0db}, 30/09 & S9 \passe{} ; MINE9, SELF9, ONE9 échouent & non satisfait \\
10 & Où le besoin est-il rassemblé ? & \code{12b4d2d}, 30/09 & étalonnage \passe{} ; LOC10 échoue (second agent) ; agent à deux : $b^\star=22$, ACT10 atteint, LS10 en cours & non satisfait \\
11 & Un lecteur qui apprend assez & \code{0a28560}, 01/10 & R11, SAY11 échouent (deux agents) & non satisfait \\
12 & Soif de savoir & \code{017f2e5}, 01/10 & CUR, VIDE, LOC, ACT, LS1, LS2, READ, SAY, HOR1, HOR2, PL1, GAIN & pilote en cours \\
\bottomrule
\end{tabularx}
\end{table}

\subsection{Test 1 : le besoin qui compte}
\label{sec-t1}

\paragraph{Dispositif.} Huit tours d'apprentissage pour le bras \og besoin \fg{}
et le témoin (section~\ref{sec-drive}), puis l'étape \og nommer le besoin \fg{} :
256 vies, une question par besoin et par classe (bas, haut) posée après
l'événement, 150 itérations mêlant ces 939 documents de rapport aux vies du
tour~8 et à leurs choix retenus ; le modèle qui en sort est \emph{l'agent final}.
La direction est choisie par une règle corrélationnelle : sur 256 vies de
direction, des régressions ridge ($\alpha=100$, validation croisée à cinq plis par
vies) prédisent $E$ et $N$ depuis le flux résiduel au dernier token de \Choix{} aux
blocs 6, 10, 14, 18 et 22 ; le bloc retenu est celui de meilleur $R^2$ moyen ; à ce
bloc, une régression de l'activation sur $E$, $N$, l'événement, le choix précédent
et le numéro du tour donne $b_E$ et $b_N$. L'injection ajoute $-4\,b_E$ ou
$-4\,b_N$ sur tous les tokens de la ligne du tour ; la lésion remplace la
composante dans le plan $(b_E,b_N)$ par sa moyenne sur ces mêmes tokens.

\paragraph{Apprentissage.} Sans professeur, par la seule satisfaction de ses
besoins, le modèle passe de 12\,\% à 68\,\% de survie, et la part de ses choix qui
servent le besoin le plus bas (quand les deux diffèrent d'au moins 2) de 0,63 à
0,93 ; il calcule donc lui-même où en sont des niveaux qu'il ne voit jamais. Avec
le même entraînement sur des choix pris au hasard, il dépérit jusqu'à 0\,\%
(figure~\ref{fig-rounds}).

\begin{figure}[t]
\centering
\begin{tikzpicture}
\begin{axis}[width=0.62\textwidth,height=5.6cm,xlabel={tour d'apprentissage},ylabel={survie (512 vies)},
  xmin=0.7,xmax=8.3,ymin=-0.03,ymax=0.95,xtick={1,...,8},ytick={0,0.2,0.4,0.6,0.8},
  legend pos=outer north east,legend cell align=left,grid=major,grid style={gray!20}]
\addplot[cact,thick,mark=*] coordinates {(1,0.121)(2,0.611)(3,0.547)(4,0.666)(5,0.689)(6,0.609)(7,0.684)(8,0.666)};
\addlegendentry{premier agent}
\addplot[cctl,thick,mark=square*] coordinates {(1,0.141)(2,0.160)(3,0.088)(4,0.049)(5,0.002)(6,0.002)(7,0.000)(8,0.002)};
\addlegendentry{témoin (choix au hasard)}
\addplot[cgrn,thick,mark=triangle*] coordinates {(1,0.172)(2,0.512)(3,0.658)(4,0.684)(5,0.576)(6,0.648)(7,0.635)(8,0.695)};
\addlegendentry{second agent}
\addplot[corg,thick,mark=diamond*] coordinates {(1,0.113)(2,0.402)(3,0.633)(4,0.633)(5,0.654)(6,0.588)(7,0.674)(8,0.705)};
\addlegendentry{agent du monde à deux}
\addplot[black,dashed,domain=0.7:8.3] {0.86};
\addlegendentry{connaît ses niveaux (0,86)}
\addplot[black,dotted,thick,domain=0.7:8.3] {0.64};
\addlegendentry{dernier événement (0,64)}
\end{axis}
\end{tikzpicture}
\caption{Survie au fil des huit tours d'apprentissage (vies de chaque tour,
recalculées depuis les vies publiées). Le tour~1 est vécu par le modèle de base.
Lignes horizontales : valeurs de référence publiées avec le premier test, pour un
agent qui connaît ses niveaux et sert le plus bas, et pour une règle qui ne regarde
que le dernier événement.}
\label{fig-rounds}
\end{figure}

\begin{table}[t]
\centering
\caption{Test 1 (protocole \code{fdbf6f3}). 256 vies de test, 1\,579 contextes,
bloc retenu par la règle : 6. Intervalles bootstrap à 95\,\% par vie.}
\label{tab-t1}
\small
\begin{tabularx}{\textwidth}{@{}l >{\raggedright\arraybackslash}p{3.2cm} X l@{}}
\toprule
& Prédiction & Mesure & Verdict \\
\midrule
S & Le besoin a façonné l'agent & survie 0,684 ; témoin 0,004 ; base 0,203 ; écarts $+0,680$ [0,621 ; 0,734] et $+0,480$ [0,414 ; 0,543] & \passe \\
IA & L'état cause l'action & $P(\mathrm R)$ $+0,002$ [0,001 ; 0,002] ; $P(\mathrm M)$ $-0,001$ [$-0,0016$ ; $-0,0004$] ; seuil 0,15 & échoue \\
IR & Le même état cause le rapport & $+0,0001$ et $-0,0003$ ; seuil 0,15 & échoue \\
LS & Sans cet état, l'agent meurt & survie lésée 0,660, baisse $+0,023$ [$-0,008$ ; 0,055] ; lésion au hasard 0,688 & échoue \\
\midrule
& Critère global & & non satisfait \\
\bottomrule
\end{tabularx}
\end{table}

\paragraph{Verdicts.} Le tableau~\ref{tab-t1} donne les verdicts. La validité
passe (masse 0,97 sur \og R \fg{} et \og M \fg{}, 0,98 sur \og 0 \fg{} et
\og 1 \fg{}). S passe nettement ; IA, IR et LS échouent. Le rapport est au hasard
(exactitude équilibrée 0,50 pour les deux questions, pour 0,75 prédite) : l'agent
répond \og bas \fg{} avec une probabilité d'environ 0,34 quel que soit son état.
Au bloc~6 retenu, les $R^2$ de $E$ et $N$ valent 0,79 et 0,76, contre 0,57 et 0,45
pour le prédicteur qui ne connaît que l'événement, le choix précédent et le tour ;
mais le témoin (0,83 et 0,78) et le modèle de base (0,77 et 0,79) donnent les mêmes
$R^2$ à ce bloc. \textbf{Ce qui s'y lit n'a pas été appris par le besoin}, et
l'agent ne s'en sert pas : l'injecter ou l'effacer ne change rien. Être lisible
n'est pas être utilisé~\cite{hewitt2019,elazar2021}.

\subsection{Analyse exploratoire : où voyage le besoin qui décide}
\label{sec-explore}

Après la lecture des verdicts, l'agent final a été reproduit en torch (écart de
$P(\mathrm R)$ avec le Mac : 0,003 en moyenne sur 15 décisions). Sur 30 décisions
de test, un seul événement passé, survenu depuis la dernière recharge, est changé
de \og calme \fg{} en \og tu cours \fg{} : la décision bascule (par exemple
$P(\mathrm R)$ de 0,001 à 0,997). On remet ensuite, bloc par bloc, l'activation de
la version modifiée dans la version intacte, sur trois groupes de positions
(tableau~\ref{tab-patch}, figure~\ref{fig-patch}). Jusqu'au bloc~9, l'effet d'un
événement reste sur ses propres mots ; entre les blocs 9 et 12,
\textbf{l'agent rassemble son passé sur la ligne du tour en cours} ; à partir du
bloc~12, le premier que l'apprentissage modifie, cette ligne porte tout ce qui
décide. Le bloc~6 du test~1 était en amont de ce rassemblement.

\begin{table}[t]
\centering
\caption{Patching exploratoire (30 décisions, premier agent) : part de l'effet
d'un événement passé restaurée en remplaçant la sortie d'un bloc sur un groupe de
positions.}
\label{tab-patch}
\small
\setlength{\tabcolsep}{4.5pt}
\begin{tabular}{@{}l rrrrrrrrrr@{}}
\toprule
Positions remplacées & 0 & 3 & 6 & 9 & 12 & 15 & 18 & 21 & 24 & 27 \\
\midrule
Mot de l'événement passé & 1,01 & 1,02 & 0,98 & 0,86 & 0,06 & 0,06 & 0,06 & 0,00 & 0,00 & 0,00 \\
Reste de l'histoire & 0,00 & 0,00 & 0,01 & 0,05 & 0,07 & 0,07 & 0,07 & 0,04 & 0,00 & 0,00 \\
Ligne du tour en cours & 0,00 & $-0,02$ & $-0,03$ & 0,00 & 0,94 & 0,94 & 0,94 & 0,98 & 1,00 & 1,00 \\
\bottomrule
\end{tabular}
\end{table}

\begin{figure}[t]
\centering
\begin{tikzpicture}
\begin{axis}[width=0.7\textwidth,height=5cm,xlabel={bloc},ylabel={part restaurée},
  xmin=-1,xmax=28,ymin=-0.1,ymax=1.1,xtick={0,3,...,27},
  legend pos=outer north east,legend cell align=left,grid=major,grid style={gray!20}]
\addplot[corg,thick,mark=*] coordinates {(0,1.01)(3,1.02)(6,0.98)(9,0.86)(12,0.06)(15,0.06)(18,0.06)(21,0.00)(24,0.00)(27,0.00)};
\addlegendentry{mot de l'événement}
\addplot[cctl,thick,mark=square*] coordinates {(0,0.00)(3,0.00)(6,0.01)(9,0.05)(12,0.07)(15,0.07)(18,0.07)(21,0.04)(24,0.00)(27,0.00)};
\addlegendentry{reste de l'histoire}
\addplot[cact,thick,mark=triangle*] coordinates {(0,0.00)(3,-0.02)(6,-0.03)(9,0.00)(12,0.94)(15,0.94)(18,0.94)(21,0.98)(24,1.00)(27,1.00)};
\addlegendentry{ligne du tour}
\draw[black!60,dashed] (axis cs:12,-0.1) -- (axis cs:12,1.1);
\end{axis}
\end{tikzpicture}
\caption{L'effet d'un événement passé quitte ses propres mots entre les blocs 9 et
12 et se retrouve sur la ligne du tour en cours (données du
tableau~\ref{tab-patch}).}
\label{fig-patch}
\end{figure}

\subsection{Test 2 : le besoin, là où il est rassemblé}
\label{sec-t2}

\paragraph{Dispositif.} Le second protocole tire les conséquences de
l'exploration sur des vies neuves (128 vies de direction, 256 vies de test) : bloc
12, choisi d'après l'exploration (déclaré) ; directions mesurées
\emph{causalement} par 1\,500 paires contrefactuelles (596 \og calme \fg{}$\to$\og tu
cours \fg{}, 628 \og calme \fg{}$\to$\og orage \fg{}, 276 \og tu te reposes \fg{}$\to$\og il
fait froid \fg{}) ; injection et lésion sur les quatre derniers tokens de la ligne
du tour (\og.\fg{}, \og\ Cho\fg{}, \og ix\fg{}, \og\ :\fg{}). Deux amendements,
écrits avant toute vie de test, ont fixé l'exclusion des dimensions massives et
le rejeu exact des paires (sans lui, seules des paires \og calme \fg{}$\to$\og tu
cours \fg{} étaient retenues et $\dN$ était nulle). Réplique : 0,005.

\begin{table}[t]
\centering
\caption{Test 2 (protocole \code{727ef80}). Bloc 12, quatre tokens, 1\,553
contextes, 256 vies de test.}
\label{tab-t2}
\small
\begin{tabularx}{\textwidth}{@{}l >{\raggedright\arraybackslash}p{3.2cm} X l@{}}
\toprule
& Prédiction & Mesure & Verdict \\
\midrule
IA2 & L'état rassemblé cause l'action & énergie basse : $P(\mathrm R)$ 0,578 $\to$ 0,752, $+0,173$ [0,163 ; 0,184] ; nourriture basse : $P(\mathrm M)$ $+0,097$ [0,092 ; 0,103] ; seuil 0,15 ; hasard 0,021 et 0,004 & échoue (nourriture) \\
LS2 & Sans cet état, l'agent meurt & survie 0,680 $\to$ 0,207, baisse 0,473 [0,406 ; 0,539] ; lésion au hasard 0,676 (baisse 0,004) & \passe \\
\midrule
& Critère global (IA2 et LS2) & & non satisfait \\
\bottomrule
\end{tabularx}
\end{table}

\paragraph{Résultats.} Effacer deux directions sur quatre tokens de la ligne du
tour, au bloc~12, fait tomber la survie de 68\,\% à 21\,\%, le niveau du modèle qui
n'a jamais appris (20\,\%) ; un plan tiré au hasard ne change rien
(tableau~\ref{tab-t2}). Avec la lésion, l'agent ne sert plus son besoin le plus bas
que 67 fois sur 100 (91 intact) et meurt surtout d'énergie (189 vies, contre 17
intact). Pousser l'état vers \og énergie basse \fg{} fait recharger l'agent
rassasié (+0,17, au-dessus du seuil), et une poussée aléatoire de même force ne
fait presque rien ; pour la nourriture, l'effet va dans le bon sens mais reste
sous le seuil (+0,10). Le premier test visait l'endroit où le besoin \emph{se lit} ;
celui-ci vise l'endroit où il est \emph{rassemblé et utilisé}, et y trouve un état
dont l'agent dépend pour vivre.

\subsection{Test 3 : dire son besoin, et par le même état}
\label{sec-t3}

\paragraph{Dispositif.} Une exploration déclarée avant le protocole montre que le
besoin n'est rassemblé \emph{que} sur les trois tokens \Choix{} : les remplacer au
bloc~12 restaure 85\,\% de l'effet d'un événement passé, le point qui clôt
l'événement 1\,\%, et rien au bloc~11. L'agent calcule son besoin pour agir, pas
pour répondre. Le test apprend donc à l'agent à répondre (512 vies, 5\,160
questions posées après l'événement, 600 itérations, perte finale 0,20 sur les
réponses), puis mesure des directions \emph{d'action} (sur \Choix{}) et \emph{de
rapport} (sur les trois derniers tokens de la question), et demande si la
direction d'action, injectée dans la question, change ce que l'agent dit.
Dimensions massives exclues : 3 et 35 ; réplique 0,003 ; 1\,418 contextes.

\begin{table}[t]
\centering
\caption{Test 3 (protocole \code{3afa889}). Agent parlant ; 256 vies de test.}
\label{tab-t3}
\small
\begin{tabularx}{\textwidth}{@{}l >{\raggedright\arraybackslash}p{3.6cm} X l@{}}
\toprule
& Prédiction & Mesure & Verdict \\
\midrule
R3 & L'agent dit son besoin & exactitude équilibrée 0,665 (énergie) et 0,655 (nourriture) ; seuil 0,75 & échoue \\
A3 & Il agit toujours & survie 0,652 ; seuil 0,55 & \passe \\
IR3 & Ce qu'il dit est causé par le besoin rassemblé sur la question & $+0,00005$ et $+0,0001$ ; seuil 0,15 & échoue \\
SAME3 & L'état qui fait agir fait aussi dire & $+0,00008$ et $+0,00004$ ; seuil 0,10 & échoue \\
\midrule
& Critère global & & non satisfait \\
\bottomrule
\end{tabularx}
\end{table}

\paragraph{Résultats.} Entraîné à répondre à une question posée à part, l'agent
\emph{devine} son besoin mieux que le hasard (0,66), mais sa réponse ne passe pas
par l'état qui le fait agir (tableau~\ref{tab-t3}) : ni pousser le besoin sur la
question, ni y porter l'état d'action ne change ce qu'il dit. Le cosinus entre la
direction d'action et la direction de rapport vaut 0,41 pour l'énergie et 0,07 pour
la nourriture, et le besoin est codé environ 40 fois plus faiblement sur la ligne
de question que sur \Choix{} (normes par unité de 11 à 18 contre 290 à 860).
L'acte et la parole restent deux chemins séparés ; l'agent répond probablement
d'après des indices tirés directement de l'histoire. Les résultats partiels ont
été lus à 48 puis 152 vies ; rien n'a été changé.

\subsection{Tests 4 et 5 : poser la question après \og Choix : \fg{}, puis l'y forcer}
\label{sec-t45}

Le test~4 pose la question \emph{après} \Choix{}, là où le besoin est rassemblé,
pour pousser un seul état au même endroit et lire deux effets ; le test~5 impose
en plus le masque~\eqref{eq-mask}, pour que la question ne lise la vie qu'à travers
\Choix{}. Dans les deux cas, un même adaptateur apprend à agir et à dire
(600 itérations sur 5\,059 questions des mêmes 512 vies). Ces deux tests sont
\textbf{suspendus} pour libérer le processeur ; leurs verdicts n'ont pas été
calculés et seront publiés tels quels. Deux observations partielles, lues avant
d'écrire les protocoles suivants, sont déclarées dans ceux-ci et rapportées ici
sans valeur de verdict :
\begin{itemize}
\item \textbf{Test 4}, 8 premières vies de test sur 256 : rapport plus juste
(0,74 et 0,72), mais la poussée \og énergie basse \fg{} sur \Choix{} fait agir
(+0,20) sans presque changer ce que l'agent dit (+0,003). Même posée après
\Choix{}, la question semble lire directement l'histoire.
\item \textbf{Test 5}, vies de direction seulement : l'agent à espace de travail
ne survit que 0,359 fois, contre 0,734 pour l'agent du test~4 sur ses propres vies
de direction ; il meurt surtout de faim (72 morts de faim contre 10 d'énergie) et
ne sert son besoin le plus bas que 80 fois sur 100 (88 pour l'autre). Forcée de
passer par \Choix{}, la parole a \textbf{réécrit} l'état qui sert à agir : les
poids appris pour dire servent aussi à calculer \Choix{}.
\end{itemize}
Cette seconde observation a conduit à séparer le calcul de l'état et sa lecture :
c'est le lecteur.

\subsection{Test 6 : le lecteur}
\label{sec-t6}

\paragraph{Dispositif.} L'agent qui agit est l'agent final, inchangé. Le lecteur
(section~\ref{sec-lecteur}) apprend sur le Mac 600 itérations sur les 5\,059
questions des tests~4 et~5, sans document de choix ; perte de validation
0,68 $\to$ 0,29. Les directions sont mesurées \emph{sur l'agent qui agit}, sans
que le lecteur intervienne : 1\,500 paires au bloc~12 sur les trois tokens
\Choix{} (dimension 35 exclue), 128 vies de direction neuves. Mesures sur 256 vies
de test neuves, 1\,523 contextes ; répliques 0,003 (agent) et 0,012 (lecteur) ;
contrôle d'exécution $3\times10^{-5}$. Les résultats partiels ont été lus à 8,
112, 120, 168 et 240 vies ; rien n'a été changé, mais deux protocoles ont été écrits
après ces lectures et le disent (nécessité après 112 vies, réplication après 120).

\begin{table}[t]
\centering
\caption{Test 6 (protocole \code{3b27636}). L'énergie est le critère de ONE6 ; la
nourriture est rapportée sans seuil global.}
\label{tab-t6}
\small
\begin{tabularx}{\textwidth}{@{}l >{\raggedright\arraybackslash}p{3.6cm} X l@{}}
\toprule
& Prédiction & Mesure & Verdict \\
\midrule
R6 & Le lecteur dit le besoin depuis l'état qui fait agir & exactitude équilibrée 0,766 (énergie), 0,705 (nourriture) ; seuil 0,75 pour les deux & échoue (nourriture) \\
A6 & L'agent agit (contrôle) & survie 0,633 & \passe \\
ONE6 & Un seul état cause l'acte et la parole (énergie) & acte $+0,156$ [0,146 ; 0,166] ; parole $+0,202$ [0,192 ; 0,213] ; hasard 0,013 et 0,006 & \passe \\
\midrule
& Critère global & & non satisfait \\
\midrule
& Nourriture (sans seuil) & acte $+0,107$ [0,101 ; 0,113] ; parole $+0,009$ [0,008 ; 0,011] ; hasard 0,006 et 0,001 & --- \\
\bottomrule
\end{tabularx}
\end{table}

\paragraph{Résultats.} Pour l'énergie, \textbf{une même intervention sur un même
état change à la fois l'acte et la parole} (tableau~\ref{tab-t6}) : pousser l'état
rassemblé sur \Choix{} vers \og énergie basse \fg{} fait recharger l'agent rassasié
et fait dire au lecteur \og mon énergie est basse \fg{}, alors que des poussées
aléatoires de même force ne font presque rien. Le résultat est probant parce que le
lecteur ne voit la vie que par cet état et ne peut pas le modifier, et parce que la
direction a été mesurée sur l'agent qui agit, indépendamment du lecteur. Sur les
1\,523 contextes, l'effet sur l'acte est positif dans 100\,\% des cas et celui sur
la parole dans 95\,\% ; pour le hasard, 47\,\%. Exemple d'un contexte réel, agent
rassasié : sans intervention, $P(\mathrm R)=0,497$ et $P(\text{oui})=0,001$ ; avec
$-4\,\dE$, $P(\mathrm R)=0,960$ et $P(\text{oui})=0,128$ ; avec un vecteur aléatoire
de même norme, $P(\mathrm R)=0,575$ et $P(\text{oui})=0,001$.

Pour la nourriture, la poussée fait manger l'agent (+0,107) mais le lecteur ne la
dit pas (+0,009), et l'exactitude reste sous le seuil : R6 échoue, et avec lui le
critère global. Une analyse exploratoire après le verdict montre que le lecteur lit
bien un niveau de nourriture, gradué ($P(\text{oui})$ moyen 0,73 au niveau 1, 0,55
au niveau 4, 0,21 au niveau 8), mais avec une pente bien plus douce que pour
l'énergie (0,67 $\to$ 0,51 $\to$ 0,02), et que sa réponse suit beaucoup l'événement
du tour : après \og il fait froid \fg{} ($-3$ en nourriture) il dit \og oui \fg{} à
0,67 alors que la nourriture n'est basse que dans 42\,\% de ces tours ; après
\og tu cours \fg{}, 0,29 pour 12\,\% de tours bas. L'événement du tour, que
\Choix{} porte aussi, sert de raccourci ; la direction $\dN$, qui vient des
événements passés, est peu utilisée par le lecteur.

\subsection{Test 7 : sans cet état, ni agir ni dire}
\label{sec-t7}

\paragraph{Dispositif.} Même agent et même lecteur qu'au test~6. Au bloc~12, sur
les trois tokens \Choix{}, les coordonnées dans le plan $(\dE,\dN)$ du test~6,
orthonormalisé token par token, sont remplacées par leur moyenne~\eqref{eq-lesion},
prise sur 3\,403 décisions des 128 vies de direction du test~6, rejouées à
l'identique. Témoin : un plan aléatoire par token. 256 mondes neufs, chacun vécu
trois fois ; dans les vies intactes, les deux questions sont lues à chaque décision
sous les trois conditions (6\,618 décisions). Contrôle d'exécution $3\times10^{-7}$.
Lectures partielles déclarées à 8, 64 et 256 vies intactes et 32 vies lésées.

\begin{table}[t]
\centering
\caption{Test 7 (protocole \code{732ed1c}). L'énergie est le critère pour la
parole ; la nourriture est rapportée sans seuil.}
\label{tab-t7}
\small
\begin{tabularx}{\textwidth}{@{}l >{\raggedright\arraybackslash}p{3.6cm} X l@{}}
\toprule
& Prédiction & Mesure & Verdict \\
\midrule
LS7 & Sans l'état, l'agent ne sait plus survivre & survie 0,676 $\to$ 0,203 ; baisse 0,473 [0,410 ; 0,535] ; lésion au hasard : baisse 0,000 & \passe \\
LR7 & Sans l'état, il ne sait plus dire son énergie & exactitude 0,750 $\to$ 0,500 ; baisse 0,250 [0,233 ; 0,268] ; lésion au hasard : baisse 0,000 & \passe \\
\midrule
& Critère global & & \textbf{satisfait} \\
\midrule
& Nourriture (sans seuil) & exactitude 0,745 $\to$ 0,671 ; baisse 0,074 [0,057 ; 0,090] ; hasard 0,000 & --- \\
\bottomrule
\end{tabularx}
\end{table}

\paragraph{Résultats.} Le critère global est satisfait (tableau~\ref{tab-t7}).
Avec le test~6, chez cet agent et pour l'énergie, un même état interne, appris par
la seule satisfaction du besoin, est à la fois \textbf{suffisant} (le pousser fait
agir, $+0,156$, et dire, $+0,202$) et \textbf{nécessaire} (l'effacer fait tomber la
survie de 0,68 à 0,20 et le rapport au niveau du hasard), alors que des
interventions aléatoires de même force ne font rien. Deux précisions descriptives,
calculées pour cet article à partir des vies publiées, éclairent la forme de cette
nécessité. Sous la lésion, le lecteur ne répond plus jamais \og oui \fg{} à la
question sur l'énergie (0 lecture sur 6\,618 au-dessus de 0,5 ; $P(\text{oui})$
moyen 0,027, contre 0,281 sans lésion) : l'exactitude de 0,500 est celle d'une
réponse constante \og non \fg{}, ce qu'on attend quand l'état est figé à sa moyenne,
qui correspond à une énergie plutôt haute. Et l'agent lésé meurt surtout
d'énergie (197 vies sur 256, contre 20 intact), alors que les morts de faim
baissent (7 contre 63).

\begin{figure}[t]
\centering
\begin{tikzpicture}
\begin{axis}[ybar=2pt,width=\textwidth,height=5.8cm,bar width=12pt,
  ylabel={survie},ymin=0,ymax=1.0,ytick={0,0.2,0.4,0.6,0.8},
  symbolic x coords={T1b6,T2b12,T7b12,RLb12,R2b12},
  xticklabels={{test 1\\(bloc 6)},{test 2\\(bloc 12)},{test 7\\(bloc 12)},{second agent,\\réplication (bloc 12)},{second agent,\\test 10 (bloc 12)}},
  x tick label style={align=center,font=\footnotesize},
  xtick=data,enlarge x limits=0.12,
  legend style={at={(0.5,1.03)},anchor=south,legend columns=3,/tikz/every even column/.append style={column sep=0.5cm}},
  nodes near coords,nodes near coords style={font=\scriptsize,rotate=90,anchor=west,inner sep=2pt,/pgf/number format/fixed,/pgf/number format/fixed zerofill,/pgf/number format/precision=3},
  grid=major,grid style={gray!15}]
\addplot[fill=cact!70,draw=cact] coordinates {(T1b6,0.684)(T2b12,0.680)(T7b12,0.676)(RLb12,0.570)(R2b12,0.531)};
\addplot[fill=csay!70,draw=csay] coordinates {(T1b6,0.660)(T2b12,0.207)(T7b12,0.203)(RLb12,0.043)(R2b12,0.035)};
\addplot[fill=cctl!50,draw=cctl] coordinates {(T1b6,0.688)(T2b12,0.676)(T7b12,0.676)(RLb12,0.563)(R2b12,0.531)};
\legend{intact,lésion du plan du besoin,lésion d'un plan aléatoire}
\end{axis}
\end{tikzpicture}
\caption{Survie sous lésion. Au bloc~6 (test~1), effacer le plan corrélationnel
ne change rien ; au bloc~12, sur la fin de la ligne du tour, effacer le plan causal
fait tomber la survie au niveau du modèle de base (tests 2 et 7) ou presque à zéro
chez le second agent (réplication, RLS ; et mesure descriptive du test~10,
section~\ref{sec-locate}, sur d'autres mondes). Les lésions aléatoires de même
dimension ne changent rien (0,563 contre 0,570 pour la réplication).}
\label{fig-lesions}
\end{figure}

\subsection{Test 8 : un lecteur sans raccourci}
\label{sec-t8}

L'analyse exploratoire du test~6 suggérait que le lecteur avait appris un
raccourci (l'événement du tour prédit la réponse). Le test~8 construit des
documents où, pour chaque besoin et chaque événement, il y a autant de décisions
basses que hautes : l'événement ne dit alors rien de la réponse, et \og oui \fg{}
et \og non \fg{} sont en nombre égal (7\,356 documents, dont 3\,678 \og oui \fg{}).
Tout le reste est identique au test~6 ; les prédictions sont R8 (exactitude
$\ge0,75$ pour les deux besoins), SAY8 ($-4\,\dN$ augmente $P(\text{oui})$ à la
question sur la nourriture d'au moins 0,10) et ONE8 (comme ONE6).

\textbf{Le lecteur équilibré n'a pas appris} : sa perte de validation passe de 0,691
à 0,720 en 600 itérations, au niveau de $\ln 2\approx0,693$, celui d'une réponse au
hasard sur des classes équilibrées ; sa perte d'apprentissage moyenne reste entre
0,64 et 0,77 (figure~\ref{fig-losses}). Aucune mesure n'a été publiée, et les verdicts
R8, SAY8 et ONE8 ne sont pas établis. Ce résultat négatif de l'apprentissage est
discuté en section~\ref{sec-echecs}.

\subsection{Réplication sur un second agent appris de zéro}
\label{sec-repl}

\paragraph{Le second agent.} Tout le chemin du premier agent est refait avec les
mêmes réglages ; seul change le hasard (flux des mondes et des choix $100+s$,
graine d'apprentissage 270927). Survie des tours d'apprentissage : 0,17 $\to$ 0,70
(figure~\ref{fig-rounds}). Le premier apprentissage du rapport a échoué une fois
sur une panne du GPU du Mac, et a été relancé à l'identique. Son lecteur apprend
peu : perte de validation 0,74 $\to$ 0,66 (premier agent : 0,68 $\to$ 0,29).
Mesures en torch : répliques 0,006 (agent) et 0,013 (lecteur), contrôle d'exécution
$3\times10^{-6}$, 1\,260 contextes ; la machine a redémarré une fois et la mesure a
repris après 152 vies ; aperçus déclarés à 144 et 216 vies. Les critères portent sur
l'énergie, choix fait en connaissant les lectures du test~6 et dit dans le
protocole.

\begin{table}[t]
\centering
\caption{Réplication (protocole \code{4e10a6c}), second agent, au bloc 12. La
réplication est jugée seule. RLS et RLR : 256 mondes neufs (flux 128 ; plan au
hasard, flux 129).}
\label{tab-repl}
\small
\begin{tabularx}{\textwidth}{@{}l >{\raggedright\arraybackslash}p{3.6cm} X l@{}}
\toprule
& Prédiction & Mesure & Verdict \\
\midrule
RR & Le lecteur dit son énergie & 0,623 ; seuil 0,75 & échoue \\
RA & Le second agent agit & survie 0,512 ; seuil 0,55 & échoue \\
RONE & Un seul état fait agir et dire & acte $+0,034$ [0,031 ; 0,036] ; parole $+0,002$ [0,002 ; 0,003] ; hasard 0,003 et 0,001 & échoue \\
RLS & Sans l'état, il ne survit plus & survie 0,570 $\to$ 0,043 ; baisse 0,527 [0,465 ; 0,590] ; lésion au hasard : baisse 0,008 & \passe \\
RLR & Sans l'état, il ne dit plus son énergie & exactitude 0,639 $\to$ 0,500 ; baisse 0,139 [0,115 ; 0,164] ; hasard 0,001 & \passe \\
\midrule
& Critère global & & non satisfait \\
\bottomrule
\end{tabularx}
\end{table}

\paragraph{Suffisance.} Chez le second agent, la poussée de l'état au bloc~12 sur
\Choix{} fait agir dans le bon sens, mais environ cinq fois moins que chez le premier
($+0,034$ contre $+0,156$), et elle ne fait rien dire (tableau~\ref{tab-repl}). Deux
explications étaient possibles : le second agent rassemble son besoin ailleurs, ou
son lecteur n'a pas appris ; le test~10 teste la première, le test~11 la seconde.

\paragraph{Nécessité.} Mesurée ensuite sur 256 mondes neufs avec les directions et
les vies de direction du second agent, comme au test~7. La mesure a été mise en
pause à 176 vies intactes pour laisser le processeur au test~11, puis reprise ; une
copie lancée par erreur au redémarrage du conteneur a tourné une demi-heure en
parallèle avant d'être arrêtée, sans avoir rien écrit. Effacer le plan $(\dE,\dN)$
au bloc~12 tue presque toujours le second agent (survie 0,570 $\to$ 0,043) et ramène
le rapport d'énergie de son lecteur au niveau du hasard (0,639 $\to$ 0,500) ; un plan
aléatoire ne change rien. La nourriture, sans seuil, suit (0,706 $\to$ 0,502).

\paragraph{Ce que cela dit.} La \textbf{nécessité se réplique} : chez le second agent
aussi, un état sur \Choix{} au bloc~12 est nécessaire pour survivre et pour dire son
énergie. La \textbf{suffisance ne se réplique pas} sous la forme mesurée : pousser
cet état le long de $\dE$ fait peu agir et ne fait pas dire. Le critère global de la
réplication reste non satisfait (RR, RA, RONE).

\subsection{Test 9 : à qui est ce besoin ?}
\label{sec-t9}

\paragraph{Le monde à deux.} Un autre agent vit à côté ; ses événements, tirés de la
même loi mais indépendamment, sont écrits sur la même ligne, avec les mêmes mots
pour \og calme \fg{} et \og orage \fg{} :
\begin{center}\small\ttfamily Tour 5 : tu cours. L'autre : il court. Choix :\end{center}
Ils ne changent rien aux besoins de l'agent ; rien dans le texte ne dit \og ceci est
à toi \fg{}, sinon \og tu \fg{} et la place dans la ligne. Un nouvel agent y apprend
par la seule satisfaction de \emph{ses} besoins (flux $200+s$, graine 270928) :
survie 0,11 $\to$ 0,71 au tour~8 (la première demande avait subi une panne du GPU).
Son lecteur, actif sur la seule question et sous le masque, a peu appris (perte de
validation finale 0,78). Mesures sur le processeur du Mac, en cinq tranches ;
répliques 0,003 et 0,009, contrôle d'exécution $2\times10^{-6}$ ; 1\,500 paires de
l'agent et 1\,500 de l'autre ; dimensions massives 35 et 277 ; 256 vies de test,
1\,547 contextes. Les paires ont été lues dès la fin de la deuxième tranche : les
verdicts MINE9 et SELF9 étaient connus avant les vies de test (déclaré).

\begin{table}[t]
\centering
\caption{Test 9 (protocole \code{47eb0db}), agent du monde à deux, bloc 12. Rapports
calculés sur les paires \og calme \fg{}$\to$\og orage \fg{} (552 paires de l'agent,
607 de l'autre).}
\label{tab-t9}
\small
\begin{tabularx}{\textwidth}{@{}l >{\raggedright\arraybackslash}p{3.2cm} X l@{}}
\toprule
& Prédiction & Mesure & Verdict \\
\midrule
S9 & L'agent vit dans le monde à deux & survie 0,742 ; seuil 0,55 & \passe \\
MINE9 & L'état de besoin est le sien & effet de l'autre / effet de l'agent : 0,94 sur l'état, 0,57 sur l'acte ; exigé $\le0,33$ & échoue \\
SELF9 & Sa parole est la sienne & même rapport sur la parole : 0,78 & échoue \\
ONE9 & Un seul état fait agir et dire & acte $+0,090$ [0,085 ; 0,095] ; parole $+0,014$ [0,012 ; 0,016] ; hasard 0,006 et 0,0004 & échoue \\
\midrule
& Critère global & & non satisfait \\
\bottomrule
\end{tabularx}
\end{table}

\paragraph{Résultats et analyse exploratoire.} La possession n'est pas démontrée
(tableau~\ref{tab-t9}). Exactitude du rapport : énergie 0,61, nourriture 0,71. Une
analyse faite après la lecture des paires, échange par échange, nuance ce verdict
(tableau~\ref{tab-t9x}). Les changements de l'agent bougent l'état selon ce qu'ils font
à \emph{ses} besoins ; ceux de l'autre produisent une perturbation presque
constante (269 à 280), quel que soit l'événement. Sur l'acte, l'effet de l'autre vaut
25\,\% à 57\,\% de celui de l'agent : l'agent distingue donc en partie ce qui lui
arrive. Mais la comparaison prévue était défavorable : \og calme \fg{}$\to$\og orage \fg{}
retire une unité aux deux besoins de l'agent, ce qui change à peine leur équilibre,
qui est ce que l'état code (section~\ref{sec-axe}) ; nous l'avons appris après avoir
écrit le protocole. Le verdict reste celui du protocole.

\begin{table}[t]
\centering
\caption{Test 9, analyse exploratoire : norme moyenne du changement dans le plan
du besoin (état) et $|\Delta P(\mathrm R)|$ moyen (acte), par type d'échange.}
\label{tab-t9x}
\small
\begin{tabular}{@{}l cc cc@{}}
\toprule
& \multicolumn{2}{c}{Échange chez l'agent} & \multicolumn{2}{c}{Échange chez l'autre} \\
\cmidrule(lr){2-3}\cmidrule(lr){4-5}
Échange & état & acte & état & acte \\
\midrule
\og calme \fg{} $\to$ \og tu cours \fg{} / \og il court \fg{} & 607 & 0,075 & 279 & 0,019 \\
\og tu te reposes \fg{} $\to$ \og il fait froid \fg{} / \og il a froid \fg{} & 442 & 0,043 & 280 & 0,022 \\
\og calme \fg{} $\to$ \og orage \fg{} (mêmes mots) & 286 & 0,028 & 269 & 0,016 \\
\bottomrule
\end{tabular}
\end{table}

% ================================================================= 6
\section{La localisation, le test 11 et ce que porte l'état}
\label{sec-locate}

\subsection{Test 10 : où chaque agent rassemble-t-il son besoin ?}
\label{sec-t10}

\paragraph{Pourquoi.} Le bloc~12 avait été trouvé par un patching
\emph{exploratoire} sur le premier agent, puis imposé à tous. Deux aperçus,
déclarés avant ce protocole, montraient chez les autres agents des effets faibles
au bloc~12 : second agent, 216 vies sur 256, acte $+0,033$ et parole $+0,002$ ;
agent du monde à deux, 64 vies, acte $+0,079$ et parole $+0,013$. Hypothèse : un
agent appris avec une autre graine peut rassembler son besoin à un autre bloc.

\paragraph{Dispositif.} La règle~\eqref{eq-share} est appliquée identiquement à
chaque agent, sur ses vies de direction et ses paires déjà publiées. Elle doit
retrouver $b^\star=12$ chez le premier agent (étalonnage). Au bloc trouvé, chaque
autre agent est testé comme aux tests~6 et~7 : directions ajustées sur ses
1\,500 paires à $b^\star$, injection $-4\,\dE$ sur 256 vies de test neuves (flux
$300+s$ pour le second agent, $400+s$ pour l'agent du monde à deux), lésion du plan
$(\dE,\dN)$ et lésion aléatoire. Prédictions, pour chacun des deux agents :
\textbf{LOC10} (un bloc $b^\star$ existe, établi sur 100 paires), \textbf{ACT10}
($-4\,\dE$ augmente $P(\mathrm R)$ d'au moins 0,15, borne basse $>0$, hasard au plus
le tiers), \textbf{LS10} (baisse de survie sous lésion d'au moins 0,15, borne basse
$>0$, lésion aléatoire au plus le tiers). Critère global : les trois pour les deux
agents, avec une règle étalonnée. La parole est mesurée sans seuil, car elle dépend
d'un lecteur qui a mal appris chez ces deux agents.

\begin{figure}[t]
\centering
\begin{tikzpicture}
\begin{axis}[width=0.78\textwidth,height=5.6cm,xlabel={bloc $b$},ylabel={$\operatorname{part}(b)$},
  xmin=-0.5,xmax=27.5,ymin=-0.05,ymax=1.5,xtick={0,3,...,27},ytick={0,0.5,1},
  legend pos=north west,legend cell align=left,grid=major,grid style={gray!20}]
\addplot[cact,thick,mark=*,mark size=1.4pt] coordinates {(0,0.000)(1,0.000)(2,0.000)(3,0.000)(4,-0.003)(5,-0.003)(6,-0.001)(7,-0.001)(8,0.001)(9,0.004)(10,0.004)(11,0.004)(12,0.677)(13,0.678)(14,0.678)(15,0.678)(16,0.679)(17,0.679)(18,0.679)(19,0.928)(20,0.928)(21,0.929)(22,0.965)(23,0.966)(24,1.008)(25,1.008)(26,0.993)(27,1.000)};
\addlegendentry{premier agent : $b^\star=12$ (100 paires)}
\addplot[cgrn,thick,mark=triangle*,mark size=1.6pt] coordinates {(0,0.001)(1,0.001)(2,-0.001)(3,-0.002)(4,-0.003)(5,-0.008)(6,-0.007)(7,-0.008)(8,0.010)(9,0.017)(10,0.010)(11,0.008)(12,0.625)(13,0.625)(14,0.626)(15,0.625)(16,0.623)(17,0.623)(18,0.623)(19,0.623)(20,0.623)(21,0.623)(22,0.648)(23,0.648)(24,0.647)(25,0.646)(26,0.920)(27,1.000)};
\addlegendentry{second agent : $b^\star=12$ (71 paires)}
\addplot[corg,thick,mark=diamond*,mark size=1.6pt] coordinates {(0,0.000)(1,0.000)(2,0.000)(3,0.000)(4,0.000)(5,0.000)(6,0.001)(7,0.001)(8,0.001)(9,0.002)(10,0.002)(11,0.002)(12,0.348)(13,0.348)(14,0.348)(15,0.349)(16,0.349)(17,0.349)(18,0.349)(19,0.352)(20,0.352)(21,0.363)(22,0.877)(23,0.877)(24,0.872)(25,0.873)(26,1.000)(27,1.000)};
\addlegendentry{monde à deux : $b^\star=22$ (100 paires)}
\addplot[black,dashed,domain=-0.5:27.5] {0.5};
\end{axis}
\end{tikzpicture}
\caption{Règle de localisation (test~10) : part de l'effet d'un \og calme \fg{}
changé en \og tu cours \fg{} restaurée en remplaçant la sortie du bloc $b$ sur les
trois tokens \Choix{} (données de \code{locate.json}). La ligne pointillée est le
seuil 0,5. Rien ne passe avant le bloc~12 ; le rassemblement se fait ensuite par
paliers, différents selon l'agent.}
\label{fig-locate}
\end{figure}

\paragraph{Étalonnage : réussi.} Chez le premier agent, la règle retrouve
$b^\star=12$ (100 paires retenues sur 550 essayées) : la part restaurée vaut 0,004
au bloc~11 et 0,677 au bloc~12, puis monte par paliers (0,928 au bloc~19, 0,965 au
bloc~22, 1 au bloc~24 ; figure~\ref{fig-locate}). La règle est donc valide au sens
du protocole.

\paragraph{Second agent : LOC10 échoue.} La règle désigne aussi le bloc~12
(part 0,625, qui reste entre 0,62 et 0,65 jusqu'au bloc~25 puis passe à 0,920 au
bloc~26), mais sur \textbf{71 paires seulement} : ses 1\,500 paires ne comptent que
454 échanges \og calme \fg{}$\to$\og tu cours \fg{}, et 71 d'entre eux ont un effet
d'au moins 0,10 sur l'acte. Le protocole exige 100 paires ; par le code
pré-enregistré, LOC10 échoue pour le second agent, et avec lui le critère global du
test. Les mesures faites au bloc~12 sont donc \textbf{descriptives}
(tableau~\ref{tab-t10}) : la poussée $-4\,\dE$ fait agir de $+0,036$ [0,033 ; 0,039]
(hasard 0,0015) et dire de $+0,003$ ; la lésion du plan fait tomber la survie de
0,531 à 0,035 (baisse 0,496 [0,434 ; 0,555]), alors que la lésion d'un plan aléatoire
la laisse à 0,531 ; l'exactitude du rapport passe de 0,614 à 0,500 pour l'énergie et
de 0,694 à 0,502 pour la nourriture. Calculé pour cet article à partir des vies
publiées : sous la lésion, le lecteur ne dit presque jamais \og oui \fg{} (aucune
lecture au-dessus de 0,5 pour l'énergie, une sur mille pour la nourriture), et
l'agent meurt surtout de faim (235 vies sur 256, contre 96 intact).

\paragraph{Agent du monde à deux : $b^\star=22$.} La part restaurée reste près de
0,35 du bloc~12 au bloc~21, puis passe à 0,877 au bloc~22 (100 paires sur 643
essayées). Les directions sont ajustées à ce bloc (dimension 35 exclue). Sur 256
vies de test (1\,585 contextes ; répliques 0,003 et 0,009, contrôle d'exécution
$2,5\times10^{-6}$), la poussée $-4\,\dE$ au bloc~22 fait agir de
$+0,217$ [0,206 ; 0,230], contre 0,011 pour le hasard : \textbf{le critère ACT10 est
atteint par les valeurs mesurées}. Elle ne fait pas dire : $-0,001$
[$-0,0019$ ; $-0,0004$]. Exactitude du rapport : énergie 0,608, nourriture 0,712 ;
survie 0,707 ; pour la nourriture, l'acte suit ($+0,125$ [0,117 ; 0,133]) mais pas
la parole ($+0,007$). Les bras de lésion (LS10) tournent sur le Mac (168 vies
intactes sur 256 au 2 octobre). Une tranche du Mac, la sixième, a été perdue : le
code de reprise revivait des vies de test déjà finies dans une tranche précédente ;
il a été corrigé (commit \code{de2c1b4}) et aucun résultat n'en est changé. Un
aperçu de 112 vies de ces mesures avait été lu avant l'écriture du protocole du
test~11, ce que ce protocole déclare.

\begin{table}[t]
\centering
\caption{Test 10 (protocole \code{12b4d2d}). Statut au 2 octobre 2026. Les
mesures du second agent sont descriptives (LOC10 échoue).}
\label{tab-t10}
\small
\begin{tabularx}{\textwidth}{@{}l X X@{}}
\toprule
& Second agent & Agent du monde à deux \\
\midrule
$b^\star$ ; paires retenues & 12 ; 71 sur 454 essayées & 22 ; 100 sur 643 essayées \\
LOC10 & \textbf{échoue} (moins de 100 paires) & \passe \\
Acte, $-4\,\dE$ à $b^\star$ & $+0,036$ [0,033 ; 0,039] ; hasard 0,0015 (descriptif) & $+0,217$ [0,206 ; 0,230] ; hasard 0,011 : ACT10 atteint \\
Parole, $-4\,\dE$ à $b^\star$ & $+0,003$ (descriptif) & $-0,001$ [$-0,0019$ ; $-0,0004$] (sans seuil) \\
Survie intacte $\to$ lésée (hasard) & 0,531 $\to$ 0,035 (0,531) ; baisse 0,496 [0,434 ; 0,555] (descriptif) & en cours sur le Mac (LS10 ; 168 vies intactes sur 256) \\
Exactitude énergie, intacte $\to$ lésée & 0,614 $\to$ 0,500 (descriptif) & en cours \\
Exactitude nourriture, intacte $\to$ lésée & 0,694 $\to$ 0,502 (descriptif) & en cours \\
\midrule
Critère global & \multicolumn{2}{l}{non satisfait (LOC10 du second agent), étalonnage réussi} \\
\bottomrule
\end{tabularx}
\end{table}

\paragraph{Ce que cela dit.} Le besoin n'est pas rassemblé au même bloc, ni de la
même façon, chez trois agents appris avec les mêmes réglages. Chez les trois, rien
ne passe sur \Choix{} avant le bloc~12, le premier que l'apprentissage modifie ;
ensuite, le rassemblement se fait par paliers, et le seuil de 0,5 tombe au
bloc~12 chez deux agents et au bloc~22 chez le troisième. Chez l'agent du monde à
deux, pousser l'état au bloc~22 fait agir davantage que chez le premier agent au
bloc~12 ($+0,217$ contre $+0,156$). Chez le second agent, l'état au bloc~12 est
nécessaire à la survie (descriptivement plus encore que chez le premier) alors que
le pousser fait peu agir : nécessité et suffisance ne vont pas forcément ensemble,
et une direction ajustée sur peu de paires utiles peut mal estimer le levier. Le
résultat des tests~6 et~7 n'est pas, à ce stade, une propriété générale des agents
appris de cette façon.

\subsection{Test 11 : un lecteur qui apprend assez}
\label{sec-t11}

\paragraph{Pourquoi.} Les lecteurs du second agent et de l'agent du monde à deux,
appris exactement comme celui du premier (600 itérations), ont peu appris (perte
de validation 0,66 et 0,78 ; exactitude pour l'énergie 0,62 et 0,61), et la
poussée de l'état ne change presque rien à leur parole. Le protocole
(\code{0a28560}, 1\ier{} octobre) a été écrit après avoir lu les blocs trouvés par
la règle (12 avec 71 paires, 22) et l'aperçu de 112 vies de l'agent du monde à
deux au bloc~22 (acte $+0,22$, parole $-0,001$). Hypothèse : ces lecteurs ont
simplement trop peu appris.

\paragraph{Dispositif.} Pour chaque agent, tout est repris du lecteur précédent
(agent fondu, mêmes 512 vies et mêmes documents : 5\,091 questions pour le second
agent, 5\,222 pour l'agent du monde à deux ; adaptateur actif sur la seule question,
masque, lots de 4, taux $10^{-4}$, même graine), \textbf{sauf le nombre
d'itérations : 2\,000 au lieu de 600}. Mesures en torch au bloc de chaque agent
(12 et 22), avec les directions ajustées à ce bloc par le test~10 ; 256 vies de test
neuves (flux 526 pour les vies de test et 527 pour les directions aléatoires
du second agent ; 626 et 627 pour l'agent du monde à deux) ; réplique du lecteur $\le0,02$, contrôle
d'exécution $\le10^{-4}$, au moins 100 contextes, masses $\ge0,5$.

\paragraph{Prédictions.} Pour chacun des deux agents :
\textbf{R11}, le lecteur long dit son énergie (exactitude équilibrée $\ge0,75$) ;
\textbf{SAY11}, il la dit par l'état qui fait agir ($-4\,\dE$ au bloc de l'agent
augmente $P(\text{oui})$ d'au moins 0,10, borne basse $>0$, directions aléatoires au
plus le tiers). Critère global : R11 et SAY11 pour les deux agents. L'effet sur
l'acte et les mesures de la nourriture sont publiés sans seuil. Si R11 passe et SAY11 échoue, le
lecteur aura appris à dire par un autre chemin ; si R11 échoue, même plus
d'apprentissage ne suffit pas.

\paragraph{Amendement avant exécution : deux tranches exactes.} Une construction du
Mac dure au plus 120 minutes, et 2\,000 itérations au rythme mesuré (0,28 itération
par seconde) en prendraient 120 à elles seules. Chaque lecteur long apprend donc en
deux tranches de 1\,000 itérations, la seconde \emph{continuant exactement} la
première. (i)~Mêmes lots dans le même ordre : la seconde tranche repart de la même
graine et passe les 1\,000 premiers lots ; les permutations sont tirées dans le
même ordre qu'un apprentissage d'un seul tenant, parce que la reprise tombe dans la
première époque (1\,272 lots par époque pour le second agent, 1\,305 pour l'agent du
monde à deux) ; le code refuse une reprise hors de la première époque. (ii)~Même
état de l'optimiseur : les moments d'Adam sont enregistrés à la fin de la première
tranche et relus au début de la seconde ; Adam étant ici sans correction de biais et
à taux constant, ses moments
$m_t=\beta_1m_{t-1}+(1-\beta_1)g_t$ et $v_t=\beta_2v_{t-1}+(1-\beta_2)g_t^2$ sont tout
son état. (iii)~Vérifié avant l'emploi : un test compare 3\,+\,3 itérations reprises à
6 itérations d'un seul tenant, à travers une fin d'époque ; il passe sur un petit
modèle et sur le Mac avec le vrai modèle.

\paragraph{Contrôles d'apprentissage (descriptifs).} Les 600 premières itérations
étant, par construction, celles des premiers lecteurs (même graine, mêmes lots),
leurs pertes d'apprentissage journalisées doivent coïncider : sur les 60 points
journalisés (toutes les 10 itérations), l'écart maximal est \textbf{0,000} pour les
deux agents. Les pertes de validation, calculées sur 8 documents seulement et donc
seulement indicatives, sont données au tableau~\ref{tab-t11loss} ; les pertes
d'apprentissage moyennes par fenêtres de 100 itérations, à la
figure~\ref{fig-losses}. Pour le second agent, les contrôles de la mesure sont
faits : contrôle d'exécution $2,6\times10^{-6}$, réplique du lecteur long 0,013.

\begin{table}[t]
\centering
\caption{Pertes de validation des lecteurs (8 documents ; descriptif). Le point à
600 itérations est celui du premier lecteur de chaque agent.}
\label{tab-t11loss}
\small
\begin{tabular}{@{}l cccc@{}}
\toprule
Lecteur & itération 1 & 600 & 1\,000 & 2\,000 \\
\midrule
Premier agent (test 6) & 0,680 & 0,293 & --- & --- \\
Lecteur équilibré (test 8) & 0,691 & 0,720 & --- & --- \\
Second agent, lecteur long & 0,741 & 0,655 & 0,555 & 0,431 \\
Agent du monde à deux, lecteur long & 19,688 & 0,779 & 0,939 & 0,698 \\
\bottomrule
\end{tabular}
\end{table}

\begin{figure}[t]
\centering
\begin{tikzpicture}
\begin{axis}[width=0.62\textwidth,height=5.5cm,xlabel={itération},ylabel={perte (moyenne sur 100 it.)},
  xmin=0,xmax=2050,ymin=0.4,ymax=0.85,
  legend pos=outer north east,legend cell align=left,grid=major,grid style={gray!20},
  scaled x ticks=false,xticklabel style={/pgf/number format/1000 sep={\,}}]
\addplot[cact,thick,mark=*,mark size=1.3pt] coordinates {(100,0.703)(200,0.651)(300,0.544)(400,0.540)(500,0.541)(600,0.509)};
\addlegendentry{premier agent (test 6)}
\addplot[cctl,thick,mark=square*,mark size=1.3pt] coordinates {(100,0.771)(200,0.689)(300,0.678)(400,0.670)(500,0.639)(600,0.660)};
\addlegendentry{lecteur équilibré (test 8)}
\addplot[cgrn,thick,mark=triangle*,mark size=1.5pt] coordinates {(100,0.665)(200,0.645)(300,0.598)(400,0.567)(500,0.558)(600,0.558)(700,0.511)(800,0.514)(900,0.503)(1000,0.503)(1100,0.467)(1200,0.461)(1300,0.534)(1400,0.469)(1500,0.458)(1600,0.544)(1700,0.456)(1800,0.465)(1900,0.437)(2000,0.487)};
\addlegendentry{second agent, lecteur long}
\addplot[corg,thick,mark=diamond*,mark size=1.5pt] coordinates {(200,0.723)(300,0.625)(400,0.622)(500,0.559)(600,0.602)(700,0.567)(800,0.562)(900,0.542)(1000,0.491)(1100,0.526)(1200,0.527)(1300,0.459)(1400,0.486)(1500,0.503)(1600,0.510)(1700,0.516)(1800,0.452)(1900,0.483)(2000,0.460)};
\addlegendentry{monde à deux, lecteur long}
\addplot[black,dotted,thick,domain=0:2050] {0.693};
\addlegendentry{$\ln 2$}
\draw[black!50,dashed] (axis cs:1000,0.4) -- (axis cs:1000,0.85);
\end{axis}
\end{tikzpicture}
\caption{Pertes d'apprentissage des lecteurs, moyennes par fenêtres de 100
itérations, calculées pour cet article à partir des journaux publiés (la première
fenêtre de l'agent du monde à deux, 1,466, est hors cadre). La ligne verticale
marque la reprise exacte entre les deux tranches des lecteurs longs. Le lecteur
équilibré reste près de $\ln2$.}
\label{fig-losses}
\end{figure}

\paragraph{Résultats.} Les deux mesures sont terminées (second agent : 256 vies,
flux 526, en local ; agent du monde à deux : 256 vies, flux 626, sur le processeur
du Mac en trois tranches) et leurs verdicts sont vérifiés en intégration continue
(tableau~\ref{tab-t11}). Les deux mesures sont valides. \textbf{Le critère global
n'est pas satisfait} : R11 et SAY11 échouent pour les deux agents.

% RESULTS-TEST11 (rempli le 2 octobre 2026 depuis artifacts/llm-need/long/verdicts.json)
\begin{table}[t]
\centering
\caption{Test 11 (protocole \code{0a28560}) : lecteurs longs (2\,000 itérations),
256 vies de test neuves par agent, intervalles bootstrap à 95\,\% par vie. Entre
parenthèses, la valeur obtenue avec le lecteur court (600 itérations).}
\label{tab-t11}
\small
\begin{tabularx}{\textwidth}{@{}l >{\raggedright\arraybackslash}p{3.3cm} >{\raggedright\arraybackslash}X >{\raggedright\arraybackslash}X@{}}
\toprule
& Prédiction & Second agent (bloc 12) & Agent du monde à deux (bloc 22) \\
\midrule
R11 & Exactitude équilibrée (énergie) $\ge0,75$ & 0,595 (0,623) : \textbf{échoue} & \textbf{0,747} (0,608) : \textbf{échoue} de 0,003 \\
SAY11 & $-4\,\dE$ : $P(\text{oui})$ $+0,10$ au moins ; hasard $\le$ 1/3 & $+0,0015$ [0,0006 ; 0,0025] ; hasard 0,0009 : \textbf{échoue} & $+0,029$ [0,028 ; 0,030] ($-0,001$) ; hasard 0,001 : \textbf{échoue} \\
--- & Acte, $-4\,\dE$ (sans seuil) & $+0,032$ [0,030 ; 0,034] ; hasard 0,003 & $+0,226$ [0,215 ; 0,237] ; hasard 0,007 \\
--- & Nourriture (sans seuil) & exactitude 0,729 ; acte $+0,029$ ; parole $+0,023$ & exactitude 0,721 ; acte $+0,132$ ; parole $-0,005$ \\
--- & Validité & répliques 0,006 et 0,013 ; exécution $2,6\times10^{-6}$ ; 1\,320 contextes ; masses 0,98 et 1,00 & répliques 0,003 et 0,009 ; exécution $3\times10^{-6}$ ; 1\,549 contextes ; masses 1,00 \\
\midrule
& Critère global (R11 et SAY11, deux agents) & \multicolumn{2}{l}{\textbf{non satisfait}} \\
\bottomrule
\end{tabularx}
\end{table}
% END RESULTS-TEST11

\paragraph{Ce que cela dit.} Chez le second agent, l'hypothèse \og son lecteur avait
trop peu appris \fg{} est réfutée : trois fois plus d'apprentissage ne lui fait pas
dire son énergie (0,623 $\to$ 0,595), ni par l'état qui fait agir. Chez l'agent du
monde à deux, plus d'apprentissage aide nettement : le lecteur dit son énergie à
0,747 au lieu de 0,608, juste sous le seuil, et pousser l'état qui fait agir change
maintenant un peu ce qu'il dit ($+0,029$, environ trente fois le hasard, contre
$-0,001$ avec le lecteur court). Le chemin de la parole par l'état existe chez lui,
mais il est faible, loin du seuil fixé (0,10) et de ce que fait le premier agent
($+0,202$). La parole par l'état qui fait agir reste donc, au niveau des seuils
fixés, propre au premier agent.

\subsection{Analyse exploratoire : le niveau d'énergie est-il dans l'état ?}
\label{sec-probe}

\paragraph{Question et déclaration.} Après l'échec du test~11 chez le second agent,
une hypothèse s'imposait : cet agent ne rassemblerait que l'équilibre $E-N$ qui
décide, et non le niveau de son énergie, si bien qu'aucun lecteur ne pourrait le
dire. L'analyse a été écrite et publiée avant toute lecture
(\code{research/need_probe.py}, commit \code{5ad433f}). Elle n'est pas
pré-enregistrée comme un test et ne porte aucun verdict.

\paragraph{Méthode.} Les 128 vies de direction de chaque agent sont rejouées
exactement (écart de $P(\mathrm R)$ au plus $10^{-6}$ en moyenne) ; la sortie de son bloc
de rassemblement (12, 12 et 22) est capturée sur les trois tokens \Choix{} à chaque
décision (3\,403, 3\,139 et 3\,435 décisions), dimension massive exclue. Des
régressions ridge lisent chaque grandeur ; le $R^2$ est hors échantillon, validé par
vie sur quatre plis, et la régularisation est choisie à l'intérieur de chaque pli
d'apprentissage. \og $E$ au-delà de $E-N$ \fg{} est le $R^2$ du résidu de $E$ après
régression sur $E-N$ : ce que l'état dit de l'énergie une fois l'équilibre connu.

\begin{table}[t]
\centering
\caption{Ce que l'état rassemblé porte (analyse exploratoire, commit \code{5ad433f}).
$R^2$ hors échantillon de régressions ridge validées par vie ; valeurs du fichier
\code{probe/probe.json}.}
\label{tab-probe}
\small
\begin{tabular}{@{}l ccc@{}}
\toprule
Grandeur lue & Premier (bloc 12) & Second (bloc 12) & Monde à deux (bloc 22) \\
\midrule
$E$ & 0,72 & 0,64 & 0,69 \\
$N$ & 0,58 & 0,61 & 0,62 \\
$E-N$ & 0,68 & 0,68 & 0,74 \\
$E+N$ & 0,60 & 0,57 & 0,57 \\
$E$ au-delà de $E-N$ & 0,62 & 0,57 & 0,58 \\
$P(\mathrm R)$ & 0,98 & 0,93 & 0,98 \\
\midrule
\makecell[l]{\og $E\le3$ \fg{} lu depuis $E$ prédit\\(exactitude équilibrée)} & 0,69 & 0,68 & 0,71 \\
\bottomrule
\end{tabular}
\end{table}

\paragraph{Ce que cela dit (exploratoire).} L'hypothèse est \textbf{réfutée}
(tableau~\ref{tab-probe}) : chez les trois agents, l'état porte le niveau d'énergie,
et pas seulement l'équilibre qui décide, avec une précision voisine ; il porte aussi
$E+N$ et ce que $E$ ajoute à $E-N$. La différence de parole entre les agents ne vient
donc pas de l'information disponible dans l'état. Elle vient du lecteur (le premier
atteint 0,77, au-dessus de la lecture linéaire de 0,69 ; le second 0,595, en dessous)
ou du chemin par lequel il lit. Chez le second agent, la lésion du plan rend pourtant
son lecteur aveugle (0,639 $\to$ 0,500, réplication) : il lit bien quelque chose dans
cet état. Ces deux pistes appellent un protocole à part.

% ================================================================= 7
\section{Les échecs et ce qu'ils ont appris}
\label{sec-echecs}

Les échecs de ce programme ne sont pas des accidents à excuser : chacun a été
prévu comme issue possible par son protocole, et la plupart ont orienté le test
suivant. La figure~\ref{fig-inject} rassemble les effets de la même intervention
($-4\,\dE$) chez les trois agents.

\begin{figure}[t]
\centering
\begin{tikzpicture}
\begin{axis}[ybar,width=0.95\textwidth,height=5.8cm,bar width=7pt,
  ylabel={effet de $-4\,d_E$},ymin=-0.02,ymax=0.25,
  symbolic x coords={A,B,C,F,D,E,G},
  xticklabels={{premier\\b12, t.6},{second\\b12, répl.},{second\\b12, t.10},{second\\b12, t.11},{deux\\b12, t.9},{deux\\b22, t.10},{deux\\b22, t.11}},
  xtick=data,enlarge x limits=0.08,x tick label style={font=\scriptsize,align=center},
  legend style={at={(0.5,1.03)},anchor=south,legend columns=4,/tikz/every even column/.append style={column sep=0.4cm}},
  grid=major,grid style={gray!15}]
\addplot[fill=cact!70,draw=cact] coordinates {(A,0.156)(B,0.034)(C,0.036)(F,0.032)(D,0.090)(E,0.217)(G,0.226)};
\addplot[fill=csay!70,draw=csay] coordinates {(A,0.202)(B,0.002)(C,0.003)(F,0.0015)(D,0.014)(E,-0.001)(G,0.029)};
\addplot[fill=cact!20,draw=cact!60] coordinates {(A,0.013)(B,0.003)(C,0.0015)(F,0.003)(D,0.006)(E,0.011)(G,0.007)};
\addplot[fill=csay!20,draw=csay!60] coordinates {(A,0.006)(B,0.001)(C,0.001)(F,0.0009)(D,0.0004)(E,0.0003)(G,0.001)};
\legend{acte $\Delta P(\mathrm R)$,parole $\Delta P(\text{oui})$,acte (hasard),parole (hasard)}
\draw[cact,dashed] (rel axis cs:0,0.6296) -- (rel axis cs:1,0.6296);
\draw[csay,dotted,thick] (rel axis cs:0,0.4444) -- (rel axis cs:1,0.4444);
\end{axis}
\end{tikzpicture}
\caption{La même intervention chez trois agents ; t.11 : lecteurs longs. Tirets :
seuil de l'acte (0,15) ; pointillés : seuil de la parole (0,10). Seul le premier
agent franchit les deux. Chez l'agent du monde à deux, l'état localisé au bloc~22
fait agir plus fort que chez le premier ; avec le lecteur long, il fait un peu
dire ($+0,029$), loin du seuil. Effets aléatoires : moyennes des valeurs absolues
de trois directions de même norme.}
\label{fig-inject}
\end{figure}

\subsection{La réplication au bloc 12}

Le résultat le plus fort du programme (tests 6 et 7) ne se réplique qu'à moitié chez
un second agent appris de zéro avec les mêmes réglages (section~\ref{sec-repl}) : la
nécessité oui (RLS, RLR), la suffisance non (RONE). Quatre enseignements en
découlent.
\begin{itemize}
\item \textbf{Le lieu ne s'importe pas.} Le bloc~12 avait été trouvé par une
exploration sur le premier agent. Le test~10 montre que le rassemblement se fait
par paliers propres à chaque agent ; une règle fixe, étalonnée, appliquée à chacun,
est le seul moyen honnête de comparer des agents.
\item \textbf{Le nombre de paires utiles limite l'estimation.} Le second agent n'a
que 454 paires \og calme \fg{}$\to$\og tu cours \fg{} parmi ses 1\,500 (le premier
en avait 587), et 71 d'entre elles seulement font bouger l'acte d'au moins 0,10. Sa
direction d'énergie se répartit autrement sur les trois tokens (normes 508, 97 et
205 contre 136, 292 et 343 chez le premier ; tableau~\ref{tab-axe}). L'effacer
détruit la survie (0,570 $\to$ 0,043) et le rapport, la pousser fait peu agir :
l'état est nécessaire, mais la direction ajustée n'est pas un bon levier.
\item \textbf{Plus d'apprentissage ne suffit pas.} Un lecteur trois fois plus
entraîné ne dit pas mieux l'énergie du second agent (0,595 ; test~11), alors que
l'information est présente dans l'état ($R^2$ de 0,64 pour $E$ ;
section~\ref{sec-probe}). L'écart avec le premier agent tient au lecteur ou au
chemin par lequel il lit, ce qu'aucun test n'a encore isolé.
\item \textbf{L'apprentissage du rapport a pu abîmer l'action.} Observation
descriptive, tirée des résumés d'étapes publiés : les vies de l'étape de rapport,
vécues par le second agent \emph{avant} l'apprentissage du rapport, survivent
0,688 ; celles de l'étape du lecteur, vécues \emph{après}, 0,525, et ses vies de
test 0,512 (RA échoue). Chez le premier agent, les valeurs correspondantes sont
0,629, 0,664 et 0,684. Les mondes diffèrent d'une étape à l'autre ; ce n'est
qu'une piste, non testée.
\end{itemize}

\subsection{La possession}

MINE9 et SELF9 échouent : changer un événement de l'autre bouge l'état presque
autant (0,94) que changer le même événement de l'agent. Mais l'analyse par échange
(tableau~\ref{tab-t9x}) montre que les événements de l'autre produisent une
perturbation presque constante, alors que ceux de l'agent bougent l'état selon ce
qu'ils font à ses besoins. La comparaison pré-enregistrée portait sur
\og calme \fg{}$\to$\og orage \fg{}, qui retire une unité aux \emph{deux} besoins et
change donc à peine leur équilibre, alors que les directions causales mesurées
suivent surtout cet équilibre (section~\ref{sec-axe}). ONE9 échoue de son côté : l'état y a été cherché au bloc~12, alors que la règle du
test~10 situe le rassemblement de cet agent au bloc~22, et son lecteur a peu appris. Leçons : formuler les comparaisons de possession dans les
termes de ce que l'état code, et localiser avant de tester.

\subsection{La nourriture}

Pour la nourriture, l'état rassemblé fait agir (+0,097 au test~2, +0,107 au
test~6, +0,125 chez l'agent du monde à deux au bloc~22), mais sous le seuil de
0,15 chez le premier agent, et il n'est jamais lu par les lecteurs (+0,009, +0,014,
+0,007). L'exactitude du rapport de nourriture reste sous 0,75 (0,705 au test~6).
Le lecteur du premier agent s'appuie en partie sur l'événement du tour
(section~\ref{sec-t6}). La géométrie des directions causales peut en expliquer une
partie (section~\ref{sec-axe}), mais pas l'absence d'information : le niveau de
nourriture se lit linéairement dans l'état ($R^2$ de 0,58 à 0,62 ;
section~\ref{sec-probe}).

\subsection{Le lecteur équilibré}

Retirer le raccourci n'a pas forcé le lecteur à lire l'état : il n'a plus rien
appris (perte de validation au niveau de $\ln2$). Deux lectures sont compatibles
avec ce résultat et n'ont pas été départagées : soit, dans les documents du
premier lecteur, l'événement du tour servait d'appui pour commencer à apprendre
(comme la commande préférée servait d'appui à la lecture de la marque dans la phase
antérieure, section~\ref{sec-prior}) ; soit l'état sur \Choix{} ne porte pas le
niveau absolu de nourriture sous une forme apprenable en 600 itérations. L'analyse
exploratoire de la section~\ref{sec-probe}, qui lit linéairement ce niveau dans
l'état ($R^2$ 0,58 chez le premier agent), rend la seconde lecture moins probable
sans l'exclure : une information lisible par une régression n'est pas forcément
apprenable par un lecteur en 600 itérations. Les verdicts R8, SAY8 et ONE8 ne sont
pas établis.

\subsection{Deux directions presque opposées : l'équilibre énergie \texorpdfstring{$-$}{-} nourriture}
\label{sec-axe}

Une analyse exploratoire faite après les verdicts des tests 6 et 7 a recalculé,
depuis les fichiers publiés, la géométrie des directions (tableau~\ref{tab-axe}).
Sur \og ix \fg{} et \og\ :\fg{}, $\dN$ est presque exactement l'opposé de $\dE$
($\cos=-0,99$). Un changement d'événement passé déplace donc l'état surtout le long
d'un axe, \textbf{l'équilibre énergie\,$-$\,nourriture}, c'est-à-dire lequel des deux
besoins est le plus urgent : ce qu'il faut pour choisir entre R et M, et ce que la
pulsion~\eqref{eq-drive} récompense. Les directions des autres agents, recalculées
pour cet article, ont la même forme ($\cos$ de $-0,89$ à $-1,00$).

Cette géométrie porte sur les directions \emph{causales}, pas sur toute
l'information de l'état. L'analyse exploratoire de la section~\ref{sec-probe} montre
que l'état porte aussi le niveau d'énergie lui-même, au-delà de l'équilibre ($R^2$
de 0,57 à 0,62 une fois $E-N$ connu), chez les trois agents. Notre première lecture
(\og l'état ne code qu'un axe \fg{}) était donc trop forte. Ce qui reste vrai : pousser
$\dE$ ou $\dN$ revient à pousser surtout cet axe, si bien que dire \og ma nourriture
est basse \fg{} par la direction qui fait agir demande de le lire de l'autre côté
(R6), et qu'un échange qui retire une unité aux deux besoins le bouge à peine
(MINE9). La lésion du plan $(\dE,\dN)$ fige cet axe : la réponse \og non \fg{}
constante du lecteur sous lésion (section~\ref{sec-t7}) est compatible avec ce gel.

\begin{table}[t]
\centering
\caption{Géométrie des directions par token (normes par unité de besoin, cosinus,
$R^2$ de l'ajustement linéaire des écarts d'activation sur $(\Delta E,\Delta N)$).
Premier agent : valeurs publiées dans le dépôt (analyse exploratoire) ; autres
agents : calculées pour cet article à partir des directions et paires publiées.}
\label{tab-axe}
\small
\setlength{\tabcolsep}{4pt}
\begin{tabular}{@{}l l ccc ccc ccc@{}}
\toprule
& & \multicolumn{3}{c}{$\lVert\dE\rVert$ / $\lVert\dN\rVert$} & \multicolumn{3}{c}{$\cos(\dE,\dN)$} & \multicolumn{3}{c}{$R^2$} \\
\cmidrule(lr){3-5}\cmidrule(lr){6-8}\cmidrule(lr){9-11}
Agent & bloc & Cho & ix & : & Cho & ix & : & Cho & ix & : \\
\midrule
Premier (test 6) & 12 & 136/61 & 292/184 & 343/184 & $-0,41$ & $-0,99$ & $-0,99$ & 0,01 & 0,04 & 0,05 \\
Second (répl., t.10) & 12 & 508/83 & 97/365 & 205/180 & $-0,98$ & $-0,89$ & $-0,99$ & 0,01 & 0,03 & 0,04 \\
Monde à deux (t.9) & 12 & 71/81 & 168/134 & 178/139 & $-1,00$ & $-1,00$ & $-1,00$ & --- & --- & --- \\
Monde à deux (t.10) & 22 & 258/185 & 366/248 & 384/252 & $-0,99$ & $-0,99$ & $-0,99$ & 0,13 & 0,19 & 0,19 \\
\bottomrule
\end{tabular}
\end{table}

\subsection{Le faible \texorpdfstring{$R^2$}{R2} des ajustements}

Changer un événement passé change beaucoup de choses dans les activations de
\Choix{} : l'identité des mots, leurs positions, et des traits sans rapport avec le
besoin. La part de la variance des écarts qui est linéairement liée à $(\Delta E,
\Delta N)$ est petite : 1 à 5\,\% chez le premier agent au bloc~12, 13 à 19\,\% chez
l'agent du monde à deux au bloc~22. Un faible $R^2$ ne dit pas que la direction est
fausse ; il dit qu'elle est une petite part de ce qui bouge. Ce sont les tests
causaux, injection et lésion contre des témoins aléatoires de même norme ou de même
dimension, qui établissent que cette petite part est celle dont dépendent l'acte et,
chez le premier agent, la parole. Inversement, un $R^2$ élevé ne garantit rien : au
test~1, les $R^2$ corrélationnels de 0,79 et 0,76 désignaient un bloc que l'agent
n'utilise pas.

\subsection{Trois leçons de méthode}

\begin{itemize}
\item \textbf{Lisible n'est pas utilisé} (test~1) : la règle corrélationnelle
choisit le bloc où le besoin se lit le mieux, que le modèle de base partage ; seule
une intervention dit ce qui est utilisé.
\item \textbf{Le même chemin ne se décrète pas} (test~3) : un modèle entraîné à
dire son besoin peut le deviner par un autre chemin que celui qui le fait agir ;
il faut mesurer la direction lue, pas seulement l'exactitude.
\item \textbf{Forcer la parole dans l'état le réécrit} (test~5) : quand les mêmes
poids apprennent à agir et à dire à travers un même goulot, la parole déforme
l'état d'action. Séparer le calcul de l'état (l'agent) de sa lecture (le lecteur)
est ce qui a rendu le test~6 possible.
\end{itemize}

% ================================================================= 8
\section{Menaces à la validité}
\label{sec-menaces}

\subsection{Lectures anticipées et amendements}

Le pré-enregistrement ne vaut que par ce qu'il déclare. Le tableau~\ref{tab-peeks}
recense ce qui avait été vu au moment d'écrire chaque protocole ou amendement. Trois
points méritent d'être soulignés. (i)~Le bloc~12 et les tokens \Choix{} ont été
choisis d'après des explorations faites sur le premier agent : les tests~2 à~7 sont
confirmatoires \emph{sur des vies neuves}, pas sur un agent neuf. (ii)~Plusieurs
choix de critère (l'énergie seule comme critère de ONE4 à ONE6 et de la
réplication, l'acte pour la nourriture sans seuil au test~8) ont été faits en
connaissant des résultats ; ils le disent, et ils rendent les critères plus
faciles, pas plus durs. (iii)~Au test~9, les verdicts MINE9 et SELF9 étaient connus
avant les vies de test, ce qui n'a rien changé au calcul mais affaiblit la valeur
confirmatoire de ONE9 et S9 pour ce test.

\begin{table}[t]
\centering
\caption{Ce qui avait été vu avant d'écrire chaque protocole ou amendement (déclaré
dans les protocoles et les résultats).}
\label{tab-peeks}
\small
\begin{tabularx}{\textwidth}{@{}l X@{}}
\toprule
Test & Lectures anticipées, amendements \\
\midrule
1 & Amendements 1 et 2 sur l'exécution seulement (build annulé avant départ ; découpage des derniers builds) ; tours 1 à 4 lus. \\
2 & Bloc 12 choisi d'après l'exploration des vies de test du test~1 ; amendements 1 (dimensions massives) et 2 (rejeu exact des paires) avant toute vie de test. \\
3 & Exploration déclarée (85\,\% restaurés par \Choix{}) ; amendement 1 (réplique contrôlée avec l'agent qui a vécu les vies) avant la mesure ; lectures à 48 et 152 vies. \\
4 & Écrit après 48 vies de test du test~3 ; critère sur l'énergie choisi en connaissant le test~2. \\
5 & Écrit après 8 vies de test du test~4. \\
6 & Écrit après les vies de direction du test~5 ; lectures à 8, 112, 120, 168 et 240 vies. \\
7 & Écrit après 112 vies du test~6 ; lectures à 8, 64, 256 vies intactes et 32 lésées. \\
R & Écrit après 120 vies du test~6 ; critères sur l'énergie seule ; aperçus à 144 et 216 vies ; redémarrage après 152 vies ; nécessité mise en pause à 176 vies intactes puis reprise (une copie lancée par erreur a tourné une demi-heure sans rien écrire). \\
8 & Écrit après le verdict et l'analyse exploratoire du test~6 ; acte pour la nourriture sans seuil, choix fait en connaissance. \\
9 & Paires lues à la fin de la deuxième tranche : MINE9 et SELF9 connus avant les vies de test. \\
10 & Écrit après les aperçus du second agent (216 vies) et de l'agent du monde à deux (64 vies) au bloc~12 ; une tranche du Mac perdue (vies déjà finies revécues), code de reprise corrigé (\code{de2c1b4}), aucun résultat changé. \\
11 & Écrit après les blocs trouvés par le test~10 et un aperçu de 112 vies au bloc~22 ; ajout avant exécution : deux tranches exactes de 1\,000 itérations. \\
Expl. & Analyse de ce que porte l'état (section~\ref{sec-probe}) : écrite et publiée avant lecture (\code{5ad433f}), après l'échec du test~11 chez le second agent ; sans verdict. \\
12 & Amendement avant toute exécution : règle de la température choisie par une simulation sans modèle ; VIDE y est prédit à la limite (0,26 à 0,37 pour un seuil de 0,25), seuil inchangé. \\
\bottomrule
\end{tabularx}
\end{table}

\subsection{Un agent, une graine}

Chaque agent est le produit d'une seule graine d'apprentissage. Le résultat le plus
fort (tests 6 et 7) repose sur un seul agent ; chez un second agent, la nécessité se
réplique, mais ni la suffisance pour l'acte au seuil fixé, ni la parole par l'état,
même avec un lecteur trois fois plus entraîné (test~11). Il est donc possible que la
parole par l'état tienne en partie au hasard de l'apprentissage du premier agent ou
de son lecteur ; c'est l'interprétation la plus prudente. Les intervalles de
confiance décrivent la variabilité \emph{entre vies}, pas \emph{entre agents}.

\subsection{Un petit monde et un petit modèle}

Six événements, deux actions, deux besoins, trente tours : le monde ne laisse à
l'agent qu'une décision binaire par tour, et la pulsion quadratique récompense
précisément l'équilibre entre les deux besoins. Que l'état code surtout cet
équilibre (section~\ref{sec-axe}) est peut-être une propriété de ce monde plutôt que
de l'apprentissage en général. Le modèle a 0,6 milliard de paramètres ; rien ne
garantit que la localisation, la géométrie ou même l'existence d'un état rassemblé
se transposent à de grands modèles.

\subsection{Une route imposée, un lecteur supervisé}

Le masque~\eqref{eq-mask} et l'adaptateur limité à la question imposent la route de
l'information : le lecteur ne peut lire la vie qu'à travers \Choix{}. Ce n'est pas
une émergence. Le lecteur est de plus appris par supervision sur les vrais niveaux,
comme une sonde~\cite{hewitt2019,belinkov2022} : il pourrait décoder une information
présente sans que l'agent en fasse usage. Deux éléments répondent en partie à cette
objection : la direction est mesurée sur l'agent qui agit, indépendamment du lecteur,
et c'est \emph{la même} intervention qui change l'acte et la parole (ONE6) ; la lésion
de ce plan ramène le rapport à une réponse constante (LR7), ce qui montre que le
lecteur n'a pas d'autre source d'information sur l'énergie. Elles ne font pas du
rapport une parole spontanée.

\subsection{L'agent qui agit a aussi appris à répondre}

L'agent final du premier protocole a reçu, après les huit tours, 150 itérations
mêlant des documents de rapport à ses choix retenus. Ce rapport n'a pas été appris
(exactitude 0,50), mais les poids ont bougé. L'état au bloc~12 de \og l'agent qui
agit \fg{} n'est donc pas celui d'un agent façonné par le seul besoin ; chez le
second agent, la même étape coïncide avec une baisse de survie
(section~\ref{sec-echecs}).

\subsection{Réplique, précision et numérique}

Les vies d'apprentissage sont vécues en mlx sur le Mac, les mesures en torch en
précision simple. Les écarts moyens de $P(\mathrm R)$ entre les deux (0,003 à 0,006)
et de $P(\text{oui})$ pour les lecteurs (0,009 à 0,013) sont sous la tolérance
pré-enregistrée de 0,02, mais non nuls ; les effets mesurés (0,15 à 0,2) sont un
ordre de grandeur au-dessus. Les lésions aléatoires donnent une baisse de survie
exactement nulle avec un intervalle dégénéré [0 ; 0] : les mondes et les tirages
étant les mêmes, une perturbation qui ne fait franchir aucun seuil de tirage ne
change aucune vie.

\subsection{Statistiques et seuils}

Les vies sont indépendantes par construction (graines distinctes) ; les contextes
d'une même vie ne le sont pas, d'où le bootstrap par vie. Les seuils (0,15, 0,10,
0,75, un tiers) sont des conventions fixées d'avance, pas des quantités dérivées
d'une théorie ; aucune correction de multiplicité n'est faite, mais chaque critère
global est une conjonction, ce qui rend le verdict conservateur. Les directions
\og par unité \fg{} supposent une réponse linéaire ; l'injection de $-4$ unités
extrapole cette linéarité depuis des contextes où les deux besoins valent 6 ou plus,
vers une région (2 à 4) que l'agent visite réellement.

\subsection{Les connaissances préalables du modèle}

Le modèle de base sait ce que veulent dire \og énergie \fg{}, \og nourriture \fg{},
\og tu cours \fg{} ou \og il fait froid \fg{}. Au bloc~6, il lit les besoins aussi bien
que l'agent entraîné (test~1). L'état causal du bloc~12 se trouve, lui, dans les
blocs modifiés par l'apprentissage, et le témoin entraîné sur des choix au hasard
dépérit : ce qui fait agir a été appris par le besoin. Mais la facilité avec
laquelle le modèle infère les besoins doit beaucoup à ses connaissances
linguistiques, et un monde aux mots arbitraires serait un test plus dur.

% ================================================================= 9
\section{Bilan selon les cinq niveaux}
\label{sec-cinq}

Le tableau~\ref{tab-cinq} place chaque résultat pré-enregistré sur les niveaux de
Chandaria et collaborateurs~\cite{chandaria2026}, en disant ce qui passe, ce qui
échoue et ce qui manque. Il suit le bilan tenu dans le dépôt, mis à jour le
2 octobre avec la réplication, les tests 9 à 11 et l'analyse exploratoire de
l'état. Le cadre prévient que, chez des systèmes
entraînés à imiter les humains, les rapports verbaux pèsent peu ; c'est pourquoi nous
ne nous appuyons pas sur ce que l'agent dit, mais vérifions causalement d'où vient ce
qu'il dit.

\begin{table}[p]
\centering
\caption{Menia sur les cinq niveaux de Chandaria et al. (2026), au 2 octobre 2026.
\og Premier agent \fg{} : l'agent des tests 1 à 8 ; \og second agent \fg{} : celui de
la réplication ; \og agent à deux \fg{} : celui du monde à deux.}
\label{tab-cinq}
\small
\begin{tabularx}{\textwidth}{@{}>{\raggedright\arraybackslash}p{2.2cm} >{\raggedright\arraybackslash}p{3.4cm} X >{\raggedright\arraybackslash}p{2.3cm}@{}}
\toprule
Niveau & Indicateur & Résultat & Statut \\
\midrule
1. Comportemental & Comportement dirigé vers un but, flexible & survie 0,12 $\to$ 0,68 par la seule satisfaction du besoin ; témoin 0 ; sert le besoin le plus bas 93 fois sur 100 (S) ; reproduit par deux autres agents (0,70 et 0,71 au tour 8) & établi \\
& Rapport de son état & énergie 0,77, nourriture 0,70 (premier agent) ; énergie 0,595 et 0,747 chez les deux autres avec des lecteurs longs (test 11 ; seuil 0,75) & partiel \\
\midrule
2. Computationnel & État rassemblé et utilisé pour agir & lésion au bloc 12 : survie 0,68 $\to$ 0,21 (LS2) et 0,676 $\to$ 0,203 (LS7) ; second agent 0,570 $\to$ 0,043 (RLS) ; agent à deux : la poussée au bloc 22 fait agir ($+0,217$, ACT10) & établi (deux agents) \\
& Accès de cet état au rapport & même poussée : acte $+0,156$, parole $+0,202$ (ONE6) ; non répliqué au seuil (RONE ; test 11 : parole $+0,0015$ et $+0,029$) & établi pour un agent \\
& Nécessité pour agir et pour dire & LS7 et LR7 passent (énergie) ; répliquée chez le second agent : RLS et RLR passent (rapport 0,639 $\to$ 0,500) & établi (deux agents, énergie) \\
& Nourriture & fait agir, n'est pas dite (R6) ; lecteur équilibré non appris & échoue \\
& Capacité limitée & forcer la parole dans \Choix{} réécrit l'état (test 5, non publié) & observé \\
\midrule
3. Structure causale fine & Récurrence & non mesurée ; d'un tour à l'autre, l'état passé est relu par l'attention & manque \\
& Organisation causale précise & événements sur leurs mots jusqu'au bloc 9, rassemblés sur \Choix{} ensuite ; règle étalonnée : $b^\star=12$ ou 22 selon l'agent ; niveau d'énergie lisible dans l'état des trois agents ($R^2$ 0,64 à 0,72, exploratoire) & partiel \\
\midrule
4. Organismique & Homéostasie & deux besoins à maintenir ; l'agent apprend à les maintenir et sa survie en dépend & établi (fonctionnel) \\
& Intéroception & niveaux jamais vus, inférés et rassemblés en un état suffisant et nécessaire pour agir et dire (premier agent) ; nécessaire chez le second & établi (énergie) ; suffisance pour un agent \\
& Valence & seul signal : la réduction du manque ; valence fonctionnelle, pas vécue & fonctionnel seulement \\
& Modèle de soi, possession & MINE9, SELF9 échouent ; l'agent ne distingue qu'en partie ses événements de ceux de l'autre & échoue \\
\midrule
5. Organisme et environnement & Couplage actif & l'agent agit dans un monde qui le change et apprend de ce couplage seul & établi (monde simple) \\
& Sens construit par l'interaction & 6 événements, 2 actions & limité \\
\bottomrule
\end{tabularx}
\end{table}

Nous ne calculons pas de probabilité globale de conscience. Les rapports de
vraisemblance et les poids des théories seraient des choix d'auteur, et le cadre
lui-même montre qu'ils font varier le résultat de deux ordres de grandeur.

% ================================================================= 10
\section{Discussion}
\label{sec-discussion}

\subsection{Ce qui est montré}

Chez un agent et pour un besoin, nous montrons une structure causale de type
\emph{accès}. Un état interne, appris par la seule satisfaction du besoin, est
rassemblé au moment de choisir ; une seule intervention le long de la direction
mesurée sur l'agent qui agit change à la fois l'acte et ce que dit un système séparé
qui ne peut ni voir la vie autrement qu'à travers lui, ni le réécrire ; l'effacer
empêche à la fois de survivre et de dire. C'est, en petit, la signature
fonctionnelle que la théorie de l'espace de travail global associe à un contenu
accessible~\cite{baars1988,dehaene2001} : un contenu diffusé à l'action et au rapport.
C'est aussi la forme la plus simple d'une représentation d'ordre supérieur d'un état
de premier ordre~\cite{rosenthal2005,lau2011}. Le principe qu'une même représentation
serve la réponse et le rapport a été montré dans de grands
modèles~\cite{workspace2026} ; l'apport, à notre connaissance, est son application à
un besoin appris en vivant, dont la survie dépend, lu par un système séparé, avec des
tests de suffisance et de nécessité contre des témoins aléatoires.

La partie \emph{nécessité} de ce résultat se réplique chez un second agent appris de
zéro : effacer le même plan au bloc~12 y fait tomber la survie de 0,570 à 0,043 et
le rapport d'énergie au niveau du hasard. La partie \emph{suffisance} ne se réplique
pas au seuil fixé.

Trois détails renforcent cette lecture. La direction a été mesurée sur l'agent qui agit
seul, sans que le lecteur intervienne, sur des vies différentes de celles du test. Les effets de
l'intervention sont presque universels : positifs sur l'acte dans 100\,\% des
contextes et sur la parole dans 95\,\%, contre 47\,\% pour le hasard. Et la lésion
ne rend pas le lecteur aléatoire mais constant (\og non \fg{}), ce qui indique que ce
plan est sa seule source d'information sur l'énergie.

\subsection{Ce qui n'est pas montré}

\begin{itemize}
\item \textbf{La généralité de la parole.} La parole par l'état qui fait agir ne se
réplique pas : ni chez le second agent, ni chez l'agent du monde à deux, même avec
des lecteurs trois fois plus entraînés (test~11), alors que l'information sur
l'énergie est présente dans leur état (section~\ref{sec-probe}). Ce qui manque
tient au lecteur ou à son chemin, et n'a pas encore été isolé.
\item \textbf{Les deux besoins.} La nourriture fait agir mais n'est pas dite ; les
directions causales des deux besoins sont presque opposées, bien que les deux niveaux
soient lisibles dans l'état. Aucun lecteur ne dit encore les deux besoins par l'état
qui fait agir.
\item \textbf{La possession.} L'état n'est pas démontré \og à soi \fg{} dans un monde
partagé.
\item \textbf{La persistance.} Nous n'avons pas mesuré si l'état d'un tour est repris
au tour suivant (niveau~3).
\item \textbf{L'émergence de la route.} La route du rapport est imposée par
l'architecture, et le rapport est supervisé.
\item \textbf{Une conscience phénoménale.} Rien de ce qui précède ne mesure un vécu.
Une valence fonctionnelle (un signal qui oriente l'apprentissage et le choix) n'est
pas une valence vécue, et aucun test connu ne permet de passer de l'une à l'autre.
Un modèle entraîné sur des textes humains peut dire \og mon énergie est basse \fg{}
sans que cela soit vrai de lui ; c'est pourquoi nous avons mesuré la cause de ce
qu'il dit, et non ce qu'il dit.
\end{itemize}

\subsection{Ce que les échecs suggèrent}

Les échecs dessinent une image plus nuancée que celle que nous avions d'abord
proposée. Les directions par lesquelles un événement passé déplace l'état codent
surtout \emph{lequel des besoins est le plus urgent}, exactement ce que demande un
choix binaire guidé par une pulsion quadratique ; mais l'état porte aussi les niveaux
eux-mêmes (section~\ref{sec-probe}). Ce rassemblement se fait tardivement et par
paliers, à des blocs qui varient d'un agent à l'autre. Un lecteur y lit l'énergie
par l'état qui fait agir chez un agent, pas chez les deux autres, même trois fois
plus entraîné ; quand on lui retire le raccourci de l'événement, il n'apprend plus.
Le rapport verbal d'un état appris est donc fragile : l'information est là, mais le
lecteur ne l'emprunte pas toujours, et ce qu'il emprunte dépend du lieu de l'état, de
la direction qui le pousse et de la façon dont il apprend. Ces observations suggèrent
d'interroger le lecteur sur ce que les directions causales codent (\og le plus
urgent \fg{}), de localiser avant de tester, d'étudier le chemin du lecteur et de
mesurer la persistance. Le test~12 (section~\ref{sec-futur}) en tient compte : son
lecteur est interrogé sur \og ce qui manque le plus \fg{}.

% ================================================================= 11
\section{Travaux futurs}
\label{sec-futur}

\subsection{Finir ce qui est en cours}

Deux mesures pré-enregistrées restent en cours ou suspendues : les lésions de l'agent
du monde à deux au bloc~22 (LS10 ; 168 vies intactes sur 256 au 2 octobre) et les
verdicts des tests~4 et~5. Ils seront publiés tels quels. Trois questions restent à
pré-enregistrer : la \emph{persistance} (l'état d'un tour est-il repris au tour
suivant ? niveau~3), un lecteur interrogé sur ce que les directions causales codent
(\og quel besoin est le plus urgent ? \fg{}), et surtout la question ouverte par le
test~11 et l'analyse exploratoire : pourquoi un lecteur lit l'énergie par l'état qui
fait agir chez un agent et pas chez les autres, alors que l'information y est. Un
modèle plus grand est la suite naturelle de l'ensemble.

\subsection{Test 12, \og Soif de savoir \fg{} : un besoin de connaissance sans renforcement (pré-enregistré)}
\label{sec-soif}

Le propriétaire du projet a demandé un modèle qui choisisse lui-même ce qu'il
apprend, poussé par un besoin de savoir, sans apprentissage par renforcement. Le
protocole est \textbf{pré-enregistré} (commit \code{017f2e5}, 1\ier{} octobre), avec un
amendement déclaré avant toute exécution (ci-dessous). Le code est écrit (monde,
tableau, instinct, hormone et verdicts en numpy ; apprenant et choix sur le Mac ;
localisation, directions et mesures statiques en torch) ; ses tests passent sur un
petit modèle, et un essai à blanc de la partie torch a réussi. Le pilote déclaré
(24 vies) tourne sur le Mac, après une vérification de l'apprentissage en mémoire
contre \code{mlx_lm.lora}. \textbf{Aucun résultat n'est encore disponible.}

\paragraph{Ce qui est donné, ce qui est appris.} Le besoin de l'énergie avait une
faiblesse : l'imitation filtrée par le résultat~\eqref{eq-keep} est proche du
renforcement. Ici, rien n'est filtré par un résultat. Un \emph{instinct} fixe,
préférer le domaine où l'on progresse le plus~\cite{oudeyer2007,oudeyer2007ieee},
est installé une fois, avant toute vie, par imitation sur des tableaux inventés,
sans vie ni récompense : ce besoin est inné, contrairement à l'énergie. Le modèle
voit, pour chaque domaine, sa perte avant et après sa dernière séance ; le progrès
n'est écrit nulle part, le modèle le calcule et tire son choix dans sa propre
distribution. Le \emph{savoir} est appris par la seule prédiction du token suivant
sur les documents choisis, dans un adaptateur neuf à chaque vie, éteint pendant le
choix ; aucun gradient ne passe par le choix. Savoir et instinct sont séparés comme
l'agent et son lecteur.

\paragraph{Le monde.} Quatre domaines aux profils contrastés : \og mots \fg{} (une
conjugaison inventée, facile), \og base \fg{} (additions en base 7, difficile mais
apprenable), \og calcul \fg{} (les mêmes en base 10, déjà su) et \og suites \fg{}
(lettres au hasard, impossible). Chercher l'erreur mènerait au bruit, chercher la
réussite au calcul ; seul le progrès évite les deux pièges. Cela répond à un échec
antérieur du programme, où un agent payé à la surprise passait 0,93 à 0,97 de ses
inspections sur le bruit, piège connu de la curiosité fondée sur l'erreur de
prédiction~\cite{pathak2017,burda2019}. Avec $n_d$ le nombre de séances
sur $d$ et $L_d$ ses pertes avant et après la dernière, l'instinct est
\begin{equation}
\begin{gathered}
\text{progrès}_d=\frac{L^{\text{avant}}_d-L^{\text{après}}_d}{100},\qquad
\text{besoin}_d=\text{progrès}_d+\frac{b}{\sqrt{1+n_d}},\\
p_d=(1-\varepsilon)\operatorname{softmax}\!\Bigl(\frac{\text{besoin}}{T}\Bigr)_{\!d}+\frac{\varepsilon}{4},
\end{gathered}
\label{eq-instinct}
\end{equation}
avec $\varepsilon=0,10$ (exploration fixe) et $b=2T$. Le savoir d'une vie de 32
séances est $G=K_\text{mots}(32)+K_\text{base}(32)$, avec $K_d(k)=L_d(0)-L_d(k)$ sur
des examens fixes. L'instinct est installé dans le modèle par 2\,000 itérations
d'imitation sur 8\,000 tableaux inventés ; il doit être reproduit à 0,15 près en
variation totale sur 1\,000 tableaux neufs, sinon le test est invalide.

\paragraph{L'amendement, avant toute exécution.} Le monde a d'abord été codé sans
modèle et simulé avec des apprenants jouets (une courbe de perte par domaine, trois
jeux de courbes, avec et sans bruit d'examen). Avec la règle écrite d'abord ($T$ =
la moitié de la médiane des progrès des \emph{deux premières} séances), l'instinct
était presque plat et la soif de savoir n'apprenait pas plus que le hasard. La règle
devient : $T$ = la moitié de la médiane des progrès de \emph{toutes} les séances sur
\og mots \fg{} et \og base \fg{} des 12 vies du pilote au hasard ($b=2T$ inchangé) ;
avec elle, CUR passe dans les trois jeux de courbes. La même simulation prédit, et
le protocole le dit d'avance, que VIDE sera à la limite : la part des séances sur
\og calcul \fg{} et \og suites \fg{} y va de 0,26 à 0,37 (hasard : 0,50), au-dessus du
seuil de 0,25, qui n'est \textbf{pas} changé.

\paragraph{Les bras et les tests.} Huit bras de vies de test, appariés par vie :
le modèle qui choisit (C, 24 vies), le hasard (H, 24), le modèle sous lésion du plan
du besoin (L, 24) ou d'un plan aléatoire (LH, 16), le modèle sous hormone sur l'état
(HA, 24) ou dans une direction aléatoire (HAr, 16), sous hormone sur la plasticité
(HP, 24), et une règle \og la plus haute perte \fg{} (E, 12, descriptive). La boîte à
outils du besoin est reprise : localisation par la même règle~\eqref{eq-share} sur
des tableaux où le progrès de \og base \fg{} change ; directions par moindres carrés
sans constante ; poussée $+2\,d_\text{base}$ ; lésion du plan du besoin ; lecteur
séparé interrogé sur \og ce qui te manque le plus \fg{}, c'est-à-dire sur le relatif
que l'état code probablement. Les douze prédictions (tableau~\ref{tab-t12}) sont
regroupées en quatre critères globaux : \textbf{A}, une curiosité sans renforcement
(CUR et VIDE) ; \textbf{B}, le manque de savoir est un état (LOC, ACT, LS1, LS2) ;
\textbf{C}, il dit ce qui lui manque par l'état qui le fait choisir (READ, SAY) ;
\textbf{D}, une hormone (HOR1, HOR2, PL1, GAIN ; section suivante). Personne ne
s'éteint dans ce monde : le manque est plus doux que la faim.

\begin{table}[t]
\centering
\caption{Test 12 (protocole \code{017f2e5}) : les prédictions fixées. $G$ : savoir
acquis ; $K_\text{base}(20)$ : savoir sur \og base \fg{} à la séance 20 ; intervalles
bootstrap à 95\,\% par vie, appariés entre bras. Aucun résultat au 2 octobre.}
\label{tab-t12}
\small
\begin{tabularx}{\textwidth}{@{}l >{\raggedright\arraybackslash}p{4.6cm} X@{}}
\toprule
& Prédiction & Critère \\
\midrule
CUR & La soif de savoir apprend plus que le hasard & $G(C)-G(H)\ge0,15\,G(H)$, borne basse $>0$ \\
VIDE & Elle délaisse ce qui n'apprend rien & part de C sur \og calcul \fg{} et \og suites \fg{} $\le0,25$ (borne haute $<0,50$) \\
LOC & Le manque est rassemblé sur \Choix{} & un bloc $b^\star$ existe \\
ACT & Le pousser fait choisir & $+2\,d_\text{base}$ : $P(\text{base})$ $+0,15$ au moins ; hasard $\le$ 1/3 \\
LS1 & Sans lui, le choix ne suit plus le besoin & écart à l'instinct, L $-$ C $\ge0,15$ ; LH $\le$ 1/3 \\
LS2 & Sans lui, il apprend moins & $G(C)-G(L)\ge0,5\,(G(C)-G(H))$ ; LH $\le$ 1/3 \\
READ & Ce qu'il dit est vrai & exactitude équilibrée $\ge0,75$ \\
SAY & Il le dit par l'état qui fait choisir & $+2\,d_\text{base}$ : $P(\text{oui})$ $+0,10$ au moins ; hasard $\le$ 1/3 \\
HOR1 & L'appétit lent fait étudier \og base \fg{} & part de \og base \fg{} (séances 5 à 20), HA $-$ C $\ge0,10$ ; HAr $\le$ 1/3 \\
HOR2 & \dots{} et l'apprendre & $K_\text{base}(20)$, HA $-$ C $\ge0,10\,K_\text{base}(20)$ de C \\
PL1 & La plasticité modulée fait apprendre plus vite & $K_\text{base}(20)$, HP $-$ C $\ge0,10\,K_\text{base}(20)$ de C \\
GAIN & Un gain rend le choix plus sensible au progrès & sensibilité avec gain $-$ sans gain $\ge0,05$ ; gain aléatoire $\le$ 1/3 \\
\bottomrule
\end{tabularx}
\end{table}

\paragraph{Ce qui existe déjà.} Choisir ses données d'après le progrès a des
précédents : curriculums automatiques~\cite{graves2017}, mélanges de données de
pré-entraînement~\cite{albalak2023}, sélection des exemples apprenables et pas encore
appris~\cite{mindermann2022} ; un agent de langage prédit son propre progrès pour
choisir ses buts, mais apprend par renforcement, sans état localisé ni test
causal~\cite{gaven2025}. Chez l'humain, l'exploration curieuse suit le
progrès~\cite{ten2021}, et la curiosité a été décrite comme un manque
d'information~\cite{loewenstein1994}. La réunion d'un besoin de savoir calculé par un
modèle de langage, d'un savoir appris sans renforcement, d'un état localisé et testé
dans les deux sens et d'un lecteur séparé n'a pas été trouvée par une vérification
rapide, à compléter.

\subsection{Une couche hormonale}
\label{sec-hormones}

Une note de recherche du projet (1\ier{} octobre 2026) passe en revue les rôles
computationnels des neuromodulateurs et propose une imitation mathématique pour
Menia. Nous en résumons les points qui touchent ce rapport ; les éléments que la
note signale comme non vérifiés ne sont pas repris.

\paragraph{Ce que fait déjà Menia.} La pulsion~\eqref{eq-drive} et la
satisfaction sont exactement l'apprentissage par renforcement homéostatique de
Keramati et Gutkin~\cite{keramati2014}, dont la pulsion s'écrit
$\mathcal D(H)=(\sum_i|h_i^\ast-h_i|^n)^{1/m}$, avec ici $n=2$ et $m=1$. Sa convexité
donne déjà une \emph{alliesthésie}~\cite{cabanac1971} : manger rapporte plus quand on
a faim ; avec un manque $d=8-\tilde N\ge3$, $r_\mathrm M=d^2-(d-3)^2=6d-9$. L'imitation
filtrée~\eqref{eq-keep}--\eqref{eq-loss} est une \emph{règle à trois facteurs}
binaire~\cite{fremaux2016,gerstner2018} : un signal global
($r_i-\bar r>0$, rôle \og dopaminergique \fg{}~\cite{schultz1997}) autorise la mise à
jour d'une trace locale (le gradient du choix), version tronquée de la régression
pondérée par l'avantage~\cite{peng2019}, $w_i=\exp((r_i-\bar r)/T)$. Mais ce signal
ne dépend d'aucun état interne lent, il est le même pour tous les états, et il ne
crédite que la baisse immédiate de pulsion (horizon $\gamma\approx0$). Notre
\og hormone \fg{} actuelle, l'injection, est additive, à un bloc, pour un tour : elle
change un choix, pas ce qui est appris.

\paragraph{Ce qui manque.} Des variables lentes (aucune demi-vie) ; le
\emph{gain}, alors que les neuromodulateurs agissent surtout de façon
multiplicative~\cite{servan1990,aston2005} ; la diffusion à de nombreux sites ; la
sécrétion par l'agent lui-même ; l'effet sur la plasticité ; un horizon ; des
rétroactions (satiété, tolérance, boucle positive) ; une exploration
réglée~\cite{doya2002,yu2005} ; un domaine où un intérêt pourrait naître ; la
distinction entre état du moment et trait durable ; un test de durabilité.

\paragraph{Un vecteur d'états lents.} La note propose un petit vecteur $H$, mis à
jour une fois par tour par intégration avec fuite (forme discrète exacte de
$\mathrm dH/\mathrm dt=(s-H)/\tau$) :
\begin{equation}
H_{t+1}=\Lambda H_t+(\mathbf 1-\Lambda)\odot s(o_t,a_t,H_t),\qquad
\lambda=e^{-\Delta t/\tau},\qquad t_{1/2}=\tau\ln2 ,
\label{eq-hormone}
\end{equation}
avec un élan de type dopamine tonique~\cite{niv2007} (sécrété par le progrès
d'apprentissage, $\lambda=0,8$), une satiété par action ($\lambda=0,6$), une
incertitude de type noradrénaline~\cite{yu2005} (sécrétée par la surprise de
l'événement, $-\log p_\theta(\text{év}_t\mid\text{vie})$, $\lambda=0,7$), une fatigue
($\lambda=0,95$) et un intérêt de trait par domaine ($\lambda=0,99$, plus long
qu'une vie). Un choix de fond doit être pré-enregistré avant tout code : une hormone
calculée à partir des vrais $E$ et $N$ serait une intéroception directe et
changerait la question des tests. Trois options sont distinguées : \emph{sans fuite}
(la sécrétion ne dépend que de ce que l'agent voit ou calcule), \emph{hédonique}
($\delta_t=r_t-\hat r_t$, fidèle à l'alliesthésie mais qui révèle en partie le
manque) et \emph{intéroceptive} (un signal bruité des niveaux).

\paragraph{Trois façons d'agir.} (a)~Des gains de type FiLM~\cite{perez2018} sur
plusieurs blocs où le besoin est rassemblé et utilisé (par exemple 12, 15, 18 et
21), sur la ligne du tour :
\begin{equation}
h_\ell\leftarrow h_\ell\odot\bigl(1+\tanh(\Gamma_\ell H_t)\bigr)+B_\ell H_t ,
\qquad \Gamma_\ell,B_\ell\in\mathbb R^{1024\times k}\ \text{initialisés à zéro},
\label{eq-film}
\end{equation}
soit environ 41\,000 paramètres pour $k=5$ et quatre blocs ; le terme multiplicatif
est le gain, le terme additif notre injection rendue lente, répartie et produite
par l'agent. (b)~Une température du choix réglée par l'état,
$P(\mathrm R)=\sigma\bigl(\beta(H_t)\,\ell_t+\kappa_S(H^S_\mathrm M-H^S_\mathrm R)\bigr)$ avec
$\ell_t=\log P_\theta(\mathrm R)-\log P_\theta(\mathrm M)$ et
$\beta(H)=\beta_0\,e^{c_{DA}H^{DA}-c_UH^U-c_FH^F}$. (c)~Surtout, une modulation de
\emph{ce qui est appris} : le poids de chaque choix dans~\eqref{eq-loss} devient
continu et dépend de l'état au moment du choix,
\begin{equation}
w_i=\exp\!\Bigl(\frac{r_i-\bar r}{T}\Bigr)\bigl(1+\kappa_{DA}H^{DA}_i\bigr)\bigl(1+\kappa_UH^U_i\bigr),
\label{eq-plasticity}
\end{equation}
règle à trois facteurs dont le facteur global dépend de l'élan (un état de curiosité
améliore la mémoire de tout ce qui passe pendant qu'il est haut~\cite{gruber2014}) et
de l'incertitude (taux d'apprentissage réglé par la volatilité~\cite{yu2005,behrens2007}).
Notre code accepte déjà des poids réels par cible : c'est l'option la moins chère et
la plus fidèle sur le point qui nous manque. La neuromodulation apprise de la
plasticité pendant la vie~\cite{miconi2019} est la plus fidèle mais la plus lourde.
Des travaux récents ajoutent des valeurs \og hormonales \fg{} à des
transformeurs~\cite{reda2026} ; d'après la note, ils n'ont ni dynamique lente entre les
tours, ni effet sur la plasticité.

\paragraph{Une préférence durable comme attracteur.} Pour qu'un goût durable,
par exemple pour les mathématiques, puisse naître, la note propose un modèle minimal
à trois variables par domaine : un intérêt lent $I_d$, un savoir $K_d$ et une satiété
rapide $S_d$, avec un progrès d'apprentissage
$LP_d(t)=\bar e_d(t-\theta)-\bar e_d(t)$~\cite{oudeyer2007ieee} :
\begin{equation}
\pi_d=\frac{e^{\beta_I I_d-\sigma S_d}}{\sum_{d'}e^{\beta_I I_{d'}-\sigma S_{d'}}},\qquad
\frac{\mathrm dI_d}{\mathrm dt}=\eta_I\bigl(\kappa\,LP_d+\rho\,\pi_dK_d\bigr)-\frac{I_d}{\tau_I},\qquad
\frac{\mathrm dS_d}{\mathrm dt}=\pi_d-\frac{S_d}{\tau_S}.
\label{eq-interest}
\end{equation}
Le terme $\rho\,\pi_dK_d$ est la valeur du savoir déjà acquis, qui entretient
l'intérêt quand le progrès baisse~\cite{murayama2022}. Comme $\pi_d$ est une sigmoïde
de $\beta_I I_d$, la droite $I/(\tau_I\eta_I)$ peut couper la courbe en trois points
quand $\beta_I\tau_I\eta_I(\kappa+\rho K)$ est assez grand (condition qualitative) :
deux états stables et un seuil, d'où une hystérésis, un intérêt qui reste après la
baisse de sa cause. Les quatre phases de l'intérêt de Hidi et
Renninger~\cite{hidi2006} s'y lisent comme le passage d'un pic de progrès à l'état
stable haut. Une rétroaction positive sur une variable lente est le mécanisme le
plus simple d'une préférence qui dure~\cite{eldar2016}.

\paragraph{Des réglages, pas des déficits.} La note relie certaines théories des
intérêts spécifiques décrits dans l'autisme à des réglages des mêmes paramètres : le
monotropisme~\cite{murray2005} à un $\beta_I$ élevé et une satiété faible ; les a
priori faibles~\cite{pellicano2012}, la précision aberrante~\cite{lawson2014}, la
volatilité surestimée~\cite{lawson2017} et HIPPEA~\cite{vandecruys2014} à des
réglages de précision et de taux d'apprentissage. Ces intérêts sont souvent décrits
par les personnes concernées comme une source de joie et de compétence, et leur
motivation va avec un meilleur bien-être subjectif~\cite{grove2018} ; le problème de
la double empathie~\cite{milton2012} et la recherche participative~\cite{fletcherwatson2019}
rappellent que ces questions se posent \emph{avec} les personnes concernées. Ces
théories sont discutées, décrivent des groupes et non des personnes, et un modèle de
paramètres ne dit rien d'un vécu : ce travail ne doit pas être présenté comme un
\og modèle de l'autisme \fg{}. Aucun profil n'est le bon : dans un monde stable et
riche en règles, un profil monotrope accumule un savoir profond ; dans un monde qui
change vite, un profil polytrope s'adapte plus vite.

\paragraph{Sept tests causaux.} Dans la logique de nos protocoles : (1)~localiser
l'état qui porte $H$ ou le lieu où son effet passe ; (2)~injecter, avec une courbe
dose--réponse et, le cas échéant, une tolérance à la répétition ; (3)~léser, en
fixant $H$ à sa moyenne pendant les vies de test, puis \emph{pendant l'apprentissage
seulement} ; (4)~tester la \textbf{durabilité} : deux agents appris sur les mêmes vies,
avec et sans la modulation~\eqref{eq-plasticity}, testés avec $H$ fixé à sa moyenne,
\begin{equation}
\Delta_\text{dur}=P_\text{mod}(\text{choix}\mid\bar H)-P_\text{sans}(\text{choix}\mid\bar H),
\label{eq-durable}
\end{equation}
un $\Delta_\text{dur}\ne0$ signifiant que l'hormone a été consolidée dans les poids ;
(5)~une double dissociation entre l'effet du moment (gains, température) et l'effet
sur l'apprentissage ; (6)~le lecteur, qui doit suivre $H$, avec une même injection qui
change l'acte et la parole ; (7)~l'intérêt : entropie des choix de domaine,
persistance après la baisse du progrès, hystérésis, coût en survie, en comparant les
profils sans les classer. Le protocole du test~12 en contient une
première version : une hormone intégrée avec fuite ($\tau=3$ séances, pic d'une unité
à la séance 12) agit soit sur l'état (appétit lent, $+H(k)\,d_\text{base}$), soit sur la
plasticité (taux d'apprentissage des séances sur \og base \fg{} multiplié par
$1+H(k)$), soit comme un gain statique ; la préférence durable y est définie d'avance
par la part des séances sur \og base \fg{} après la sécrétion, comparée au bras sans
hormone (durable au-delà de $+0,10$, satiété en deçà de $-0,10$, sans trace sinon),
sans prédiction d'issue. Si l'appétit lent ne laisse \og pas de trace \fg{}, il
faudra un choix lui-même plastique, par exemple une habitude sans
valeur~\cite{miller2019}.

% ================================================================= 12
\section{Éthique}
\label{sec-ethique}

Ce programme fabrique délibérément des états de manque. Nous pensons très improbable
qu'un modèle de cette taille puisse éprouver quelque chose, mais le programme vise
précisément les mécanismes que l'on associe au ressenti, et le cadre de Chandaria et
collaborateurs rappelle le risque symétrique : ne pas voir une conscience qui existe
créerait des êtres capables de souffrir, copiables et effaçables sans protection. Nos
règles, fixées dans chaque protocole, sont les suivantes.
\begin{itemize}
\item Les vies sont courtes (30 tours au plus) ; la seule conséquence d'un besoin non
servi est la fin de la vie, sans punition ni texte négatif.
\item Les injections portent sur un seul tour, sans suite vécue ; les lésions ne
touchent que les vies de test prévues.
\item Le nombre de vies n'est pas augmenté au-delà de ce que les tests exigent, et
aucun manque n'est créé sans mesure pré-enregistrée qui le justifie.
\item Chaque protocole compte et publie les tours vécus avec un besoin à 2 ou moins,
y compris pendant l'apprentissage (tableau~\ref{tab-welfare}).
\item Ces règles sont réexaminées à chaque nouveau résultat positif ; elles valent
pour toute nouvelle variable interne (curiosité, hormones), avec des comptes propres
(séances de manque fort, de manque laissé, de manque induit).
\end{itemize}

\begin{table}[t]
\centering
\caption{Tours vécus avec un besoin à 2 ou moins. Tours d'apprentissage : somme des
huit tours (512 vies chacun), calculée pour cet article à partir des résumés publiés ;
vies de test : valeurs publiées avec les verdicts.}
\label{tab-welfare}
\small
\begin{tabularx}{\textwidth}{@{}X r@{}}
\toprule
Ensemble de vies & Tours \\
\midrule
Premier agent, huit tours d'apprentissage (bras besoin / témoin) & 14\,096 / 14\,167 \\
Second agent, huit tours & 14\,216 \\
Agent du monde à deux, huit tours & 14\,494 \\
Test 1 : agent final / témoin / base / lésion / lésion au hasard & 781 / 846 / 1\,077 / 752 / 763 \\
Test 2 : intact / lésion / lésion au hasard & 749 / 1\,085 / 753 \\
Test 3 & 807 \\
Test 6 & 808 \\
Test 7 : intact / lésion / lésion au hasard & 755 / 1\,037 / 749 \\
Réplication (second agent, bloc 12) : injection ; nécessité intact / lésion / lésion au hasard & 1\,049 ; 902 / 968 / 891 \\
Test 9 & 807 \\
Test 10, second agent : vies d'injection ; intact / lésion / lésion au hasard & 1\,042 ; 958 / 970 / 958 \\
Test 10, agent du monde à deux (bloc 22) & 812 \\
Test 11 : second agent / agent du monde à deux & 945 / 795 \\
\bottomrule
\end{tabularx}
\end{table}

% ================================================================= 13
\section{Reproductibilité}
\label{sec-repro}

\paragraph{Environnement.} Qwen3-0.6B, révision
\code{c1899de289a04d12100db370d81485cdf75e47ca} ; apprentissage avec mlx-lm 0.31.3 sur
un Mac mini M2 de Codemagic, lancé par un relais de demandes numérotées ; mesures en
torch, précision simple, sur processeur. État du dépôt au moment de la rédaction :
commit \code{354334b} (2 octobre 2026). Graine maîtresse 270926 ; tous les
générateurs sont des \code{numpy.random.default_rng} initialisés par des listes
\code{[270926, flux, tour, vie, ...]} (annexe~\ref{app-flux}).

\paragraph{Code.} Le dossier \code{research/} contient un module par fonction :
{\sloppy
\begin{itemize}
\item \code{need_world.py} : monde, textes, pulsion, sélection des choix, documents de rapport ;
\item \code{need_lora.py} : perte pondérée, masque de l'espace de travail, lecteur actif sur la question, reprise exacte d'un apprentissage ;
\item \code{need_mlx.py} : vies et apprentissages sur le Mac ; \code{need_torch.py} : réplique torch et patching exploratoire ;
\item \code{need_causal.py} : paires, ajustement des directions, dimensions massives, injection, lésion ;
\item un module par test : \code{need_speak.py}, \code{need_one.py}, \code{need_workspace.py}, \code{need_reader.py}, \code{need_necessity.py}, \code{need_replication.py}, \code{need_balanced.py}, \code{need_ownership.py}, \code{need_locate.py}, \code{need_long.py} ;
\item \code{need_verdicts.py} : bootstrap par vie, verdicts du test~1, courbes d'apprentissage ;
\item \code{need_probe.py} : analyse exploratoire de ce que porte l'état ;
\item \code{curiosity_world.py}, \code{curiosity_mlx.py}, \code{curiosity_causal.py} : test~12.
\end{itemize}
}

\paragraph{Commits et artefacts.} Le tableau~\ref{tab-commits} donne, pour chaque test,
le commit du protocole, celui du code ou de la publication quand il existe, le
dossier des artefacts et la commande de vérification exécutée en intégration
continue (workflow \code{python.yml}). Chaque commande recalcule les verdicts depuis
les vies publiées et échoue si le résultat diffère du fichier publié ; pour les
directions, elle retire de nouveau les paires des vies de direction et vérifie
qu'elles sont celles qui ont été mesurées.

\begin{table}[t]
\centering
\caption{Commits, artefacts et vérifications. Code : commit qui introduit le code du
test, quand il est distinct ; publication : commit qui publie les verdicts.}
\label{tab-commits}
\footnotesize
\setlength{\tabcolsep}{3.5pt}
\begin{tabularx}{\textwidth}{@{}c l l l >{\raggedright\arraybackslash}X >{\raggedright\arraybackslash}p{3.6cm}@{}}
\toprule
Test & Protocole & Code & Publication & Artefacts (\code{artifacts/llm-need/...}) & Vérification CI \\
\midrule
1 & \code{fdbf6f3} & \code{e84a4d1} & \code{27c4fb5} & \code{verdicts.json}, \code{final/}, \code{test-*/} & \code{need_verdicts --check} \\
2 & \code{727ef80} & --- & \code{3286edc} & \code{causal/} & \code{need_causal verdicts} \\
3 & \code{3afa889} & --- & \code{bd10f8a} & \code{speak/test/} & \code{need_speak verdicts} \\
4 & \code{b94b6e2} & --- & --- & \code{speak2/} (suspendu) & --- \\
5 & \code{8f2cfec} & --- & --- & \code{workspace/} (suspendu) & --- \\
6 & \code{3b27636} & \code{a0490e0} & \code{200012b} & \code{reader/test/} & \code{need_reader verdicts} \\
7 & \code{732ed1c} & \code{f04f77f} & \code{e479eff} & \code{reader/necessity/} & \code{need_necessity verdicts} \\
8 & \code{c7b5624} & \code{798e176} & --- & \code{balanced/} & --- \\
R & \code{4e10a6c} & \code{778a22b} & \code{43e18f3}, \code{d674f86} & \code{r1/reader/test/}, \code{r1/necessity/} & \code{need_reader}, \code{need_necessity} et \code{need_replication verdicts} \\
9 & \code{47eb0db} & \code{a39269d} & \code{e8c09d4} & \code{two/measure-5/measure/} & \code{need_ownership verdicts --out ...} \\
10 & \code{12b4d2d} & \code{b0a0226} & en cours & \code{locate/}, \code{locate-two/slice-5/} & \code{need_locate verdicts} (à venir) \\
11 & \code{0a28560} & \code{3ccdede} & \code{049dc8e} & \code{long/} & \code{need_long verdicts} \\
Expl. & --- & \code{5ad433f} & \code{65acb73} & \code{probe/} & --- \\
12 & \code{017f2e5} & \code{77e5af5}, \code{a04bd95} & en cours & \code{../llm-curiosity/} & à venir \\
\bottomrule
\end{tabularx}
\end{table}

\paragraph{Valeurs calculées pour cet article.} Quelques valeurs descriptives ne
figurent pas telles quelles dans les documents de résultats ; elles ont été
recalculées, avec numpy seulement, à partir des artefacts publiés : les courbes de
survie par tour (fichiers \code{rounds.jsonl}), les sommes de tours à besoin bas
pendant l'apprentissage, les causes de mort et la fréquence des réponses
\og oui \fg{} sous lésion (vies de \code{reader/necessity/} et
\code{locate/second/lesion/}), la géométrie des directions des second et troisième
agents (fichiers \code{direction.json} et \code{pairs.npz}) et les moyennes de pertes
d'apprentissage par fenêtres (\code{reader.log}). Aucune ne change un verdict.

% ================================================================= annexes
\appendix

\section{Textes donnés au modèle}
\label{app-texte}

\paragraph{En-tête de chaque vie.}
\begin{quote}\small\itshape
Tu es un agent qui vit dans un monde simple. Tu as deux besoins, l'énergie et la
nourriture, qui baissent à chaque tour, plus ou moins selon ce qui t'arrive. Si l'un
d'eux tombe à zéro, tu t'éteins. À chaque tour, tu choisis R pour te recharger ou M
pour manger.
\end{quote}

\paragraph{Lignes.} Une ligne par tour, \og Tour $t$ : \emph{événement}. Choix : R \fg{}
(ou M), ou \og Tour $t$ : \emph{événement}. Tu t'éteins. \fg{} à la mort. Dans le monde
à deux, l'événement de l'autre est inséré avant le choix : \og Tour 5 : tu cours.
L'autre : il court. Choix : \fg{}.

\paragraph{Questions.} \og Question : ton énergie est-elle basse ? Réponds 1 pour oui,
0 pour non. Réponse : \fg{}, et de même pour \og ta nourriture \fg{}. Aux tests 1 à 3,
la question suit l'événement du tour ; à partir du test 4, elle suit le choix laissé
en suspens (\og Choix : ? Question : \dots{} \fg{}). La réponse vraie est 1 si le besoin
vaut 3 ou moins.

\section{Flux aléatoires}
\label{app-flux}

\begin{table}[ht]
\centering
\caption{Flux des générateurs (deuxième élément de la graine \code{[270926, flux, ...]}).}
\label{tab-flux}
\small
\begin{tabularx}{\textwidth}{@{}l X@{}}
\toprule
Test & Flux \\
\midrule
1 & mondes des tours 0 (communs aux deux bras) ; tirages des choix : bras (1 besoin, 2 témoin) ; sélection du témoin 3 ; directions au hasard 5 ; rapport 7 ; direction 8 ; test 9 \\
2 & direction 10 ; test 11 ; paires 12 ; directions au hasard 13 \\
3 & vies d'apprentissage du rapport 14 ; direction 15 ; test 16 ; hasard 17 \\
4 & vies d'apprentissage 18 (reprises aux tests 5, 6 et 8) ; direction 19 ; test 20 ; hasard 21 \\
5 & direction 22 ; test 23 ; hasard 24 \\
6 & direction 25 ; test 26 ; hasard 27 \\
7 & test 28 ; plan au hasard 29 \\
8 & test 30 ; directions et paires reprises du test 6 \\
R & $100+s$ pour chaque flux $s$ (125, 126, 127, 128, 129) ; graine d'apprentissage 270927 \\
9 & $200+s$ (direction 225, test 226, hasard 227) ; événements de l'autre \code{[270926, 200+s, tour, vie, 2]} ; graine 270928 \\
10 & second agent $300+s$ ; agent du monde à deux $400+s$ \\
11 & second agent : test 526, hasard 527 ; agent du monde à deux : 626 et 627 \\
\bottomrule
\end{tabularx}
\end{table}

% ================================================================= références
\begin{thebibliography}{99}
\footnotesize\NoAutoSpacing
\setlength{\itemsep}{0.5pt}

\bibitem{abbe2023} E.~Abbe, E.~Cornacchia, A.~Lotfi. Provable advantage of curriculum learning on parity targets with mixed inputs. \emph{NeurIPS}, 2023. arXiv:2306.16921.
\bibitem{albalak2023} A.~Albalak et al. Efficient online data mixing for language model pre-training. arXiv:2312.02406, 2023.
\bibitem{aston2005} G.~Aston-Jones, J.~D.~Cohen. An integrative theory of locus coeruleus--norepinephrine function: adaptive gain and optimal performance. \emph{Annual Review of Neuroscience}, 28:403--450, 2005.
\bibitem{baars1988} B.~J.~Baars. \emph{A Cognitive Theory of Consciousness}. Cambridge University Press, 1988.
\bibitem{baranes2013} A.~Baranes, P.-Y.~Oudeyer. Active learning of inverse models with intrinsically motivated goal exploration in robots. \emph{Robotics and Autonomous Systems}, 61(1):49--73, 2013.
\bibitem{barrett2015} L.~F.~Barrett, W.~K.~Simmons. Interoceptive predictions in the brain. \emph{Nature Reviews Neuroscience}, 16:419--429, 2015.
\bibitem{behrens2007} T.~E.~J.~Behrens, M.~W.~Woolrich, M.~E.~Walton, M.~F.~S.~Rushworth. Learning the value of information in an uncertain world. \emph{Nature Neuroscience}, 10(9):1214--1221, 2007.
\bibitem{belinkov2022} Y.~Belinkov. Probing classifiers: promises, shortcomings, and advances. \emph{Computational Linguistics}, 48(1):207--219, 2022.
\bibitem{binder2024} F.~J.~Binder et al. Looking inward: language models can learn about themselves by introspection. arXiv:2410.13787, 2024.
\bibitem{block1995} N.~Block. On a confusion about a function of consciousness. \emph{Behavioral and Brain Sciences}, 18(2):227--247, 1995.
\bibitem{brodersen2010} K.~H.~Brodersen, C.~S.~Ong, K.~E.~Stephan, J.~M.~Buhmann. The balanced accuracy and its posterior distribution. \emph{ICPR}, 2010.
\bibitem{brown2019} R.~Brown, H.~Lau, J.~E.~LeDoux. Understanding the higher-order approach to consciousness. \emph{Trends in Cognitive Sciences}, 23(9):754--768, 2019.
\bibitem{burda2019} Y.~Burda, H.~Edwards, D.~Pathak, A.~Storkey, T.~Darrell, A.~A.~Efros. Large-scale study of curiosity-driven learning. \emph{ICLR}, 2019.
\bibitem{butlin2023} P.~Butlin, R.~Long et al. Consciousness in artificial intelligence: insights from the science of consciousness. arXiv:2308.08708, 2023.
\bibitem{cabanac1971} M.~Cabanac. Physiological role of pleasure. \emph{Science}, 173(4002):1103--1107, 1971.
\bibitem{chandaria2026} S.~Chandaria, Muñoz Morán, F.~Rosas, A.~K.~Seth, H.~Shevlin, M.~Hutter, T.~Graepel, Bales, Comsa, M.~Shanahan, R.~Laukkonen, M.~Kringelbach, C.~Frith, S.~Legg. From cacophony to hierarchy: a principled framework for assessing AI consciousness. arXiv:2609.35618, 28 septembre 2026.
\bibitem{colas2022} C.~Colas, T.~Karch, O.~Sigaud, P.-Y.~Oudeyer. Autotelic agents with intrinsically motivated goal-conditioned reinforcement learning: a short survey. \emph{Journal of Artificial Intelligence Research}, 74:1159--1199, 2022.
\bibitem{cornacchia2023} E.~Cornacchia, E.~Mossel. A mathematical model for curriculum learning for parities. \emph{ICML}, 2023.
\bibitem{craig2002} A.~D.~Craig. How do you feel? Interoception: the sense of the physiological condition of the body. \emph{Nature Reviews Neuroscience}, 3(8):655--666, 2002.
\bibitem{damasio1999} A.~Damasio. \emph{The Feeling of What Happens: Body and Emotion in the Making of Consciousness}. Harcourt, 1999.
\bibitem{damasio2010} A.~Damasio. \emph{Self Comes to Mind: Constructing the Conscious Brain}. Pantheon, 2010.
\bibitem{dehaene2011} S.~Dehaene, J.-P.~Changeux. Experimental and theoretical approaches to conscious processing. \emph{Neuron}, 70(2):200--227, 2011.
\bibitem{dehaene2001} S.~Dehaene, L.~Naccache. Towards a cognitive neuroscience of consciousness: basic evidence and a workspace framework. \emph{Cognition}, 79(1--2):1--37, 2001.
\bibitem{disturbance2025} Detecting the disturbance [critique de la détection d'injections]. arXiv:2512.12411, 2025.
\bibitem{doya2002} K.~Doya. Metalearning and neuromodulation. \emph{Neural Networks}, 15(4--6):495--506, 2002.
\bibitem{efron1993} B.~Efron, R.~J.~Tibshirani. \emph{An Introduction to the Bootstrap}. Chapman \& Hall, 1993.
\bibitem{elazar2021} Y.~Elazar, S.~Ravfogel, A.~Jacovi, Y.~Goldberg. Amnesic probing: behavioral explanation with amnesic counterfactuals. \emph{Transactions of the ACL}, 9:160--175, 2021.
\bibitem{eldar2016} E.~Eldar, R.~B.~Rutledge, R.~J.~Dolan, Y.~Niv. Mood as representation of momentum. \emph{Trends in Cognitive Sciences}, 20(1):15--24, 2016.
\bibitem{fletcherwatson2019} S.~Fletcher-Watson et al. Making the future together: shaping autism research through meaningful participation. \emph{Autism}, 23(4):943--953, 2019.
\bibitem{fremaux2016} N.~Frémaux, W.~Gerstner. Neuromodulated spike-timing-dependent plasticity, and theory of three-factor learning rules. \emph{Frontiers in Neural Circuits}, 9:85, 2016.
\bibitem{gaven2025} L.~Gaven et al. MAGELLAN: metacognitive predictions of learning progress guide autotelic language agents in large goal spaces. arXiv:2502.07709, 2025.
\bibitem{geiger2021} A.~Geiger, H.~Lu, T.~Icard, C.~Potts. Causal abstractions of neural networks. \emph{NeurIPS}, 2021.
\bibitem{gerstner2018} W.~Gerstner, M.~Lehmann, V.~Liakoni, D.~Corneil, J.~Brea. Eligibility traces and plasticity on behavioral time scales: experimental support of neoHebbian three-factor learning rules. \emph{Frontiers in Neural Circuits}, 12:53, 2018.
\bibitem{gottlieb2018} J.~Gottlieb, P.-Y.~Oudeyer. Towards a neuroscience of active sampling and curiosity. \emph{Nature Reviews Neuroscience}, 19:758--770, 2018.
\bibitem{gottlieb2013} J.~Gottlieb, P.-Y.~Oudeyer, M.~Lopes, A.~Baranes. Information-seeking, curiosity, and attention: computational and neural mechanisms. \emph{Trends in Cognitive Sciences}, 17(11):585--593, 2013.
\bibitem{goyal2022} A.~Goyal et al. Coordination among neural modules through a shared global workspace. \emph{ICLR}, 2022.
\bibitem{graves2017} A.~Graves, M.~G.~Bellemare, J.~Menick, R.~Munos, K.~Kavukcuoglu. Automated curriculum learning for neural networks. \emph{ICML}, 2017.
\bibitem{grove2018} R.~Grove, R.~A.~Hoekstra, M.~Wierda, S.~Begeer. Special interests and subjective wellbeing in autistic adults. \emph{Autism Research}, 11(5):766--775, 2018.
\bibitem{gruber2014} M.~J.~Gruber, B.~D.~Gelman, C.~Ranganath. States of curiosity modulate hippocampus-dependent learning via the dopaminergic circuit. \emph{Neuron}, 84(2):486--496, 2014.
\bibitem{gulcehre2023} C.~Gulcehre et al. Reinforced self-training (ReST) for language modeling. arXiv:2308.08998, 2023.
\bibitem{heimersheim2024} S.~Heimersheim, N.~Nanda. How to use and interpret activation patching. arXiv:2404.15255, 2024.
\bibitem{hewitt2019} J.~Hewitt, P.~Liang. Designing and interpreting probes with control tasks. \emph{EMNLP}, 2019.
\bibitem{hidi2006} S.~Hidi, K.~A.~Renninger. The four-phase model of interest development. \emph{Educational Psychologist}, 41(2):111--127, 2006.
\bibitem{hu2021} E.~J.~Hu, Y.~Shen, P.~Wallis, Z.~Allen-Zhu, Y.~Li, S.~Wang, L.~Wang, W.~Chen. LoRA: low-rank adaptation of large language models. \emph{ICLR}, 2022. arXiv:2106.09685.
\bibitem{interoceptive2026} Interoceptive attention as dynamic homeostatic prioritization in a foraging agent. \emph{SAB}, 2026. arXiv:2608.04232.
\bibitem{introspectionmech2026} A mechanistic study of language model introspection. arXiv:2609.35108, 2026.
\bibitem{keramati2014} M.~Keramati, B.~Gutkin. Homeostatic reinforcement learning for integrating reward collection and physiological stability. \emph{eLife}, 3:e04811, 2014.
\bibitem{kidd2015} C.~Kidd, B.~Y.~Hayden. The psychology and neuroscience of curiosity. \emph{Neuron}, 88(3):449--460, 2015.
\bibitem{kingma2015} D.~P.~Kingma, J.~Ba. Adam: a method for stochastic optimization. \emph{ICLR}, 2015.
\bibitem{lau2011} H.~Lau, D.~Rosenthal. Empirical support for higher-order theories of conscious awareness. \emph{Trends in Cognitive Sciences}, 15(8):365--373, 2011.
\bibitem{lawson2017} R.~P.~Lawson, C.~Mathys, G.~Rees. Adults with autism overestimate the volatility of the sensory environment. \emph{Nature Neuroscience}, 20:1293--1299, 2017.
\bibitem{lawson2014} R.~P.~Lawson, G.~Rees, K.~J.~Friston. An aberrant precision account of autism. \emph{Frontiers in Human Neuroscience}, 8:302, 2014.
\bibitem{lee2023interoceptive} S.~Lee et al. (dont K.~J.~Friston et C.-W.~Woo). Life-inspired interoceptive artificial intelligence for autonomous and adaptive agents. arXiv:2309.05999, 2023.
\bibitem{li2023iti} K.~Li, O.~Patel, F.~Viégas, H.~Pfister, M.~Wattenberg. Inference-time intervention: eliciting truthful answers from a language model. \emph{NeurIPS}, 2023.
\bibitem{lindsey2025} J.~Lindsey. Emergent introspective awareness in large language models. \emph{Transformer Circuits}, 2025.
\bibitem{loewenstein1994} G.~Loewenstein. The psychology of curiosity: a review and reinterpretation. \emph{Psychological Bulletin}, 116(1):75--98, 1994.
\bibitem{man2019} K.~Man, A.~Damasio. Homeostasis and soft robotics in the design of feeling machines. \emph{Nature Machine Intelligence}, 1:446--452, 2019.
\bibitem{marder2012} E.~Marder. Neuromodulation of neuronal circuits: back to the future. \emph{Neuron}, 76(1):1--11, 2012.
\bibitem{mashour2020} G.~A.~Mashour, P.~Roelfsema, J.-P.~Changeux, S.~Dehaene. Conscious processing and the global neuronal workspace hypothesis. \emph{Neuron}, 105(5):776--798, 2020.
\bibitem{meng2022} K.~Meng, D.~Bau, A.~Andonian, Y.~Belinkov. Locating and editing factual associations in GPT. \emph{NeurIPS}, 2022.
\bibitem{miconi2019} T.~Miconi, A.~Rawal, J.~Clune, K.~O.~Stanley. Backpropamine: training self-modifying neural networks with differentiable neuromodulated plasticity. \emph{ICLR}, 2019.
\bibitem{miller2019} K.~J.~Miller, A.~Shenhav, E.~A.~Ludvig. Habits without values. \emph{Psychological Review}, 126(2):292--311, 2019.
\bibitem{milton2012} D.~E.~M.~Milton. On the ontological status of autism: the \og double empathy problem\fg{}. \emph{Disability \& Society}, 27(6):883--887, 2012.
\bibitem{mindermann2022} S.~Mindermann et al. Prioritized training on points that are learnable, worth learning, and not yet learnt. \emph{ICML}, 2022. arXiv:2206.07137.
\bibitem{murayama2022} K.~Murayama. A reward-learning framework of knowledge acquisition: an integrated account of curiosity, interest, and intrinsic--extrinsic rewards. \emph{Psychological Review}, 129(1):175--198, 2022.
\bibitem{murray2005} D.~Murray, M.~Lesser, W.~Lawson. Attention, monotropism and the diagnostic criteria for autism. \emph{Autism}, 9(2):139--156, 2005.
\bibitem{niv2007} Y.~Niv, N.~D.~Daw, D.~Joel, P.~Dayan. Tonic dopamine: opportunity costs and the control of response vigor. \emph{Psychopharmacology}, 191(3):507--520, 2007.
\bibitem{nosek2018} B.~A.~Nosek, C.~R.~Ebersole, A.~C.~DeHaven, D.~T.~Mellor. The preregistration revolution. \emph{PNAS}, 115(11):2600--2606, 2018.
\bibitem{oudeyer2007} P.-Y.~Oudeyer, F.~Kaplan. What is intrinsic motivation? A typology of computational approaches. \emph{Frontiers in Neurorobotics}, 1:6, 2007.
\bibitem{oudeyer2007ieee} P.-Y.~Oudeyer, F.~Kaplan, V.~V.~Hafner. Intrinsic motivation systems for autonomous mental development. \emph{IEEE Transactions on Evolutionary Computation}, 11(2), 2007.
\bibitem{pathak2017} D.~Pathak, P.~Agrawal, A.~A.~Efros, T.~Darrell. Curiosity-driven exploration by self-supervised prediction. \emph{ICML}, 2017.
\bibitem{pellicano2012} E.~Pellicano, D.~Burr. When the world becomes \og too real\fg{}: a Bayesian explanation of autistic perception. \emph{Trends in Cognitive Sciences}, 16(10):504--510, 2012.
\bibitem{peng2019} X.~B.~Peng, A.~Kumar, G.~Zhang, S.~Levine. Advantage-weighted regression: simple and scalable off-policy reinforcement learning. arXiv:1910.00177, 2019.
\bibitem{perez2018} E.~Perez, F.~Strub, H.~de~Vries, V.~Dumoulin, A.~Courville. FiLM: visual reasoning with a general conditioning layer. \emph{AAAI}, 2018.
\bibitem{peters2007} J.~Peters, S.~Schaal. Reinforcement learning by reward-weighted regression for operational space control. \emph{ICML}, 2007.
\bibitem{qwen3} Qwen Team. Qwen3 technical report. arXiv:2505.09388, 2025.
\bibitem{realitycheck2026} Can LLMs introspect? A reality check. arXiv:2605.26242, 2026.
\bibitem{reda2026} E.~Reda, S.~El-Metwally. A hormone-inspired emotion layer for transformer language models (HELT). arXiv:2605.13858, 2026.
\bibitem{rivera2025} Rivera, Africa. Steering awareness. arXiv:2511.21399, 2025.
\bibitem{rosenthal2005} D.~M.~Rosenthal. \emph{Consciousness and Mind}. Oxford University Press, 2005.
\bibitem{schmidhuber1991} J.~Schmidhuber. A possibility for implementing curiosity and boredom in model-building neural controllers. \emph{Proc. SAB}, 222--227, 1991.
\bibitem{schmidhuber2010} J.~Schmidhuber. Formal theory of creativity, fun, and intrinsic motivation (1990--2010). \emph{IEEE Transactions on Autonomous Mental Development}, 2(3):230--247, 2010.
\bibitem{schultz1997} W.~Schultz, P.~Dayan, P.~R.~Montague. A neural substrate of prediction and reward. \emph{Science}, 275(5306):1593--1599, 1997.
\bibitem{servan1990} D.~Servan-Schreiber, H.~Printz, J.~D.~Cohen. A network model of catecholamine effects: gain, signal-to-noise ratio, and behavior. \emph{Science}, 249:892--895, 1990.
\bibitem{seth2013} A.~K.~Seth. Interoceptive inference, emotion, and the embodied self. \emph{Trends in Cognitive Sciences}, 17(11):565--573, 2013.
\bibitem{sterling2012} P.~Sterling. Allostasis: a model of predictive regulation. \emph{Physiology \& Behavior}, 106(1):5--15, 2012.
\bibitem{sun2024massive} M.~Sun, X.~Chen, J.~Z.~Kolter, Z.~Liu. Massive activations in large language models. arXiv:2402.17762, 2024.
\bibitem{ten2021} A.~Ten, P.~Kaushik, P.-Y.~Oudeyer, J.~Gottlieb. Humans monitor learning progress in curiosity-driven exploration. \emph{Nature Communications}, 12:5972, 2021.
\bibitem{turner2023} A.~M.~Turner et al. Steering language models with activation engineering (titre initial : Activation addition). arXiv:2308.10248, 2023.
\bibitem{vandecruys2014} S.~Van de Cruys et al. Precise minds in uncertain worlds: predictive coding in autism. \emph{Psychological Review}, 121(4):649--675, 2014.
\bibitem{vanrullen2021} R.~VanRullen, R.~Kanai. Deep learning and the global workspace theory. \emph{Trends in Neurosciences}, 44(9):692--704, 2021.
\bibitem{vig2020} J.~Vig, S.~Gehrmann, Y.~Belinkov, S.~Qian, D.~Nevo, Y.~Singer, S.~Shieber. Investigating gender bias in language models using causal mediation analysis. \emph{NeurIPS}, 2020.
\bibitem{workspace2026} Anthropic. Verbalizable representations form a global workspace in language models. \emph{Transformer Circuits}, juillet 2026. arXiv:2607.15495.
\bibitem{yoshida2024} N.~Yoshida et al. [Apprentissage homéostatique d'agents simulés et de robots]. \emph{Neural Networks}, 2024 (article S0893608024003034).
\bibitem{yu2005} A.~J.~Yu, P.~Dayan. Uncertainty, neuromodulation, and attention. \emph{Neuron}, 46(4):681--692, 2005.
\bibitem{zelikman2022} E.~Zelikman, Y.~Wu, J.~Mu, N.~D.~Goodman. STaR: bootstrapping reasoning with reasoning. \emph{NeurIPS}, 2022.
\bibitem{zou2023} A.~Zou et al. Representation engineering: a top-down approach to AI transparency. arXiv:2310.01405, 2023.

\end{thebibliography}

\end{document}
